\documentclass[11pt]{article}

\usepackage[a4paper, margin=1in]{geometry}

\usepackage{amsmath}
\usepackage{amsfonts}
\usepackage{amssymb}
\usepackage{mathtools}
\usepackage{microtype}
\usepackage{lmodern}
\usepackage{graphicx}
\usepackage{float}
\usepackage{subcaption}
\usepackage{booktabs}
\usepackage{longtable}
\usepackage{xcolor}
\usepackage{hyperref}

\hypersetup{
colorlinks=false,
linkcolor=blue,
urlcolor=blue,
citecolor=blue
}

\usepackage{fancyhdr}
\pagestyle{fancy}
\fancyhf{}
\lhead{Neurological Foundations of the Mind}
\rhead{\thepage}
\renewcommand{\headrulewidth}{0.4pt}

\usepackage{titlesec}

\titleformat{\section}
{\large\bfseries}
{\thesection}{1em}{}

\titleformat{\subsection}
{\normalsize\bfseries}
{\thesubsection}{1em}{}

\usepackage{listings}
\lstset{
basicstyle=\ttfamily\small,
breaklines=true,
frame=single
}

\newcommand{\R}{\mathbb{R}}
\newcommand{\vct}[1]{\mathbf{#1}}
\newcommand{\mat}[1]{\mathbf{#1}}

\title{\textbf{From Potentials and Chemicals to the Mind} \\
\large \textit{The Neurological and Psychological Foundations of the Mind}}
\author{Jack Wetherell}
\date{}

\begin{document}

\maketitle
\tableofcontents
\newpage

\phantomsection
\section*{Introduction}
\addcontentsline{toc}{section}{Introduction}

The title of these notes is meant literally. They begin with ions crossing a membrane and they intend to arrive, by stages and without any step that is not physical, at seeing, remembering, wanting, fearing, and knowing oneself to be a self. Whether that journey can be completed is a real question. What is not in doubt is that a great deal of it has already been made, and the point of writing these notes is to see how far the completed part actually goes.

It is worth being precise, once, about which question is being asked, because two quite different ones travel under the same name. Chalmers drew the distinction that has stuck. The \textit{easy} problems of consciousness are the ones that ask how a physical system manages to do the things minds do: how it discriminates one stimulus from another, integrates separate streams into a single scene, stores what happened and retrieves it, attends to one thing rather than another, assigns value, controls behaviour in the service of a goal, reports on its own states, and builds a model of itself. They are called easy not because they are simple, since several of them have taken a century and are not finished, but because it is clear what an answer would look like: a mechanism, specified well enough that one could see how it does the job. The \textit{hard} problem asks something else. It asks why any of that machinery, however completely described, should be accompanied by experience at all, why there is something it is like to be the system rather than nothing.

These notes pursue the easy problems. They do so because that is where the neuroscience is, and because the progress there is genuinely remarkable and deserves to be laid out in order. The hard problem is named here and then set aside; it is not raised again, and nothing in what follows should be read as an attempt on it. A reader who thinks the hard problem dissolves once the easy ones are solved, and a reader who thinks it survives untouched, can both use this document, because it is an account of mechanism either way.

The method is to work up through levels. An ion has a concentration gradient; a channel has a conductance; a synapse has a strength; a neuron has a firing rate; a circuit has a computation; a nucleus has a function; a system has a capacity. Each level has its own vocabulary, and the vocabulary of one level does not appear at the level below: there is no fear in a sodium channel and no attention in a vesicle. What there is instead is a chain of explanation, in which each level is accounted for by the one beneath it. This is why a document about the mind opens with the Nernst equation and closes with a bear on a woodland path, and why the equation turns out to be doing work in the bear.

One claim recurs so often in what follows that it is worth naming in advance, so that it can be recognised each time it appears. It is that nothing about the mind is given. Seeing feels like opening the eyes and receiving the world; it is a construction performed in stages by cortex, and it fails in strikingly specific ways when particular stages are damaged. The sense of occupying a body in a place feels primitive and unanalysable; it is built, and it can be taken apart. Feeling an emotion feels like the most direct knowledge one has; it is assembled, late, out of a reading of the body's own state. In every case the immediacy is real as an experience and misleading as evidence, because the construction is invisible from the inside. That the seams do not show is a fact about the machinery, not about what the machinery is doing.

The order of the sections follows the ascent. The first four establish the substrate: the layout of the nervous system, the structure of the neuron, the electrical and chemical events by which neurons signal, and the glial cells without which none of it works. The next five walk the cortex, and are where the constructive claim is made repeatedly and in detail, as sensation becomes perception, perception becomes recognition, and recognition becomes a world with a body in it. The last four go beneath the cortex, to the structures that select actions, generate drives, gate what reaches awareness, and assign to events the significance that makes any of them worth responding to.

Some things are deliberately absent. There is no developmental account, little on the cortical microcircuit, and no systematic treatment of psychiatric illness, though all three are touched where they illuminate something. Clinical material appears throughout, but as evidence rather than as medicine: a deficit that follows a specific lesion is the closest thing this subject has to a controlled experiment on the human mind, and the argument leans on that evidence heavily.

\section{Nervous System Structure}

Before we can talk about thoughts, feelings, memories, or the mind itself, we need a map. The nervous system is the organ that does all of this, it senses, decides, remembers, imagines, and like any complicated machine it helps enormously to know, at least roughly, where its parts are and what each part is for.

At the highest level the system splits into two halves. The \textbf{central nervous system (CNS)} is the brain and the spinal cord: the headquarters, where information is integrated and decisions are made. The \textbf{peripheral nervous system (PNS)} is everything else, the sprawling network of nerves that links the CNS to muscles, glands, and sensory organs. Every sensation and every voluntary movement crosses the boundary between these two systems.

The rest of this section walks through the major structures of the CNS and PNS from the top down, establishing where each of the regions discussed later actually sits.

\subsection{The Cerebrum}

The cerebrum is the familiar part of the brain, the wrinkled, walnut-shaped dome that sits on top of everything else. It is also, by some distance, the largest part of the brain and the part most closely associated with what we usually mean by ``the mind.'' Thinking, reasoning, memory, language, voluntary movement, the experience of being a self, all of these depend on cerebral machinery.

Structurally it is symmetrical. Two hemispheres, left and right, are joined by a thick band of white matter called the \textbf{corpus callosum}, which lets them share information. Each hemisphere is then divided into four lobes, and each lobe has a characteristic job:

\begin{itemize}
    \item \textbf{Frontal lobe}, voluntary movement, planning, decision-making, and much of what we call personality.
    \item \textbf{Parietal lobe}, processing of bodily sensations such as touch, temperature, and pain, and spatial awareness.
    \item \textbf{Temporal lobe}, hearing, language comprehension, and the formation of long-term memories.
    \item \textbf{Occipital lobe}, vision.
\end{itemize}

The outer layer of the cerebrum, the \textbf{cerebral cortex}, is only a few millimetres thick but is crumpled into deep folds. The folding is not cosmetic: by crumpling the sheet, evolution has packed far more cortical area into the skull than would otherwise fit, and cortical area turns out to matter enormously for cognitive capacity.

\subsection{The Cerebellum}

Tucked beneath the back of the cerebrum, the cerebellum (``little brain'') looks almost like a miniature version of its larger sibling. Despite its size, it contains roughly half of all the neurons in the brain, and it is indispensable for smooth, coordinated movement.

Crucially, the cerebellum does not initiate actions. Instead it receives a copy of the motor commands being issued by the cortex together with sensory feedback about what the body is actually doing, and it compares the two. When the intention and the reality diverge, the cerebellum quietly corrects the discrepancy in real time. That is why skilled movements, reaching for a mug, playing a scale on the piano, walking in a straight line, feel effortless.

When the cerebellum is damaged, this silent correction falls apart. The classic picture includes \textit{ataxia} (loss of coordination), tremor, and difficulty maintaining balance. Patients often describe it as feeling perpetually drunk, which is itself a hint: alcohol disrupts cerebellar function, and that is exactly why it impairs coordination.

\subsection{The Diencephalon}

Between the cerebrum above and the brainstem below sits the \textbf{diencephalon}, a compact group of structures wrapped around the third ventricle. It is small relative to the hemispheres that enclose it, and it is easy to overlook on a surface view because almost none of it is visible from outside, but very little reaches the cortex without passing through it or being regulated by it.

It comprises four parts:

\begin{itemize}
    \item The \textbf{thalamus} is by far the largest, a pair of egg-shaped masses of grey matter sitting either side of the midline. It is the great relay station of the brain: with the single exception of smell, every sensory stream, vision, hearing, touch, taste, and the sense of the body's own position, synapses in a specific thalamic nucleus before continuing to the cortex. It is not a passive junction box. The thalamus gates what reaches the cortex according to attention and arousal, and it carries return traffic from the cortex back down as well, taking part in the motor and cognitive loops discussed later.
    \item The \textbf{hypothalamus} lies beneath the thalamus and is, for its size, the most consequential structure in the body. It is the master regulator of \textbf{homeostasis}, holding body temperature, fluid balance, hunger, thirst, and the sleep-wake cycle within viable limits. It commands the autonomic nervous system, and through its connection to the \textbf{pituitary gland} it also commands the endocrine system, which makes it the point at which the nervous and hormonal systems become one control problem rather than two.
    \item The \textbf{epithalamus} is a small posterior region whose principal component is the \textbf{pineal gland}. The pineal secretes \textbf{melatonin} on a daily rhythm and is central to the timing of the sleep-wake cycle.
    \item The \textbf{subthalamus} lies below the thalamus and contains the \textbf{subthalamic nucleus}, a small structure that plays a disproportionately large role in the control of movement, as the section on the basal ganglia sets out.
\end{itemize}

The diencephalon is therefore best understood not as one organ but as the junction at which the traffic between cortex, brainstem, and body is routed, filtered, and regulated. Damage here rarely produces a single clean deficit, precisely because so many separate pathways are packed into so small a volume.

\subsection{The Midbrain and Brainstem}

If the cerebrum is where thought happens and the cerebellum is where movement is refined, the brainstem is where life is sustained. It connects the brain to the spinal cord and is the conduit for almost every signal passing between them. It is also the home of the circuits that drive breathing, set the rhythm of the heart, and maintain consciousness itself.

It is traditionally divided into three levels, each with distinct responsibilities:

\begin{itemize}
    \item \textbf{Midbrain}, visual and auditory reflexes, aspects of motor control, and arousal.
    \item \textbf{Pons}, a relay between the cerebrum and cerebellum, and a participant in sleep and respiration.
    \item \textbf{Medulla oblongata}, the control centre for vital autonomic functions: heart rate, blood pressure, breathing, coughing, swallowing.
\end{itemize}

Because the brainstem handles these non-negotiable tasks, even small lesions there can be catastrophic. A few millimetres in the wrong place can interrupt consciousness or stop breathing. This is also why brainstem death is used as a clinical definition of death.

\subsection{Cranial Nerves}

Most nerves reach the body via the spinal cord, but the twelve \textbf{cranial nerves} bypass that route and leave the CNS directly, mostly from the brainstem. They handle the close work of the head and neck, vision, eye movement, facial sensation and expression, taste, hearing, balance, swallowing, and, in the case of the vagus nerve, reach deep into the thorax and abdomen to regulate several visceral functions.

Each is conventionally numbered with a Roman numeral, running from front to back in the order the nerves emerge. Each is also classified as \textit{sensory}, \textit{motor}, or \textit{both}, according to the direction of the information it carries:

\begin{itemize}
    \item \textbf{I, Olfactory} (sensory). Carries smell. Uniquely among the cranial nerves it does not arise from the brainstem, projecting instead directly to the cerebrum, a peculiarity taken up in the section on the temporal lobe.
    \item \textbf{II, Optic} (sensory). Carries vision from the retina, partially crossing at the optic chiasm so that each hemisphere receives the opposite half of the visual field.
    \item \textbf{III, Oculomotor} (motor). Drives four of the six muscles that move the eye, raises the upper eyelid, and carries the parasympathetic fibres that constrict the pupil and focus the lens.
    \item \textbf{IV, Trochlear} (motor). Supplies a single muscle, the superior oblique, which turns the eye downward and inward. It is the thinnest cranial nerve and the only one to emerge from the back of the brainstem.
    \item \textbf{V, Trigeminal} (both). The great sensory nerve of the face, carrying touch, pain, and temperature from the skin, the cornea, the sinuses, and the front of the tongue through its three divisions, the ophthalmic, maxillary, and mandibular. Its motor fibres drive the muscles of chewing.
    \item \textbf{VI, Abducens} (motor). Supplies the lateral rectus, which turns the eye outward.
    \item \textbf{VII, Facial} (both). Drives the muscles of facial expression, carries taste from the front two-thirds of the tongue, and supplies the glands that produce tears and saliva.
    \item \textbf{VIII, Vestibulocochlear} (sensory). Carries hearing from the cochlea and balance from the vestibular apparatus, in two divisions that travel together but report on quite different things.
    \item \textbf{IX, Glossopharyngeal} (both). Carries taste and sensation from the back of the tongue and the pharynx, including the blood pressure and oxygen sensors of the carotid body, and contributes to swallowing.
    \item \textbf{X, Vagus} (both). By far the most widely distributed, wandering (as its name says) out of the head entirely to supply the heart, lungs, and most of the digestive tract. It is the principal parasympathetic nerve of the body, and it also drives the muscles of the larynx and pharynx that make speech and swallowing possible.
    \item \textbf{XI, Accessory} (motor). Supplies the sternocleidomastoid and trapezius, the muscles that turn the head and shrug the shoulders.
    \item \textbf{XII, Hypoglossal} (motor). Supplies the muscles of the tongue, and so underlies both speech and the manipulation of food in the mouth.
\end{itemize}

Their arrangement is diagnostically valuable. Because the nerves emerge in an orderly sequence along the brainstem, and because each can be tested individually at the bedside, the pattern of which nerves fail and which are spared localises a brainstem lesion with considerable precision. A deficit combining several adjacent nerves points to a lesion at the level from which they emerge together.

\subsection{Ventricles and Cerebrospinal Fluid}

Inside the brain is a branching system of cavities called the \textbf{ventricles}: two lateral ventricles in the cerebral hemispheres, a third between them, and a fourth near the brainstem. They are filled with \textbf{cerebrospinal fluid (CSF)}, a clear fluid produced mostly by the \textit{choroid plexus} that lines them.

It is tempting to think of CSF as just a liquid cushion, and that is certainly part of its job, the brain effectively floats in it, which dramatically reduces its effective weight and protects it against mechanical shock. But CSF also helps maintain a stable chemical environment around neurons and carries metabolic waste away from nervous tissue. It is manufactured in the ventricles, circulates out over the surface of the brain and spinal cord, and is eventually reabsorbed into the venous blood.

When CSF flow is obstructed, the ventricles swell under pressure, a condition called \textbf{hydrocephalus}, which we will meet again when we talk about ependymal cells.

\subsection{The Spinal Cord}

The spinal cord is a bundle of nervous tissue about the thickness of a finger that runs down inside the vertebral column. It does two things at once. First, it is a \textbf{highway}: sensory information travels upward to the brain, and motor commands travel downward to the body. Second, and less obviously, it is a \textbf{simple decision-making device} in its own right. Reflexes such as withdrawing a hand from a hot surface are mediated by the spinal cord alone, without waiting for the brain to get involved. By the time the pain is consciously registered, the hand has already moved.

The cord is segmented, and each segment gives rise to a pair of \textbf{spinal nerves} that serve a particular territory of the body. This segmental organisation is why spinal injuries at different levels produce such different patterns of loss: a cervical lesion affects the arms and everything below, while a lumbar lesion may spare the arms completely.

\subsection{Peripheral Nerves and the Autonomic Nervous System}

The peripheral nervous system extends from the CNS out into the rest of the body. Functionally, peripheral nerves are either \textbf{sensory (afferent)}, carrying information from the body back to the CNS, or \textbf{motor (efferent)}, carrying commands from the CNS out to muscles and glands.

Motor control itself splits into two worlds. The \textbf{somatic nervous system} governs deliberate action: moving an arm, speaking, walking. The \textbf{autonomic nervous system} governs what proceeds without deliberation, heart rate, digestion, pupil size, sweating, blood pressure, and it does so through two opposing branches:

\begin{itemize}
    \item The \textbf{sympathetic division} prepares the body for action: heart rate rises, pupils dilate, blood is shunted toward muscle. This is the ``fight or flight'' response.
    \item The \textbf{parasympathetic division} does the opposite: it slows the heart, promotes digestion, and generally encourages the body to conserve and repair. This is ``rest and digest.''
\end{itemize}

This division matters far beyond the textbook. The balance between sympathetic and parasympathetic activity is the physiological signature of emotional arousal, stress, and relaxation, and it is where neurology and psychology start to meet in earnest: a frightened mind produces a sympathetically driven body, and a calm body feeds back into a calmer mind.

\subsection{Summary}

The nervous system is hierarchical and deeply integrated. The cerebrum handles the most elaborate cognition; the cerebellum refines motor output; the diencephalon relays and gates almost everything passing between them and the body below, and regulates the internal state of the organism; the brainstem keeps the lights on and gives rise to the cranial nerves that serve the head and neck; the spinal cord is both cable and reflex circuit; and peripheral nerves extend the system out into every organ and muscle in the body.

It is worth being clear about what this map does and does not provide. It is an inventory of parts and their locations, and nothing in it yet explains anything mental. Knowing that the occipital lobe is concerned with vision is not knowing how seeing is done, and knowing that the hypothalamus regulates hunger is not knowing why anything is wanted. What the map establishes is only where to look, which is a real thing to have and is the necessary first step. Everything after this section is the attempt to replace a location with a mechanism, and it begins by zooming in on the individual building blocks, the neurons themselves.

\section{Neuron Structure}

Neurons are the cells that make the nervous system a nervous system. Every tissue in the body signals in one way or another, but neurons have specialised for one job to an extreme degree: they move information, quickly and reliably, over long distances. Everything else about their biology, their shape, their metabolism, their internal traffic system, is bent toward that single purpose.

The most striking consequence of this specialisation is that neurons are \textbf{polarised}. They have a clearly defined input side and output side, and the whole architecture of the cell enforces one-way traffic. Understanding that shape is the first step toward understanding how neurons fire and communicate.

A neuron has three main regions, the \textbf{soma}, the \textbf{dendrites}, and the \textbf{axon}, plus a suite of internal machinery adapted to its unusual demands.

\subsection{Soma, Dendrites, and the Axon}

\subsubsection{Soma (Cell Body)}

The soma is the neuron's headquarters. It contains the nucleus and most of the cell's organelles, and it is where proteins are manufactured, energy is managed, and the general business of being a living cell is transacted. In a real sense the soma is the only part of the neuron that looks like a generic animal cell, everything else is specialised out.

Its most important computational job is \textbf{integration}. The soma, together with the axon hillock that we will meet in a moment, is where the neuron sums the inputs arriving on its dendrites and decides whether to fire. It is the place where ``many inputs'' becomes ``one decision.''

\subsubsection{Dendrites}

Dendrites are the branches that fan out from the soma like the root system of a tree, and they are the neuron's input antennae. A single neuron may have thousands of them, covered in small protrusions called \textit{dendritic spines}, and each spine can host a synapse from another neuron. The sheer geometry of this arrangement means that one neuron can listen to thousands of others at once.

What arrives at a dendrite is not, at first, an electrical spike. It is a \textit{chemical} signal, neurotransmitter released from the presynaptic cell, that is converted into a small, local change in membrane voltage called a \textbf{graded potential}. These graded potentials spread passively toward the soma, and they can either push the cell toward firing or pull it back from the brink. The dendritic tree is, in effect, the neuron's ear and its first stage of processing all at once.

\subsubsection{Axon}

The axon is the output. It is a single long projection that leaves the soma and carries electrical signals away to whatever the neuron talks to next, another neuron, a muscle cell, a gland. Axons can be astonishingly long: the motor neurons running from the spinal cord to the foot are over a metre in length in an adult, all in a single cell.

Three regions of the axon are worth naming explicitly:

\begin{itemize}
    \item The \textbf{axon hillock} is the short region where the axon leaves the soma. It has an unusually high density of voltage-gated sodium channels, and it is here that action potentials are typically triggered. The hillock is, in effect, the neuron's ``go / no-go'' switch.
    \item The \textbf{axon shaft} conducts the action potential along its length.
    \item The \textbf{axon terminals} (or synaptic boutons) are the tiny swellings at the far end, where neurotransmitter is released onto the next cell.
\end{itemize}

Many axons are wrapped in \textbf{myelin}, a fatty insulating sheath produced by glial cells, which dramatically speeds up conduction. The glial cells section takes this up in detail.

\subsection{Internal Machinery: Nucleus, Rough ER, and Golgi}

Neurons are unusually demanding cells. They need a constant supply of ion channels, receptors, transporters, structural proteins, and neurotransmitter-related machinery, and they need to ship all of this out to parts of the cell that may be a metre away. The internal factory is built accordingly.

\subsubsection{Nucleus}

The nucleus holds the cell's DNA and controls gene expression. In neurons, that expression has to be finely tuned: the mix of receptors and channels a neuron produces is a large part of what determines its electrical character, which transmitters it responds to, how strongly, and what it releases in turn.

Because a neuron is a permanent cell that is never replaced, this transcriptional control operates over the whole of a lifetime rather than across cell divisions. Gene expression in a neuron is also unusually responsive to the cell's own activity: patterns of firing feed back on which genes are transcribed, which is one of the routes by which experience leaves a lasting mark on the cell's properties.

\subsubsection{Rough Endoplasmic Reticulum (Nissl Substance)}

In a neuron, the rough endoplasmic reticulum is so abundant and so densely studded with ribosomes that it stains as a distinct pattern under the microscope. Historically this pattern was named \textbf{Nissl substance}. It is the neuron's protein factory: it synthesises the membrane proteins, ion channels, receptors, and secreted proteins that the cell needs in industrial quantities.

The amount of Nissl substance a neuron contains is a reasonable proxy for how metabolically active it is, and its dispersal (\textit{chromatolysis}) is a classic sign of neuronal stress after injury.

\subsubsection{Golgi Apparatus}

The Golgi apparatus takes proteins fresh from the rough ER, modifies them, sorts them, and packages them into vesicles for delivery. Its sorting role is what matters most in a neuron. A protein leaving the Golgi must be labelled for a destination, and the destinations are unusually far apart: a channel bound for a dendritic membrane, a transporter bound for the axon, and a vesicle protein bound for a terminal a metre away all leave from the same place and must not be confused. Mis-sorting in a cell this asymmetric is not a small error, since a protein delivered to the wrong compartment is both absent where it is needed and disruptive where it arrives.

\subsection{Microtubules, Kinesin, and Dynein}

Because a neuron's axon can be so long, ordinary diffusion is nowhere near fast enough to move molecules around inside it. Diffusion time scales with the square of distance, so a journey of a metre is not a hundred times slower than one of a centimetre but ten thousand times slower, on the order of decades. Neurons solve this by building a system of oriented tracks along which motor proteins actively haul their cargo.

\subsubsection{Microtubules}

The tracks are \textbf{microtubules}: rigid, polar tubes of the cytoskeleton that run the length of the axon and into the dendrites. They provide mechanical support and, crucially, a surface that motor proteins can walk along. The polarity matters, microtubules have distinct ``plus'' and ``minus'' ends, because it allows the cell to enforce direction on intracellular traffic.

\subsubsection{Kinesin and Dynein}

The motors themselves are two families of protein, and they differ in the direction they travel:

\begin{itemize}
    \item \textbf{Kinesin} walks toward the plus end, which in an axon means away from the soma and out toward the terminals. This is \textit{anterograde} transport, and it is how freshly made proteins, vesicles, and organelles are delivered to the synapse.
    \item \textbf{Dynein} walks the other way, toward the minus end and back toward the soma. This is \textit{retrograde} transport, and it is how the cell recycles components, clears damaged cargo, and, sometimes, how signals from distant synapses make it back to the nucleus.
\end{itemize}

This bidirectional traffic is non-negotiable for neuronal life. A synapse at the end of a long axon is essentially a consumer that produces very little for itself; everything it uses has to be manufactured in the soma and delivered by kinesin, and everything it must dispose of has to be carried back by dynein. Disruption of axonal transport is implicated in several neurodegenerative diseases, and the vulnerability is structural: a cell whose furthest working parts are a metre from its only factory has no redundancy if the delivery system fails.

\subsection{Vesicles and SNARE Proteins}

At the axon terminal, information has to cross from one cell to another, and for most neurons it does so chemically. Neurotransmitter is stored in advance inside small, membrane-bound \textbf{synaptic vesicles} that cluster just inside the terminal membrane, primed and ready to release.

Release has to be fast, on the scale of fractions of a millisecond, and it has to be tightly coupled to the arrival of an action potential, or else the signal would be hopelessly noisy. The machinery that accomplishes this is a family of proteins called \textbf{SNAREs}.

SNAREs on the vesicle and SNAREs on the presynaptic membrane interlock like a zip, pulling the vesicle tight against the membrane and priming it for fusion. The final trigger is a rapid influx of calcium through voltage-gated channels that open when the action potential arrives. Calcium binds to a sensor protein, the SNARE complex completes fusion, and the contents of the vesicle are dumped into the \textbf{synaptic cleft}, the tiny gap between the two neurons, in a fraction of a millisecond.

This machinery is so central to nervous-system function that several of the most potent biological toxins known target it. Botulinum toxin and tetanus toxin, for example, both cleave specific SNARE proteins; the difference in their clinical effects comes down to which synapses they disable.

\subsection{Neurotransmitters}

A \textbf{neurotransmitter} is simply a chemical messenger released from one neuron that acts on receptors belonging to another. In principle the definition is simple; in practice the list of neurotransmitters is long and the consequences are enormous, because neurotransmitters are where neuroscience starts bleeding directly into psychology and psychiatry.

At the crudest level they fall into two groups:

\begin{itemize}
    \item \textbf{Excitatory} neurotransmitters, the prototypical example being \textit{glutamate}, push the postsynaptic cell toward firing.
    \item \textbf{Inhibitory} neurotransmitters, the prototypical example being \textit{GABA}, pull it away from firing.
\end{itemize}

Most CNS signalling, by volume, is glutamate and GABA. Two further amino acids belong alongside them: \textbf{glycine}, which is the main inhibitory transmitter of the spinal cord and brainstem and works much as GABA does, opening chloride channels; and \textbf{aspartate}, which is excitatory and closely related to glutamate. These four do the bulk of fast, point-to-point signalling.

Layered on top of this are the \textbf{neuromodulators}, which typically do not simply flip a downstream cell on or off but tune the way whole circuits respond to input. These are the systems most often implicated in mood, attention, reward, and arousal, and they are also the systems most commonly targeted by psychiatric drugs. The best characterised are the \textbf{monoamines}, so called because each is built around a single amine group:

\begin{itemize}
    \item \textbf{Dopamine} governs the initiation and selection of movement, and signals reward and motivation. Its motor role is the subject of the basal ganglia section, where the same transmitter is shown to promote or suppress an action depending only on which receptor it reaches.
    \item \textbf{Serotonin} is associated with mood, sleep, and appetite, and is the target of the most widely prescribed class of antidepressant.
    \item \textbf{Noradrenaline} sets the level of arousal and alertness, driving the systems that raise vigilance in response to threat or salience.
    \item \textbf{Adrenaline} is chemically the next step beyond noradrenaline and acts mainly as a circulating hormone, though a small population of brainstem neurons use it as a transmitter in the control of blood pressure and respiration.
    \item \textbf{Histamine} is released from a small hypothalamic population that projects widely across the brain and helps maintain wakefulness. Its role is most obvious in reverse: antihistamines that cross into the brain cause drowsiness.
\end{itemize}

\textbf{Acetylcholine} sits outside the monoamine family but belongs with the modulators. It does two quite separate jobs. In the periphery it is the transmitter of the neuromuscular junction, where it activates skeletal muscle directly. In the brain it is a modulator, strongly implicated in attention and in the formation of memory.

Beyond these classical transmitters are several groups that break the mould in one way or another:

\begin{itemize}
    \item The \textbf{neuropeptides} are short chains of amino acids, larger and slower-acting than the transmitters above, and typically released alongside them rather than instead of them. They include \textbf{substance P}, which carries pain signals into the spinal cord, and the \textbf{endorphins} and \textbf{enkephalins}, the body's own opioids, which suppress those same signals and act on the receptors that morphine and its relatives exploit.
    \item The \textbf{purines}, chiefly \textbf{ATP} and its breakdown product \textbf{adenosine}, act as transmitters in their own right. Adenosine accumulates during waking and promotes sleep, which is the mechanism caffeine blocks.
    \item \textbf{Nitric oxide} is a dissolved gas rather than a packaged molecule. It cannot be stored in vesicles, so it is synthesised on demand and simply diffuses out through the membrane to act on neighbouring cells.
    \item The \textbf{endocannabinoids} are lipid messengers, and like nitric oxide they are made when needed rather than stored. Their most striking feature is the direction in which they travel, which is the subject of the final part of the next section.
\end{itemize}

Whenever we later talk about a psychological phenomenon with a pharmacological correlate, depression treated with SSRIs, Parkinson's with L-DOPA, anxiety with benzodiazepines, we will be talking, in the end, about how a drug has tweaked one of these systems.

\subsection{Summary}

A neuron's shape is not incidental. The soma integrates, the dendrites listen, the axon transmits, and the terminals release chemistry onto the next cell. Inside the cell, a network of ER, Golgi, microtubules, and motor proteins keeps this long, demanding architecture supplied with everything it needs. At the synapse, vesicles wait ready and SNAREs stand by, so that when an action potential arrives, chemical communication can happen in a fraction of a millisecond.

The feature to carry forward is the polarity. A neuron enforces one-way traffic, from dendrite to soma to hillock to terminal, and that constraint is what makes it a computational device rather than merely an excitable one. A cell that conducted equally in all directions could be a signal, but it could not be a step in an argument. Because the traffic runs one way, thousands of inputs can be summed into a single output, that output can become an input elsewhere, and elements of this kind can be chained into circuits in which the order of operations means something. Every capacity discussed in the later sections is built out of that arrangement repeated on a very large scale. With the architecture in place we are ready to ask the next question: how does a neuron actually fire?

\section{Neuron Firing}

So far we have been saying that neurons ``send signals.'' It is time to be concrete about what that actually means. A neuron does not have tiny wires or its own battery. What it has is a membrane, some ions, and some very cleverly controlled channels through which those ions can move. Everything we call neural activity, every perception, every thought, every movement, is ultimately some pattern of ions crossing membranes.

This section builds the story up from the ground. We begin with the resting state: why an idle neuron has a voltage across its membrane at all. We then see how small, local disturbances (graded potentials) can add up to a threshold crossing, at which point a stereotyped, all-or-nothing \textbf{action potential} is born. Finally we follow that action potential down the axon to the synapse, where it triggers the chemical release we just described.

\subsection{The Resting Membrane Potential}

Place a very fine electrode inside an unstimulated neuron and another in the extracellular fluid, and the voltage difference between them is roughly:

\[
V_m \approx -70 \, \text{mV}
\]

The inside of the cell is negative relative to the outside. This is not an accident, and it is not ``at rest'' in the sense of nothing happening. Maintaining the resting potential is one of the most energy-hungry things a neuron ever does.

Three facts, taken together, explain it.

First, ions are distributed unequally across the membrane. Sodium (Na\textsuperscript{+}) is much more concentrated outside the cell than inside; potassium (K\textsuperscript{+}) is much more concentrated inside than outside; chloride (Cl\textsuperscript{--}) is more concentrated outside.

Second, the membrane is selectively permeable. At rest it is much more permeable to K\textsuperscript{+} than to any other ion, because of leak channels that are always open.

Third, an active pump, the \textbf{Na\textsuperscript{+}/K\textsuperscript{+} ATPase}, burns ATP to push three sodium ions out and pull two potassium ions in with every cycle. This both maintains the concentration gradients and contributes slightly to the negative interior.

Put these three facts together and the resting membrane potential emerges more or less inevitably. Because the membrane is dominated by potassium permeability, the cell's voltage sits close to the equilibrium potential for potassium.

\subsubsection{The Nernst Equation}

We can make this quantitative. For any single ion species, there is a particular membrane voltage at which the electrical force on the ion exactly cancels the diffusional force from its concentration gradient. This voltage is the ion's \textbf{equilibrium potential}, and it is given by the \textbf{Nernst equation}:

\[
E_{ion} = \frac{RT}{zF} \ln\!\left(\frac{[ion]_{out}}{[ion]_{in}}\right)
\]

Here $R$ is the gas constant, $T$ is temperature in kelvin, $F$ is Faraday's constant, and $z$ is the ion's charge. At body temperature ($37^\circ$C), plugging in the constants and converting from the natural logarithm to the base-ten logarithm gives the convenient working form:

\[
E_{ion} \approx \frac{61}{z} \log_{10}\!\left(\frac{[ion]_{out}}{[ion]_{in}}\right) \, \text{mV}
\]

\subsubsection{Typical Equilibrium Potentials}

For a mammalian neuron under physiological conditions:

\begin{itemize}
    \item $E_K \approx -90$ mV
    \item $E_{Na} \approx +60$ mV
    \item $E_{Cl} \approx -65$ mV
\end{itemize}

Notice where $-70$ mV sits in that list. The resting potential is close to $E_K$ because at rest, potassium is effectively the ion whose voice dominates. If a cell were instead to open a large number of sodium channels, the voltage would race toward $+60$ mV instead. This is exactly what happens during an action potential.

\subsection{Graded Potentials}

Before we reach the action potential, we need a smaller kind of signal: the \textbf{graded potential}. Graded potentials are the small, local changes in membrane voltage produced in dendrites and the soma by incoming synaptic input, and they are the raw material out of which the decision to fire is built.

Three properties define them. They are \textit{graded}: their size depends on the size of the stimulus, not an all-or-nothing switch. They \textit{decay} with distance: they are passive spreads of voltage, and they get weaker the further they travel. And they \textit{summate}: if several graded potentials arrive at roughly the same time (temporal summation) or at nearby spots on the dendritic tree (spatial summation), their effects add together.

But where does the voltage change actually come from? A neurotransmitter molecule is not electrically interesting in itself; it carries no current across the membrane and it does not touch the ions on either side. Everything depends on what the transmitter binds to. Postsynaptic receptors come in two architecturally different kinds, and the difference between them accounts for most of the variety in how neurons influence one another.

\subsubsection{Ligand-Gated Ion Channels}

The first kind is the \textbf{ligand-gated ion channel}, also called an \textbf{ionotropic} receptor. Here the receptor \textit{is} the channel: a single protein assembly that spans the membrane, encloses a water-filled pore, and carries a binding site for the transmitter on its outer face. When the transmitter (the \textit{ligand}) binds, the protein changes shape, the pore opens, and ions move through it under the forces described in the previous subsection. When the transmitter unbinds, the pore shuts.

Nothing else is involved. There is no intermediate messenger and no chain of enzymes, which is why this arrangement is fast: the delay between transmitter binding and current flowing is a fraction of a millisecond, and the whole event is usually over within a few milliseconds. It is also strictly local, confined to the patch of membrane holding the channel.

What determines whether such a channel excites or inhibits is not the transmitter but the \textbf{selectivity of the pore}, that is, which ions it will let through:

\begin{itemize}
    \item \textbf{Glutamate} acts on \textbf{AMPA} and \textbf{NMDA} receptors, both of which pass cations, chiefly Na\textsuperscript{+} inward. The result is depolarisation. This is why glutamate is the workhorse excitatory transmitter.
    \item \textbf{Acetylcholine} acts on \textbf{nicotinic} receptors, which likewise pass cations inward, and which are the receptors of the neuromuscular junction.
    \item \textbf{GABA} acts on \textbf{GABA\textsubscript{A}} receptors, which pass Cl\textsuperscript{--}, and \textbf{glycine} acts on receptors that do the same. The result is inhibition.
\end{itemize}

It is worth stating plainly, because it is a common confusion, that no transmitter is intrinsically excitatory or inhibitory. Glutamate is excitatory because the receptors that happen to bind it happen to pass sodium. Put glutamate on a chloride-selective channel and it would inhibit. The transmitter selects which receptors respond; the receptor decides what that response is.

\subsubsection{G Protein-Coupled Receptors}

The second kind is the \textbf{G protein-coupled receptor (GPCR)}, also called a \textbf{metabotropic} receptor, and it works on an entirely different principle. A GPCR is not a channel and cannot pass a current at all. It is a single protein that threads back and forth through the membrane seven times, which is why the family is sometimes called the seven-transmembrane receptors, with its binding site facing outward and its business end facing the cytoplasm. What it does when activated is start a chain of chemistry inside the cell.

The sequence runs as follows.

\begin{enumerate}
    \item The transmitter binds to the receptor's extracellular face, and the protein changes shape.
    \item That shape change is felt on the inside, where the receptor is associated with a \textbf{G protein}, a three-part assembly of $\alpha$, $\beta$, and $\gamma$ subunits. At rest the $\alpha$ subunit holds a molecule of \textbf{GDP}.
    \item The activated receptor causes the $\alpha$ subunit to swap that GDP for \textbf{GTP}. This is the switch being thrown.
    \item The GTP-bound $\alpha$ subunit separates from the $\beta\gamma$ pair, and both halves are now free to act.
    \item The $\alpha$ subunit acts on an \textbf{effector enzyme}, commonly \textbf{adenylyl cyclase}, which manufactures a \textbf{second messenger}, \textbf{cyclic AMP (cAMP)}, that diffuses through the cytoplasm.
    \item cAMP activates \textbf{protein kinase A (PKA)}, an enzyme that phosphorylates other proteins, including ion channels already sitting in the membrane. Phosphorylation changes how readily those channels open.
    \item The $\alpha$ subunit eventually hydrolyses its GTP back to GDP, reassembles with $\beta\gamma$, and the system resets.
\end{enumerate}

Which direction the cascade pushes depends on which G protein the receptor is coupled to. Receptors coupled to \textbf{G\textsubscript{s}} (stimulatory) activate adenylyl cyclase, so cAMP rises and PKA is switched on. Receptors coupled to \textbf{G\textsubscript{i}} (inhibitory) suppress it, so cAMP falls. The same receptor architecture and the same transmitter can therefore produce opposite effects in two different cells, depending only on which G protein sits behind the receptor.

The $\beta\gamma$ pair is not merely left over. It can act on its own, and in one case that matters a great deal it acts directly on ion channels: $\beta\gamma$ released by GABA\textsubscript{B} and by various monoamine receptors binds and opens a class of potassium channels known as \textbf{GIRK} channels. Potassium leaves, the cell hyperpolarises, and the membrane voltage has been changed by a metabotropic receptor without any second messenger being involved at all. This is a useful corrective to the tidy picture: metabotropic does not mean slow-and-indirect-only, it means the receptor and the channel are separate proteins.

Three consequences of the whole arrangement matter. It is \textit{slow}: a cascade of several enzymatic steps takes hundreds of milliseconds to seconds, against the millisecond of an ionotropic synapse. It is \textit{amplifying}: one bound receptor activates many G proteins, each activated enzyme makes many molecules of second messenger, and so a small amount of transmitter produces a large intracellular effect. And it is \textit{modulatory} rather than commanding: the usual end result is not that the cell fires, but that the channels it already has are tuned to open more or less readily. A GPCR does not so much send a message as change how the cell reads its other messages.

Most transmitters act on both kinds of receptor. Glutamate has metabotropic receptors alongside AMPA and NMDA; GABA has GABA\textsubscript{B} alongside GABA\textsubscript{A}; acetylcholine has muscarinic receptors alongside nicotinic. The neuromodulators met in the previous section, dopamine, serotonin, and noradrenaline, act almost entirely through GPCRs, which is precisely why they modulate rather than command.

\subsubsection{How These Change the Membrane Voltage}

We can now be exact about the step that both mechanisms end in: an open channel, and a voltage that moves. The key idea is that opening a channel does not \textit{set} a voltage. It passes a \textbf{current}, and the size and direction of that current depend on how far the membrane currently sits from that ion's equilibrium potential:

\[
I_{ion} = g_{ion}\,(V_m - E_{ion})
\]

Here $g_{ion}$ is the conductance, essentially how many channels for that ion are open, and the bracketed term $(V_m - E_{ion})$ is the \textbf{driving force}. The equation says something simple: an ion flows in whichever direction moves the membrane potential \textit{toward} that ion's equilibrium potential, and it flows harder the further away the membrane currently is. An open channel is a pull toward its own ion's $E_{ion}$, and the strength of the pull is the distance still to travel.

Read against the equilibrium potentials tabulated earlier, every kind of postsynaptic response follows directly:

\begin{itemize}
    \item Open a \textbf{sodium}-permeable channel at rest. The membrane sits at $-70$ mV and $E_{Na} \approx +60$ mV, so the driving force is large and inward. Sodium pours in, the membrane depolarises toward $+60$ mV, and the result is an \textbf{EPSP}.
    \item Open a \textbf{potassium}-permeable channel, as $\beta\gamma$ does at GIRK channels. Since $E_K \approx -90$ mV lies below rest, potassium leaves, the membrane hyperpolarises toward $-90$ mV, and the result is an \textbf{IPSP}.
    \item Open a \textbf{chloride}-permeable channel, as GABA\textsubscript{A} does. Here $E_{Cl} \approx -65$ mV sits very close to rest, so remarkably little chloride actually moves and the voltage barely changes. The inhibition is real nonetheless: with a large chloride conductance open, the membrane is effectively clamped near $-65$ mV, and any excitatory current arriving has to fight that clamp. This is \textbf{shunting inhibition}, which holds the cell down rather than pushing it down.
\end{itemize}

This also explains the defining properties listed at the start of the subsection. The response is \textit{graded} because more transmitter binds more receptors, which opens more channels, which is a larger $g_{ion}$ and so a larger current and a bigger deflection; nothing here is all-or-nothing. It \textit{decays} with distance because the charge spreading away from the synapse leaks back out across the membrane as it goes, so less and less of it reaches the axon hillock. And it \textit{summates} because these are currents flowing into a common pool of charge, and currents simply add.

Finally, the two receptor families leave different signatures on the voltage. An ionotropic receptor produces a fast, sharp deflection lasting a few milliseconds, the discrete blip of a single synaptic event. A metabotropic receptor produces a slow shift, taking hundreds of milliseconds to develop and persisting for seconds, which does not so much signal an event as change the terms on which every other input is judged. When we come to the basal ganglia we will find dopamine doing exactly this, setting the gain on a circuit rather than driving it.

With that machinery in place, the two kinds of graded potential can be named:

\begin{itemize}
    \item \textbf{Excitatory postsynaptic potentials (EPSPs)} depolarise the cell, moving it closer to threshold.
    \item \textbf{Inhibitory postsynaptic potentials (IPSPs)} hyperpolarise the cell, moving it away from threshold.
\end{itemize}

The axon hillock tallies the EPSPs and IPSPs arriving at every moment, and when the net result is a depolarisation above threshold, an action potential is launched.

\subsection{The Action Potential}

The action potential is a rapid, stereotyped, all-or-nothing spike in membrane voltage that propagates along the axon without decay. It is the only signal a neuron has that can travel the full length of the cell and arrive undiminished.

\subsubsection{Threshold}

The action potential begins when the membrane at the axon hillock depolarises to about:

\[
V_{\text{threshold}} \approx -55 \, \text{mV}
\]

At this voltage, voltage-gated sodium channels open in large numbers, and what follows is a runaway.

\subsubsection{Phases of the Action Potential}

The spike unfolds as a short sequence of stages, each dominated by a different set of channels.

\textbf{1. Resting state} ($\approx -70$ mV). Voltage-gated channels are closed. Leak channels and the Na\textsuperscript{+}/K\textsuperscript{+} pump maintain the resting potential.

\textbf{2. Depolarisation} (up to $\approx +30$ mV). Voltage-gated Na\textsuperscript{+} channels open in a positive-feedback avalanche: more sodium in means more depolarisation, which means more channels open. The voltage races upward toward $E_{Na}$.

\textbf{3. Repolarisation.} Two things happen, slightly delayed. Sodium channels \textit{inactivate}, they physically plug themselves closed with an intrinsic gate, and voltage-gated potassium channels open. Potassium flows out and drags the voltage back down.

\textbf{4. Hyperpolarisation} ($\approx -80$ to $-90$ mV). Potassium channels are slow to shut, so K\textsuperscript{+} efflux briefly overshoots the resting potential, pulling the cell closer to $E_K$.

\textbf{5. Return to rest.} Potassium channels close. The Na\textsuperscript{+}/K\textsuperscript{+} pump continues its unending job of restoring the gradients, and the membrane settles back to $-70$ mV.

The entire spike, from threshold to recovery, takes only a few milliseconds. And because it is triggered by a self-reinforcing avalanche, it is all-or-nothing: a neuron either fires or it does not. There is no such thing as a partial action potential.

\subsection{Neurotransmitter Release at the Synapse}

When the action potential reaches the axon terminal, it finds yet another class of voltage-gated channels: \textbf{calcium channels}. These open, calcium floods in, and, as we saw in the previous section, that calcium is the trigger the SNARE machinery has been waiting for. Vesicles fuse with the presynaptic membrane and dump their neurotransmitter into the synaptic cleft, where it diffuses across and binds to receptors on the next cell.

The amount of neurotransmitter released depends on how much calcium enters, so the strength of the chemical signal is tightly linked to the arrival of the electrical one. In this way a digital, all-or-nothing spike is translated back into a graded chemical signal, which becomes a new graded potential in the next neuron, and the whole cycle repeats.

\subsection{Clearing the Signal: Reuptake and Degradation}

A signal is only useful if it can also be turned off. Neurotransmitter that is left to linger in the synaptic cleft blurs subsequent signals and, in the case of excitatory transmitters, can actually damage the postsynaptic cell (a phenomenon called \textit{excitotoxicity}). There are three main ways the cleft is cleaned up:

\begin{itemize}
    \item \textbf{Reuptake.} Specialised transporters on the presynaptic neuron (or on nearby glia) pump neurotransmitter back out of the cleft, often to be repackaged for reuse. This is the dominant mechanism for monoamines like serotonin and dopamine.
    \item \textbf{Enzymatic degradation.} At some synapses an enzyme in the cleft breaks the transmitter down. The classic example is \textit{acetylcholinesterase}, which cleaves acetylcholine almost as soon as it is released.
    \item \textbf{Diffusion.} Simple loss of molecules out of the narrow cleft also contributes.
\end{itemize}

Many psychoactive drugs act here. Selective serotonin reuptake inhibitors (SSRIs) block the reuptake transporter for serotonin, raising its concentration in the cleft. Amphetamines interfere with dopamine reuptake and release. Nerve agents block acetylcholinesterase, producing catastrophic overstimulation. Whenever a drug is said to ``increase'' or ``decrease'' a neurotransmitter, it is almost always doing something very concrete and very local at the synapse.

\subsection{Refractory Periods}

After an action potential, the neuron is not instantly ready to fire another. For a brief interval it refuses absolutely, and for a slightly longer interval it only reluctantly agrees. These intervals are the refractory periods, and they are not bugs, they are what keep nervous-system signalling clean.

\subsubsection{Absolute Refractory Period}

Immediately after a spike, no new action potential can be generated at any price. This is because the voltage-gated sodium channels that just fired are \textit{inactivated}: their inactivation gate is closed, and it will not reopen until the membrane has returned close to rest. No amount of stimulation will fire them. The absolute refractory period effectively sets an upper bound on how fast a neuron can fire.

\subsubsection{Relative Refractory Period}

During the hyperpolarisation that follows, sodium channels have recovered but the cell is further from threshold than usual. Firing is possible, but it takes a stronger-than-normal stimulus to get there, and the deficit shrinks as the membrane drifts back toward rest. The consequence is that the relative refractory period does not block firing so much as price it: the strength of input required to trigger a second spike is a direct measure of how recently the first one occurred. Only an unusually strong input will fire the cell early in this window, so a neuron's firing rate ends up reflecting the intensity of what is driving it.

Together, these two periods do two crucial jobs. They ensure that action potentials propagate \textbf{in one direction only}, the region just behind a travelling spike is refractory, so the spike cannot double back, and they cap the maximum firing rate, forcing the nervous system to encode intensity through \textit{rate} and \textit{pattern} rather than amplitude.

\subsection{Retrograde Signalling: Endocannabinoids}

Everything described so far in this section runs in one direction. The presynaptic terminal releases, the postsynaptic membrane receives, and the polarity of the neuron enforces the arrangement from dendrite to hillock to terminal. There is, however, a class of messenger that runs the other way, from the postsynaptic cell back to the presynaptic terminal that has just signalled to it. The best understood are the \textbf{endocannabinoids}, so named because they act on the receptors that the active compounds of cannabis also bind.

Two of them account for most of what is known: \textbf{anandamide} and \textbf{2-arachidonoylglycerol (2-AG)}. Both are lipids, derived from fatty acid components of the cell membrane itself, and that chemistry dictates almost everything else about how they behave.

\subsubsection{Synthesis on Demand}

A conventional transmitter is manufactured in advance, loaded into vesicles, and held at the terminal awaiting a calcium signal. Endocannabinoids cannot be handled this way. Being lipids, they would dissolve straight through the wall of any vesicle intended to contain them, so there is nothing to store and nothing to release in the usual sense.

Instead they are synthesised at the moment they are required. When a postsynaptic neuron is strongly depolarised, or when certain of its metabotropic receptors are activated, calcium entering the cell switches on the enzymes that cleave endocannabinoids out of membrane lipids. The molecule is produced on the spot, in the postsynaptic membrane, in response to that cell's own activity.

\subsubsection{Travelling Backwards}

Having been made in the postsynaptic membrane, the endocannabinoid diffuses the short distance back across the synaptic cleft, in the opposite direction to the traffic described in the rest of this section. On the presynaptic terminal it binds \textbf{CB1 receptors}, which are among the most abundant G protein-coupled receptors in the brain.

The consequence is inhibitory. CB1 is coupled to G\textsubscript{i}, and its activation suppresses the calcium channels on which vesicle release depends, described earlier as the trigger the SNARE machinery waits for. With calcium entry reduced, fewer vesicles fuse and less transmitter is released. The postsynaptic cell has, in effect, reached back across the synapse and instructed the presynaptic terminal to be quieter.

\subsubsection{Why a Backwards Signal Is Useful}

The arrangement gives a synapse something it otherwise lacks: a way for the receiving cell to regulate how much it is being sent. Every other mechanism in this section leaves the postsynaptic neuron a passive recipient, able to decide only what to do with the input it is given. Retrograde signalling turns that into a negotiation.

Its clearest use is as a brake against overexcitation. A postsynaptic neuron being driven hard produces endocannabinoids in proportion to that drive, and the resulting suppression of transmitter release turns the input down. The effect can be brief, lasting seconds, or sustained, and in the latter case it becomes one of the mechanisms by which the strength of a synapse is lastingly adjusted, rather than merely a moment-to-moment correction.

This also explains why the cannabinoid system has such diffuse effects when it is engaged from outside. A drug acting at CB1 does not target one circuit; it interferes with a general-purpose feedback mechanism operating at synapses throughout the brain, which is why the consequences reach across appetite, pain, mood, movement, and memory at once.

\subsection{Autoreceptors}

Endocannabinoids are not the only brake operating on the presynaptic terminal, and the other one comes from a different direction entirely. An \textbf{autoreceptor} is a receptor sitting on a neuron's own membrane that binds the very transmitter that neuron releases. The cell is, in effect, listening to itself: what it has just released comes back and binds to it, and the result is a negative feedback loop closed entirely within a single neuron. This distinguishes autoreceptors from \textbf{heteroreceptors}, which sit in the same presynaptic position but respond to some other transmitter released by some other cell. It also distinguishes them from the retrograde signalling just described, where the message originates in the postsynaptic cell. Endocannabinoids are the receiving neuron asking for less; an autoreceptor is the sending neuron regulating itself.

The mechanism divides according to where on the neuron the receptor sits, and the two locations do different jobs.

\textbf{Terminal autoreceptors} sit on the axon terminal alongside the release machinery. When transmitter accumulates in the cleft, these receptors detect it and suppress the voltage-gated calcium channels that vesicle fusion depends on, by exactly the route CB1 uses. Less calcium enters with each arriving action potential, fewer vesicles fuse, and the amount of transmitter released per impulse falls. The neuron continues to fire at the same rate; it simply says less each time it does.

\textbf{Somatodendritic autoreceptors} sit at the other end of the cell, on the soma and dendrites, where they detect transmitter released from the neuron's own dendrites or from neighbouring cells of the same population. Here the effect is on firing itself. These receptors act through the $\beta\gamma$ pathway described earlier, opening GIRK potassium channels; potassium leaves, the membrane hyperpolarises toward $E_K$, and the cell is carried further from threshold. The neuron now fires less often. Between the two locations, a neuron can throttle its output either by speaking less often or by saying less each time, and it is the somatodendritic population that regulates firing rate proper.

Both kinds are almost always metabotropic, and almost always coupled to G\textsubscript{i/o}. This follows from the job rather than being an accident of which receptors happened to be recruited. A feedback brake needs to be slow enough to reflect accumulated activity rather than a single vesicle, amplified enough that modest transmitter concentrations produce a real effect, and modulatory rather than commanding, since its purpose is to adjust the cell's responsiveness and not to drive it. Those are precisely the three properties of G protein-coupled signalling set out earlier, and they are what an ionotropic receptor could not provide.

The particular receptors involved are familiar ones. Dopamine neurons carry \textbf{D2} autoreceptors, noradrenergic neurons carry \textbf{$\alpha$2}, serotonergic neurons carry \textbf{5-HT\textsubscript{1A}} on their cell bodies and \textbf{5-HT\textsubscript{1B/1D}} on their terminals, and GABAergic terminals carry \textbf{GABA\textsubscript{B}}. Each of these subtypes also occurs postsynaptically elsewhere in the brain, doing quite different work. What makes a receptor an autoreceptor is therefore its location, not its identity: the same molecule is a feedback sensor on the cell that releases the transmitter and an ordinary receiver on the cell that listens.

The clinical consequences of this are best seen in the delayed action of the SSRIs. As noted earlier, these drugs block the serotonin reuptake transporter, and that blockade is immediate: within hours of the first dose, serotonin is clearing from the cleft more slowly and its concentration is raised. Yet patients do not feel better for weeks, which is a genuine puzzle if raised serotonin were the whole story. The autoreceptor resolves it. Reuptake is blocked everywhere at once, including around the raphe nuclei where the serotonergic cell bodies live, and the resulting rise in serotonin there activates their somatodendritic 5-HT\textsubscript{1A} autoreceptors. Those receptors do what autoreceptors do, and the raphe neurons slow their firing. Less serotonin therefore arrives at the distant terminals where the therapeutic effect is thought to arise, and the two changes very nearly cancel: the drug is working perfectly at the molecule it targets while the system as a whole absorbs the change.

What breaks the deadlock is adaptation. Under sustained stimulation the somatodendritic autoreceptors gradually desensitise and reduce in number over a period of two to four weeks. As their restraint weakens, raphe firing returns toward normal, but the reuptake blockade is still in place, so terminal serotonin now genuinely rises for the first time. The clinical improvement tracks this slow adaptation rather than the drug's immediate pharmacology, which is why the delay is measured in weeks and why stopping too early can look like failure when it is only earliness.

The general lesson runs beyond antidepressants. Because an autoreceptor and its postsynaptic counterpart are frequently the same molecule, a drug directed at one will ordinarily reach the other, and the effect observed in a patient depends on which site dominates and on how each adapts over time. A receptor system is not a fixed target but one that responds to being interfered with.

\subsection{Neurotransmitter Mimicry}

Almost everything in this section has been described in terms of endogenous molecules: the transmitters the brain makes, the receptors they bind, and the machinery that clears them away. But a receptor cannot tell where the molecule filling it came from. It responds to shape and chemistry, not provenance, and this indifference is the opening that pharmacology exploits. A great many drugs, from the ones a psychiatrist prescribes to the ones sold on a street corner, work for the simple reason that they resemble a neurotransmitter closely enough to be mistaken for it, or resemble it closely enough to sit in its place without doing its job.

It is worth being precise about the vocabulary, because the differences between these cases are exactly the differences between the drugs. A molecule that binds a receptor's normal transmitter site, the \textbf{orthosteric} site, and switches the receptor on as the natural transmitter would is an \textbf{agonist}. This is mimicry in the strict sense: the drug is a counterfeit key that turns the lock. A molecule that occupies the same site but fails to activate the receptor is an \textbf{antagonist}; it turns nothing, and by squatting in the site it keeps the real transmitter out, so its effect is blockade rather than imitation. Between these lies the \textbf{partial agonist}, which activates the receptor but only submaximally, so that against a background of high natural signalling it can even look like a brake. And a molecule need not touch the orthosteric site at all: an \textbf{allosteric modulator} binds elsewhere on the receptor and changes how the receptor answers its true transmitter, turning the response up (a positive modulator) or down (a negative one) without ever triggering it directly. Only the first of these is mimicry; the others are worth setting beside it precisely because they are so often confused with it.

\subsubsection{Serotonin: LSD and the Psychedelics}

The canonical example is \textbf{LSD}. Its behavioural effects are extraordinary, but its mechanism is mundane in the sense that matters here: the molecule is close enough to serotonin to bind serotonin's receptors, and in particular it is an agonist at the \textbf{5-HT\textsubscript{2A}} receptor, a G protein-coupled receptor concentrated on the pyramidal cells of the cortex. \textbf{Psilocybin}, the active principle of the relevant mushrooms, does the same thing by way of \textbf{psilocin}, the compound it is converted into in the body. Both are partial agonists, and the cortical 5-HT\textsubscript{2A} receptor is now the agreed common denominator of the classical psychedelics: block it and the experience does not occur.

Note how specific this is. Serotonin acts across a whole family of receptors, among them the 5-HT\textsubscript{1A} autoreceptor met in the previous subsection, whose job is to quiet the raphe neurons that release serotonin in the first place. LSD's psychedelic action is not a matter of raising serotonin everywhere; it is a matter of pretending to be serotonin at one particular subtype in one particular place. The lesson that receptor identity and location decide the effect, already drawn for autoreceptors, is drawn again here from the other side.

\subsubsection{Acetylcholine: Nicotine and Muscarine}

Acetylcholine offers the cleanest teaching example, because two natural products carve its two receptor families neatly in half and even supply their names. \textbf{Nicotine} is an agonist at the \textbf{nicotinic} acetylcholine receptor, the ionotropic, cation-passing receptor of the neuromuscular junction described earlier. Being ionotropic, it acts fast, and because these receptors also sit on the dopamine-releasing neurons of the reward pathway, nicotine's agonism there is a large part of why it is addictive. \textbf{Muscarine}, from the fly agaric mushroom, is an agonist at the other family, the metabotropic \textbf{muscarinic} receptors; and \textbf{atropine}, from deadly nightshade, is a muscarinic \textbf{antagonist}. Here the whole vocabulary is on display in a single transmitter: two agonists selective for two different receptor families, and an antagonist that blocks one of them, which is why atropine dries secretions and races the heart, doing the opposite of what muscarinic activation would.

\subsubsection{Opioids: Morphine and Heroin}

The body makes its own opioids, the \textbf{endorphins} and \textbf{enkephalins} met among the neuropeptides, which damp down pain signalling. \textbf{Morphine}, and \textbf{heroin} once the body converts it into morphine, are agonists at the \textbf{$\mu$-opioid receptor} those peptides normally use, a G protein-coupled receptor working through G\textsubscript{i}. This is straightforward mimicry of an endogenous transmitter, and its two faces, powerful analgesia and powerful addiction, follow from the same fact: the drug occupies both the pain-suppressing circuits and, as with nicotine, the reward circuitry, more completely and for longer than the body's own opioids ever would.

\subsubsection{Cannabinoids: THC and CBD}

The cannabis plant furnishes the neatest contrast in the whole section, because its two best-known constituents act in opposite manners. \textbf{THC} (tetrahydrocannabinol) is a partial agonist at the \textbf{CB1 receptor}, the presynaptic G protein-coupled receptor introduced when endocannabinoids were discussed. It is, quite literally, a mimic of anandamide and 2-AG: it engages the brain's own retrograde-signalling brake from outside, which is exactly why its effects are so diffuse, reaching appetite, mood, memory, and pain at once, as anticipated at the end of the endocannabinoid subsection.

\textbf{CBD} (cannabidiol) is the instructive exception. It is not an appreciable CB1 agonist and does not produce the intoxication THC does; instead it behaves as a \textbf{negative allosteric modulator} at CB1, tuning down that receptor's response rather than driving it, and it exerts much of its remaining action elsewhere, including agonism at the \textbf{5-HT\textsubscript{1A}} receptor. Two molecules from the same plant, one an agonist and one an allosteric modulator, thus illustrate in a single example the difference the opening paragraphs laboured to establish.

\subsubsection{The Non-Mimics: Caffeine and Alcohol}

It is worth ending with two of the most widely consumed psychoactive drugs in the world, precisely because neither is a mimic, and seeing why sharpens the definition. \textbf{Caffeine} works chiefly as an \textbf{antagonist} at \textbf{adenosine} receptors, the \textbf{A\textsubscript{1}} and \textbf{A\textsubscript{2A}} subtypes. Adenosine accumulates through the waking day and, acting on these receptors, promotes sleepiness; caffeine occupies the receptors without activating them, so the drowsiness signal is blocked. Caffeine does not imitate a transmitter, then. It removes a brake, and alertness is what is left when the brake is gone.

\textbf{Alcohol} is diffuse in a different way, acting not at one clean site but across several, and again by modulation rather than mimicry. Its two principal actions are as a \textbf{positive allosteric modulator} of the \textbf{GABA\textsubscript{A}} receptor, the ionotropic chloride channel that is the brain's main inhibitory receptor, potentiating inhibition and producing sedation, and as an \textbf{antagonist} at the \textbf{NMDA} glutamate receptor, damping excitation. The \textbf{benzodiazepines} belong beside alcohol here as the tidier case: they too are positive allosteric modulators of GABA\textsubscript{A}, binding their own distinct site to make the receptor more responsive to the GABA already present, which is why they calm without directly opening the channel themselves. The general point closing the section is that ``acting at a receptor'' covers three quite different things, agonism, antagonism, and allosteric tuning, and only the first deserves the name of mimicry; the confusion of the three is the confusion of a drug that imitates a signal with one that merely blocks or bends it.

\subsection{Summary}

The electrical life of a neuron is a story of ions crossing a membrane. The resting potential is set mostly by potassium permeability. Small graded potentials from synaptic input summate at the axon hillock; if they push the membrane past threshold, a self-reinforcing avalanche of sodium entry produces an action potential, which propagates to the axon terminal. There, calcium influx triggers vesicle fusion, neurotransmitter is released, and the process starts again in the next cell. Refractory periods keep the signal unidirectional and bounded. Layered on top of this forward flow are two presynaptic brakes. Endocannabinoids, synthesised on demand in the postsynaptic membrane, travel back across the cleft to suppress release from the terminal that just fired, giving the receiving cell a say in how much it is sent; autoreceptors let the sending neuron regulate itself, damping release at the terminal and firing rate at the soma according to how much transmitter it has already put out. These same receptors are also where drugs act, mimicking the natural signal as agonists, blocking it as antagonists, or allosterically tuning the response up or down, and the effect a drug produces follows from which of these it does and at which receptor. The Nernst equation, which at first looks like a piece of physical chemistry, turns out to be quietly responsible for most of it.

One result of this section deserves to be carried further than the others, because it is the first point at which anything resembling \textit{meaning} enters the account. A neurotransmitter has no intrinsic significance. Glutamate is not an excitatory molecule and GABA is not an inhibitory one; glutamate excites because the receptors that bind it happen to pass sodium, and put on a chloride-selective channel it would inhibit. The signal is identical in every case, and what the signal means is decided entirely by the receiver. This is a small fact about pore selectivity, and it is also the seed of something much larger, because it establishes that significance in a nervous system is not carried by a message but assigned by whatever the message reaches. The same principle appears twice more in these notes and does progressively heavier work each time: at the autoreceptor, where a receptor's location rather than its identity determines whether it is a feedback sensor or an ordinary receiver, and in the basal ganglia, where one diffuse broadcast of dopamine promotes movement and suppresses it simultaneously, because the two populations it lands on carry receptors of opposite sign. A system built this way can make one chemical mean many things, which is a large part of how so much is achieved with so short a list of transmitters.

The transmitters and receptor families met across this section are gathered in Table~\ref{tab:transmitters} below, extended with their principal receptor subtypes to serve as a single-glance reference.

{\footnotesize
\renewcommand{\arraystretch}{1.3}
\begin{longtable}{p{2.8cm} p{2.4cm} p{1.8cm} p{1.9cm} p{4.0cm}}
\caption{The principal neurotransmitters and the receptor families through which they act, with their synaptic actions and the drugs that bind them. ``Type'' distinguishes fast ionotropic channels from slower metabotropic G protein-coupled receptors (GPCRs), with the G protein noted where relevant. The receptors drawn on in this section are discussed in the prose above; the remaining subtypes are included to complete the picture. A dash marks a receptor for which no relevant drug is listed.}
\label{tab:transmitters} \\
\toprule
\textbf{Neurotransmitter} & \textbf{Receptor} & \textbf{Type} & \textbf{Action} & \textbf{Drug (effect)} \\
\midrule
\endfirsthead
\multicolumn{5}{l}{\emph{\small Table~\ref{tab:transmitters} continued from previous page}} \\
\toprule
\textbf{Neurotransmitter} & \textbf{Receptor} & \textbf{Type} & \textbf{Action} & \textbf{Drug (effect)} \\
\midrule
\endhead
\bottomrule
\endfoot
\bottomrule
\endlastfoot
Glutamate & AMPA & Ionotropic & Excitatory & --- \\
 & NMDA & Ionotropic & Excitatory & Alcohol (antagonist) \\
 & Kainate & Ionotropic & Excitatory & --- \\
 & mGluR & GPCR & Modulatory & --- \\
\midrule
GABA & GABA\textsubscript{A} & Ionotropic & Inhibitory & Alcohol, benzodiazepines (positive allosteric modulator) \\
 & GABA\textsubscript{B} & GPCR & Inhibitory & --- (autoreceptor) \\
\midrule
Acetylcholine & Nicotinic & Ionotropic & Excitatory & Nicotine (agonist) \\
 & Muscarinic & GPCR & Modulatory & Muscarine (agonist); atropine (antagonist) \\
\midrule
Dopamine & D1 & GPCR (G\textsubscript{s}) & Excitatory & --- \\
 & D2 & GPCR (G\textsubscript{i}) & Inhibitory & --- (autoreceptor) \\
\midrule
Serotonin (5-HT) & 5-HT\textsubscript{1A} & GPCR & Inhibitory & CBD (agonist); somatodendritic autoreceptor \\
 & 5-HT\textsubscript{1B/1D} & GPCR & Inhibitory & --- (terminal autoreceptor) \\
 & 5-HT\textsubscript{2A} & GPCR & Excitatory & LSD, psilocin (partial agonist) \\
 & 5-HT\textsubscript{3} & Ionotropic & Excitatory & Ondansetron (antagonist) \\
\midrule
Noradrenaline & $\alpha$\textsubscript{1} & GPCR (G\textsubscript{q}) & Excitatory & --- \\
 & $\alpha$\textsubscript{2} & GPCR & Inhibitory & --- (autoreceptor) \\
 & $\beta$\textsubscript{1}/$\beta$\textsubscript{2} & GPCR (G\textsubscript{s}) & Modulatory & --- \\
\midrule
Endocannabinoids (anandamide, 2-AG) & CB1 & GPCR (G\textsubscript{i}) & Inhibitory & THC (partial agonist); CBD (negative allosteric modulator) \\
 & CB2 & GPCR (G\textsubscript{i}) & Modulatory & --- \\
\midrule
Opioid peptides (endorphins, enkephalins) & $\mu$-opioid & GPCR (G\textsubscript{i}) & Inhibitory & Morphine, heroin (agonist) \\
\midrule
Adenosine & A\textsubscript{1}/A\textsubscript{2A} & GPCR & Modulatory & Caffeine (antagonist) \\
\end{longtable}
}

\section{Glial Cells}

Read a textbook written fifty years ago and it might describe neurons as the stars of the nervous system and glial cells as their stagehands, quiet, supportive, uninteresting. That view is now dead. Modern neuroscience treats glia as active participants: they shape synapses, regulate ions, feed neurons, defend the brain, build the insulation that makes fast signalling possible, and fail or misbehave in many of the diseases we most dread. A surprisingly large portion of what we might loosely call ``brain function'' is not neuronal at all.

The word ``neuroglia'' literally means ``nerve glue,'' which was a reasonable guess in the nineteenth century and is a misleading one now. Structural support turns out to be among the least of what these cells do.

This section introduces the main types of glial cell, what each of them does, and where glial failure sits in neurological disease. The theme running through all of it is that a neuron cannot do its job alone: every neuron in the body depends on glia to regulate its environment, recycle its transmitters, insulate its axon, and protect it from damage.

The main glial types, divided by location, are:

\begin{itemize}
    \item \textbf{CNS}: astrocytes, oligodendrocytes, ependymal cells, microglia.
    \item \textbf{PNS}: Schwann cells, satellite cells.
\end{itemize}

\subsection{Astrocytes and Satellite Cells}

\subsubsection{Astrocytes}

Astrocytes are the workhorses of the CNS glial family. Under the microscope they are star-shaped (hence the name), with long processes radiating in every direction. Those processes contact neurons, synapses, blood vessels, and the surface of the brain, and that connectivity is the key to everything they do. An astrocyte sits at the intersection of almost everything important in nervous tissue.

Their responsibilities divide into seven roles.

\paragraph{1. Structural support.} Astrocytes hold the architecture of nervous tissue together. Their processes fill the spaces between neurons, dividing the tissue into domains that barely overlap, so that each astrocyte takes responsibility for a defined territory and its synapses. Nervous tissue has no fibrous connective framework of the kind found elsewhere in the body; this cellular scaffolding is what stands in for it.

\paragraph{2. Ion buffering.} Every time a neuron fires, potassium leaves the cell and accumulates briefly in the extracellular space. If this were allowed to build up, nearby neurons would become abnormally depolarised and fire uncontrollably. Astrocytes take up the excess potassium and redistribute it through their own network, a process called \textbf{spatial buffering}. Without it, neuronal activity would rapidly become chaotic.

\paragraph{3. Neurotransmitter recycling.} Astrocytes are the main clean-up crew for glutamate, the principal excitatory transmitter of the brain. They pull glutamate out of the cleft through dedicated transporters, convert it to glutamine (which is inert and therefore safe), and ship the glutamine back to neurons to be re-converted into glutamate. This loop is called the \textbf{glutamate--glutamine cycle}, and it is the main defence against excitotoxicity.

\paragraph{4. Metabolic support.} Neurons have extraordinary energy demands and almost no energy storage, so someone has to keep them fed. That someone is the astrocyte. Astrocytes take up glucose from blood vessels, store glycogen, and supply lactate and other substrates to neurons, an arrangement often called the \textit{astrocyte--neuron lactate shuttle}.

\paragraph{5. The blood-brain barrier.} The CNS is protected from the rest of the body by the \textbf{blood-brain barrier (BBB)}, a highly selective wall that only specific molecules can cross. The BBB is built chiefly by tight junctions between endothelial cells in brain capillaries, but astrocytic end-feet wrap around those vessels and are essential for inducing and maintaining the barrier's properties. Without astrocytes, the BBB does not form properly.

\paragraph{6. Modulation of synapses.} The ``tripartite synapse'' concept treats every synapse as involving three cells, not two: the presynaptic neuron, the postsynaptic neuron, and an astrocytic process sitting alongside them. The astrocyte can detect neurotransmitter release, respond with its own calcium signals, and release \textit{gliotransmitters} that feed back onto nearby neurons. How large a role this plays in moment-to-moment information processing is still an active research question, but the days of astrocytes as inert scaffolding are firmly over.

\paragraph{7. Response to injury.} After any significant CNS insult, astrocytes become \textit{reactive}. They enlarge, proliferate, change which genes they express, and in severe cases help form a dense \textbf{glial scar}. This response is double-edged: it walls off damage and limits inflammation, but the scar itself is a significant obstacle to the regeneration of neurons, which is one of the reasons CNS injuries heal so poorly.

\subsubsection{Satellite Cells}

Satellite cells are the PNS cousins of astrocytes. They form a close-fitting envelope around the cell bodies of sensory and autonomic neurons inside peripheral \textbf{ganglia}, such as the dorsal root ganglia. Their functions overlap substantially with those of astrocytes, structural support, control of the local chemical environment, regulation of ion concentrations, modulation of excitability, but they act on a much smaller, more localised scale.

Interest in satellite cells has grown sharply in recent years because of their role in chronic pain. Changes in satellite cell signalling are thought to contribute to the abnormal excitability of sensory neurons in persistent pain states, which makes them one of the more surprising targets on the emerging map of pain biology.

\subsection{Oligodendrocytes and Schwann Cells}

\subsubsection{Oligodendrocytes}

Oligodendrocytes are the myelinating glia of the CNS. Their name means ``cells with few branches,'' a deliberate contrast with the exuberant astrocyte, and their central job is to wrap axons in \textbf{myelin}, the lipid-rich insulation that makes fast signalling possible.

A single oligodendrocyte can myelinate segments of several different axons at once, reaching out a process to each and wrapping it repeatedly. This is one of the most important differences between CNS and PNS myelin: a Schwann cell commits itself to a single internode on a single axon, whereas one oligodendrocyte serves many. The arrangement is efficient, but it also means that the loss of a single oligodendrocyte strips myelin from several axons at once.

Their jobs are closely linked. They \textbf{produce myelin}, layer upon layer of plasma membrane wrapped tightly around the axon. They also \textbf{support axonal health} in ways that go beyond simple insulation, damage to oligodendrocytes tends to cause damage to the axons they serve, even when the axons themselves are initially untouched. And they \textbf{organise the axonal membrane} into alternating \textbf{internodes} (myelinated) and \textbf{nodes of Ranvier} (bare gaps where voltage-gated sodium channels cluster).

The clustering of sodium channels at the nodes is what makes \textbf{saltatory conduction} possible. Instead of the action potential regenerating continuously at every point along the axon, it regenerates only at the nodes and effectively jumps from one to the next. The result is faster, more energy efficient, and more reliable.

\subsubsection{Schwann Cells}

Schwann cells are the PNS counterparts of oligodendrocytes, and they come in two forms.

A myelinating Schwann cell wraps itself repeatedly around a single segment of a single axon, producing one \textbf{internode} of myelin. One Schwann cell per internode, one axon at a time. Multiple Schwann cells are therefore needed along the length of any given PNS axon, separated by nodes of Ranvier just as in the CNS.

Not every peripheral axon is myelinated. Small, slowly conducting axons (including many pain fibres and autonomic fibres) are instead embedded in non-myelinating Schwann cells, often multiple axons per cell, in structures called \textbf{Remak bundles}. These Schwann cells do not build concentric wrappings, but they still provide mechanical support, metabolic support, and protection.

The single most striking difference between the PNS and the CNS is that peripheral nerves can regenerate and central nerves largely cannot. Schwann cells are a large part of why. After injury to a peripheral nerve, Schwann cells clear away the debris, proliferate, and form longitudinal tracks that actively guide regrowing axons back toward their targets, while secreting growth factors that promote regrowth. A cut peripheral nerve has a real chance of reconnecting; a cut spinal cord generally does not.

\subsubsection{Myelin}

Myelin is not a cell type in its own right, it is built by oligodendrocytes and Schwann cells, but its electrical role is specific enough to be worth setting out separately.

Physically, myelin is a multilayered sheath of plasma membrane wrapped tightly around an axon, squeezing out almost all of its cytoplasm so that what remains is essentially a stack of lipid bilayers. The stack is heavily lipid-rich, which is why myelinated tissue looks pale and is called \textit{white matter}, and it contains a characteristic set of structural proteins.

There are two electrical tricks at work. First, wrapping an axon in many layers of membrane enormously \textbf{increases its resistance} to transverse current flow, which means less current leaks out across the membrane. Second, and equally importantly, it \textbf{decreases its capacitance}: a stack of membranes in series stores less charge than a single membrane, so less charge needs to be moved to change the voltage. Both effects mean that a given current spreads further and faster along the inside of the axon.

Instead of regenerating at every point, the action potential now jumps almost passively along the internodes, pausing at the nodes of Ranvier only long enough for sodium channels there to boost it before the next leap. The net result is dramatic: a well-myelinated axon can conduct at roughly 100 m/s, versus perhaps 1 m/s for an unmyelinated one of the same diameter.

When myelin is stripped away, the electrical tricks disappear. Current leaks across the bare membrane, conduction slows, and action potentials may simply fail to propagate. This can cause severe neurological deficits even when the axons underneath are, for the time being, structurally intact, which is precisely what is so frightening about demyelinating diseases. Fix the myelin and the function can return; let the damage accumulate and the underlying axons may eventually die too.

\subsection{Ependymal Cells}

Ependymal cells form an epithelial-like lining along the walls of the ventricles and the central canal of the spinal cord. They are glial cells of a rather different flavour from astrocytes or oligodendrocytes: their job, essentially, is plumbing.

A specialised subset of ependymal cells, together with associated capillaries and connective tissue, forms the \textbf{choroid plexus}, the structure inside the ventricles that actively secretes cerebrospinal fluid. CSF is the fluid we met earlier: it cushions the brain, keeps its chemical environment stable, and removes waste products. Without ependymal cells, there is no CSF.

Beyond production, ependymal cells also help move CSF. Many of them have \textbf{cilia} on their apical surface, and the beating of these cilia drives circulation through the ventricular system. The ependymal lining also contributes to the regulation of exchange between CSF and brain tissue, though it is not as tight a barrier as the blood-brain barrier proper.

Clinically, anything that obstructs CSF flow or overwhelms its reabsorption can produce \textbf{hydrocephalus}, in which CSF accumulates inside the ventricles. The pressure rises, the ventricles balloon, and, especially in infants, whose skulls are still flexible, the head itself can visibly enlarge. Treatment usually involves surgically diverting CSF with a shunt.

\subsection{Microglia}

Microglia are the immune system of the CNS. They are unusual among glial cells in that they are not of ectodermal origin: they descend from the same mesodermal lineage as macrophages and other immune cells, and they migrate into the brain during early development. Once installed, they are essentially permanent residents.

In their ``resting'' state, microglia are not really resting at all. Their processes extend, retract, and extend again, continuously surveying the local tissue for anything out of place, pathogens, debris, abnormally folded proteins, damaged cells, inappropriate synaptic structures. When they find something, they respond.

\subsubsection{What Microglia Do}

\paragraph{Immune surveillance.} This is the microglial baseline, and it is continuous rather than occasional. The processes described above sweep the surrounding tissue without pause, so that the whole volume of the brain falls under inspection on a timescale of hours despite the cells themselves barely moving. What they are testing for is anything that should not be there: a pathogen, the debris of a dead cell, a protein that has misfolded, a synapse that is no longer appropriate. Surveillance is therefore not a passive readiness to respond but the activity that occupies most of a microglial cell's life, and everything else in this list is what happens on the occasions when it finds something.

\paragraph{Phagocytosis.} When activated, microglia engulf and digest dead cells, cellular debris, pathogens, and protein aggregates. They are the only professional phagocytes the CNS has.

\paragraph{Inflammatory signalling.} Microglia release cytokines and chemokines to coordinate local inflammatory responses. In acute infection or injury this is protective. In chronic or excessive activation it is corrosive, and chronic microglial inflammation is increasingly implicated in neurodegenerative disease.

\paragraph{Synaptic pruning.} During development, microglia physically remove unnecessary synapses, helping to refine neural circuits into their mature form. Remarkably, this pruning role continues into adulthood, and disruptions to microglial pruning are now suspected of contributing to several disorders of neurodevelopment.

\subsubsection{Activation States}

Microglial activation is not a simple switch. A resting microglial cell looks highly branched (``ramified''); an activated one retracts its branches and takes on a more amoeboid shape appropriate to its phagocytic role. But the states in between span a wide spectrum, and depending on context the same cell can be protective or destructive. The current picture is more one of \textit{roles} than of discrete states.

\subsection{Comparing CNS and PNS Glia}

The two compartments solve the same problems with different cells.

\paragraph{In the CNS:} astrocytes maintain homeostasis and support the BBB; oligodendrocytes myelinate; ependymal cells line the ventricles and make CSF; microglia provide immune defence.

\paragraph{In the PNS:} Schwann cells myelinate (and non-myelinating Schwann cells envelop small axons); satellite cells support the cell bodies of ganglion neurons.

The parallels are not perfect, but the functional principles are the same, control the environment, insulate the axons, defend the tissue, and each system has specialised cells for each role.

\subsection{Summary}

Glia are not the supporting cast of the nervous system. Astrocytes and satellite cells regulate the environment in which neurons live, holding ion concentrations, transmitter levels, and extracellular composition within the narrow bounds that electrical signalling requires, and supplying the energy substrates that neurons cannot store for themselves. Oligodendrocytes and Schwann cells produce the myelin without which fast conduction is impossible. Ependymal cells build and circulate the CSF. Microglia provide the standing immune defence of a compartment that the rest of the body's immune system cannot easily reach, and, with astrocytes, actively shape synapses rather than merely surrounding them. Glia also carry out most of the maintenance and repair of nervous tissue, including the regenerative capacity that distinguishes peripheral nerves from central ones. Neuronal activity is the visible half of the story, but none of it would be possible without this infrastructure, and any serious account of how the brain works, or of how it fails in disease, has to take the glial half seriously.

One glial contribution matters more than the others for what follows, and it is myelin. The figures given above, roughly a hundred metres per second along a myelinated axon against perhaps one along a bare one of the same width, are the difference between a nervous system that can act on the world and one that can only comment on it afterwards. A signal from the foot reaches the brain in something like twenty milliseconds, and a response returns on a similar schedule, which is why balance is possible and why a stumble can be caught. Without that speed the same computations would still be correct and would arrive too late to be of use. Much of what the later sections describe, the binding of vision to touch into a single scene, the interception of a movement before it is completed, the fast appraisal of a threat, presupposes that signals cross the brain in milliseconds rather than seconds. Cognition happens in real time because glia make the wiring fast enough, and that is an achievement of insulation rather than of computation.

\section{The Frontal Lobe}

We now have a working model of individual neurons, the signals they generate, and the glial tissue that supports them. The remaining sections move up to the systems level, beginning with the region that has, more than any other, been identified with personhood itself.

The \textbf{frontal lobe} is the largest of the cerebral lobes and sits just behind the forehead. It is responsible for voluntary movement, the planning and execution of actions, language production, and a cluster of higher cognitive and behavioural functions we collectively call \textit{executive function}, judgment, decision-making, inhibition, personality. If the occipital lobe is where vision lives and the temporal lobe is where memory and hearing live, the frontal lobe is where thought is translated into action, and where action is held back in the service of longer-term goals.

It is also the region most implicated in the overlap between neurology and psychology, which makes it the natural place to begin the survey of the cortex. Damage to different parts of the frontal lobe produces characteristic and sometimes startling changes in behaviour and personality, changes that force us to confront how thoroughly the ``mind'' is rooted in very specific pieces of tissue.

\subsection{Anatomical Landmarks}

The frontal lobe is bounded by three anatomical landmarks:

\begin{itemize}
    \item The \textbf{central sulcus} (Rolandic fissure) runs roughly down the top of the hemisphere and separates the frontal lobe from the parietal lobe. It is the most clinically important sulcus in the brain, because it divides the \textit{motor} cortex in front from the \textit{somatosensory} cortex behind.
    \item The \textbf{lateral sulcus} (Sylvian fissure) runs along the side of the hemisphere and separates the frontal lobe from the temporal lobe below.
    \item The \textbf{longitudinal fissure} runs down the midline and separates the two hemispheres from each other.
\end{itemize}

Within the frontal lobe, a few gyri (ridges) are worth knowing by name. The \textbf{precentral gyrus}, running immediately in front of the central sulcus, contains the primary motor cortex. The \textbf{superior, middle, and inferior frontal gyri}, running along the lateral surface, house progressively more abstract motor and cognitive areas toward the front. Broca's area, for example, sits in part of the inferior frontal gyrus.

The rest of this section walks forward from the most concrete function, the movement of muscles, to the most abstract, the governance of conduct and character.

\subsection{Primary Motor Cortex (M1)}

The \textbf{primary motor cortex}, traditionally labelled \textbf{M1} and corresponding to Brodmann area 4, occupies the precentral gyrus just in front of the central sulcus. It is the final cortical stage of voluntary movement: the place where the abstract intention to move becomes the concrete commands that drive muscles.

Its output leaves the cortex mainly via the \textbf{corticospinal tract} (for the body) and the \textbf{corticobulbar tract} (for the face and head), eventually reaching lower motor neurons in the brainstem and spinal cord that directly innervate muscles. Crucially, the corticospinal tract crosses to the opposite side of the body as it descends, so the left M1 controls the right side of the body and vice versa. This \textbf{contralateral} organisation is why a stroke in one hemisphere weakens the opposite side.

\subsubsection{The Motor Homunculus}

M1 is not a featureless sheet. It is organised \textit{somatotopically}: different patches of cortex control different parts of the body, arranged in a roughly orderly map along the gyrus. This map is often drawn as a small distorted human figure, the \textbf{motor homunculus}, draped along the cortex.

Walking from medial to lateral along M1:

\begin{itemize}
    \item The \textbf{medial} regions (down in the longitudinal fissure) control the lower limb.
    \item The \textbf{superior lateral} regions control the trunk and upper limb.
    \item The \textbf{lateral} regions control the face, lips, and tongue.
\end{itemize}

What makes the homunculus so striking visually is that it is wildly disproportionate. The hands and face occupy enormous stretches of cortex, while the trunk and legs get relatively little. The representation is proportional not to body size but to the fineness of motor control that part of the body demands. Hands and lips perform feats of precision on a millisecond timescale; the muscles of the lower back do not.

\subsubsection{Microanatomy}

M1 contains exceptionally large pyramidal neurons called \textbf{Betz cells}, which project long axons into the corticospinal tract described above. These are among the biggest neurons in the entire CNS, and the size is not incidental: a cell body must support an axon reaching from the cortex to the lower spinal cord, over a metre in an adult, and sustain the transport traffic that such a projection demands.

Their size is best understood against the transport problem described earlier: a cell sustaining an axon of that length must manufacture everything its distant terminals consume and keep the kinesin and dynein traffic running along the whole of it, and the soma is scaled to that burden. Betz cells are nonetheless only part of the story, accounting for a small fraction of corticospinal fibres; most of the tract arises from smaller pyramidal neurons alongside them.

\subsection{Motor Association Cortex: Premotor and Supplementary Motor Areas}

Immediately anterior to M1 lies the \textbf{motor association cortex} (Brodmann area 6), which is itself divided into two main regions: the \textbf{premotor cortex} on the lateral surface and the \textbf{supplementary motor area (SMA)} on the medial surface. These are the areas where movements are \textit{planned} rather than executed. A useful way to think about them is that M1 decides how to move, and the motor association cortex decides what to move and in what order.

Their jobs include sequencing complex movements, coordinating the body so that posture and proximal muscles are prepared before the main action begins, and integrating sensory input into motor plans. There is a division of labour between them that is clinically worth remembering: the \textbf{supplementary motor area} is particularly important for \textit{internally generated} and \textit{bimanually coordinated} movements, movements arising from internal intention rather than external prompting, and those requiring both hands to cooperate, while the \textbf{premotor cortex} is more involved in movements \textit{cued by external stimuli}, such as reaching for an object that has just come into view.

Damage to these areas typically does not produce paralysis. It produces \textbf{apraxia}: a striking dissociation in which strength and coordination are intact, but the patient cannot perform learned purposeful movements on request. Asked to wave goodbye or mime brushing their teeth, they fumble. They still have the ``how,'' in the low-level sense, but they have lost the ``what'', the motor plan. It is one of the cleanest demonstrations in neurology that planning and executing a movement live in different places.

\subsection{Frontal Eye Field}

Tucked into the posterior part of the middle frontal gyrus, and corresponding to Brodmann area 8, is the \textbf{frontal eye field (FEF)}. Its job is voluntary eye movement, in particular the rapid, ballistic eye movements called \textbf{saccades} that shift gaze from one point to another.

The FEF is the cortical route by which attention is directed through the eyes: when a decision is made to look at something, the FEF is a major node in the network that carries it out. It connects extensively to brainstem gaze centres and to the cranial nerve nuclei that actually drive the eye muscles. Stimulating one frontal eye field causes the eyes to move toward the opposite side; damaging it produces the opposite pattern, with the eyes drifting toward the side of the lesion because the healthy FEF on the other side is now unopposed. This ``eyes toward the lesion'' picture is a classic sign in the early hours of a frontal stroke.

\subsection{Broca's Area}

In the \textbf{inferior frontal gyrus} of the dominant hemisphere, which is almost always the left, even in most left-handers, lies \textbf{Broca's area}, corresponding to Brodmann areas 44 and 45. It is the part of the brain most closely associated with the \textit{production} of language. It plans the motor sequences of speech, handles much of the grammar and sentence structure, and coordinates the muscles of the tongue, lips, larynx, and diaphragm that actually do the talking.

Damage to Broca's area produces \textbf{Broca's aphasia} (also called expressive or non-fluent aphasia). The clinical picture is characteristic and often heartbreaking: speech is slow, effortful, and telegraphic, content words survive, grammatical glue words drop out, and patients know, agonisingly, that what they are saying is not what they mean. Comprehension is relatively preserved, which is what distinguishes Broca's aphasia from the fluent but empty speech of Wernicke's aphasia, whose lesion lies further back in the temporal lobe.

Broca's area is also a reminder that cortical function is \textbf{lateralised}. The two hemispheres are not functional mirror images of one another. Language production sits overwhelmingly on the left for most people, while analogous regions on the right handle things like prosody and emotional tone.

\subsection{Prefrontal Cortex}

Moving still further forward, we arrive at the \textbf{prefrontal cortex (PFC)}: the portion of the frontal lobe that lies in front of the motor and premotor areas. This is the most recently evolved part of the frontal lobe and, in humans, proportionally the largest. It is also where the neurology of the frontal lobe bleeds most clearly into psychology.

The PFC does not produce movement directly. It produces the conditions under which movement makes sense. It holds goals in mind over time, compares possible courses of action, suppresses impulses that would be counterproductive, and integrates emotion, memory, and sensation into coherent behaviour. The whole constellation of abilities we call \textbf{executive function} lives here: working memory, planning, flexible switching between tasks, inhibition of prepotent responses, and the capacity to act in accordance with long-term goals rather than short-term reward.

\subsubsection{Subdivisions}

The PFC divides into three functional regions, each with its own characteristic contribution.

\paragraph{Dorsolateral prefrontal cortex (DLPFC).} This is the region most associated with classical ``cold'' executive function: working memory, problem-solving, planning, and flexible control of attention. It is the part of the brain that holds a phone number in mind while it is being dialled, keeps a multi-step instruction in order, or juggles the possibilities in a chess position. DLPFC dysfunction is implicated in attention deficits and is heavily involved in schizophrenia.

\paragraph{Ventromedial prefrontal cortex (VMPFC).} This region is more concerned with ``hot'' cognition: emotion, risk, and reward. It integrates signals from the limbic system into decision-making, which is why damage here can leave formal reasoning intact while devastating the patient's ability to make real-life choices. Patients with VMPFC lesions can pass IQ tests and still ruin their lives through impulsive, poorly weighed decisions.

\paragraph{Orbitofrontal cortex (OFC).} Sitting just above the orbits of the eyes, the OFC is closely involved in social behaviour, impulse control, and the moment-to-moment evaluation of rewards and punishments. It is especially important for recognising when something that used to be rewarding has stopped being so, and for holding back socially inappropriate behaviour.

\subsubsection{Connections}

The PFC is one of the most richly connected regions of the brain. It has strong reciprocal links with the limbic system (amygdala, hippocampus), with sensory and motor cortices, with the thalamus, and with the basal ganglia. This connectivity is what lets it perform its integrative role: executive function is, in a very literal sense, the business of pulling together information from everywhere else.

\subsection{Summary}

The frontal lobe is where the nervous system turns thought into action, and where action is disciplined in the service of goals. Its organisation reflects a rough gradient. The primary motor cortex at the back of the lobe executes movement; the premotor and supplementary motor areas plan and sequence it; the frontal eye field and Broca's area carry out specialised motor behaviours of gaze and speech; and the prefrontal cortex, sitting in front of all of these, supplies the executive control, judgment, and personality that make any of those actions meaningful. Damage at each level produces a characteristic and clinically revealing deficit, and together these deficits provide some of the most compelling evidence we have that the mind is, in the end, a function of specific pieces of brain tissue, which is both the starting point of neurology and the doorway through which psychology must pass.

This is also the first section in which the account reaches something a person would recognise as themselves. Judgment, restraint, the weighing of a long-term goal against an immediate temptation, the capacity to be reliably one sort of person rather than another: these are not capacities of neurons or of cortex in general but of particular regions, and they are lost in particular ways when those regions are damaged. A patient with a ventromedial lesion who reasons perfectly and chooses disastrously is difficult to fit into any account in which character sits above and apart from tissue. Everything up to this point has been mechanism with no obvious bearing on personhood. From here on the two are hard to keep separate.

\section{The Parietal Lobe}

Having mapped the frontal lobe, the cortex of action and intention, we now move backwards across the central sulcus into territory of a different character. The \textbf{parietal lobe} is the cortex of \textit{sensation} and \textit{space}: where the body's tactile world is registered, where the position of every limb is tracked moment by moment, and where the brain assembles its sense of where things are, including where the body itself sits within the world.

Where the frontal lobe pushes outward, sending commands to muscles, the parietal lobe pulls inward, taking the river of information arriving from the skin, joints, muscles, and viscera and turning it into a coherent picture of the body and its surroundings. It is also the place where touch meets vision, where felt position meets seen position, and where the multiple sensory streams are bound into the unified spatial scene that we treat, naively, as simply ``the world.''

\subsection{Anatomical Landmarks}

The parietal lobe sits behind the frontal lobe and above the temporal lobe, occupying roughly the upper-rear quadrant of each hemisphere. Its boundaries are:

\begin{itemize}
    \item Anteriorly, the \textbf{central sulcus} separates it from the frontal lobe. The first ridge behind the sulcus, the \textbf{postcentral gyrus}, mirrors the precentral gyrus in front and contains the primary somatosensory cortex.
    \item Posteriorly, the \textbf{parieto-occipital sulcus}, visible most clearly on the medial surface, marks the boundary with the occipital lobe.
    \item Inferiorly, the \textbf{lateral sulcus} (Sylvian fissure) separates it from the temporal lobe.
    \item Medially, the \textbf{longitudinal fissure} divides it from its counterpart on the other side.
\end{itemize}

Within the parietal lobe, the most important subdivision is made by the \textbf{intraparietal sulcus}, which runs roughly horizontally and divides the surface into a \textbf{superior parietal lobule} above and an \textbf{inferior parietal lobule} below. The inferior parietal lobule is itself made up of two named gyri: the \textbf{supramarginal gyrus}, which curls around the end of the lateral sulcus, and the \textbf{angular gyrus}, which sits behind it and curls around the end of the superior temporal sulcus. These are landmarks worth keeping in mind, as they mark the functional subdivisions the rest of this section works through.

Functionally, the parietal lobe walks from the most concrete to the most abstract from front to back. The postcentral gyrus, just behind the central sulcus, registers raw sensation; the cortex behind it interprets and integrates that sensation; and the cortex further back still uses sensation to construct space, attention, and bodily self-awareness.

\subsection{Primary Somatosensory Cortex (S1)}

The \textbf{primary somatosensory cortex}, traditionally labelled \textbf{S1} and corresponding to \textbf{Brodmann areas 3, 1, and 2}, occupies the postcentral gyrus immediately behind the central sulcus. It is the cortical destination of the body's sensory traffic, the place where information about touch, pressure, vibration, joint position, and pain finally becomes conscious sensation.

The information arriving at S1 has come a long way. \textbf{Discriminative touch, vibration, and proprioception} ascend through the \textbf{dorsal column--medial lemniscus pathway}; \textbf{pain and temperature} ascend through the \textbf{spinothalamic tract}. Both pathways relay through the \textbf{thalamus}, specifically the \textbf{ventral posterolateral (VPL) nucleus} for the body and the \textbf{ventral posteromedial (VPM) nucleus} for the face, before fanning out to S1. As with the motor system, the projection is \textbf{contralateral}: the right S1 represents the left side of the body and vice versa.

\subsubsection{The Sensory Homunculus}

S1 is somatotopically organised in just the same way M1 is, and the resulting map is called the \textbf{sensory homunculus}. Walking from medial to lateral along the postcentral gyrus:

\begin{itemize}
    \item The \textbf{medial} regions, tucked into the longitudinal fissure, represent the lower limb and genitals.
    \item The \textbf{superior lateral} regions represent the trunk and upper limb.
    \item The \textbf{lateral} regions represent the face, lips, and tongue.
\end{itemize}

Like its motor counterpart, the sensory homunculus is grotesquely disproportionate. The lips, tongue, fingertips, and face occupy enormous swathes of cortex; the back, trunk, and legs make do with much less. The principle is the same as in M1: cortical real estate scales with the \textit{density of receptors} and the \textit{fineness of discrimination} demanded of that body part, not with its physical size. The fingertips can distinguish points a millimetre or two apart, and they require the cortex to do it.

\subsubsection{Microanatomy}

The three Brodmann subdivisions of S1 are not redundant. \textbf{Area 3} is the most anterior strip, sitting in the wall of the central sulcus, and it is the principal recipient of thalamic input. It is itself split into \textbf{3a}, which receives proprioceptive information from muscles and joints, and \textbf{3b}, which receives cutaneous information from the skin. Behind these, \textbf{area 1} processes mainly cutaneous input as well, and \textbf{area 2} integrates cutaneous and proprioceptive signals to support more complex perceptions such as size, shape, and texture. There is, in other words, an early stage of integration happening already within S1 itself.

\subsection{Somatosensory Association Cortex}

Behind S1, on the \textbf{superior parietal lobule}, lies the \textbf{somatosensory association cortex}, corresponding to \textbf{Brodmann areas 5 and 7}. Where S1 registers raw sensation, the association cortex \textit{interprets} it. It takes the disjointed feel of edges, textures, weights, and contours arriving from S1 and assembles them into the recognition of objects, the sense of one's body shape, and the awareness of where each limb is in space.

A useful way to put it is that S1 registers \textit{that} something is touching the body and \textit{where}, while the association cortex determines \textit{what it is}. Identifying a key in a pocket without looking, distinguishing a coin from a button by feel, judging the heft of a tool well enough to use it, all of this is somatosensory association cortex at work. It also stores tactile memories, so that the feel of familiar objects becomes immediately recognisable rather than having to be reconstructed from scratch each time.

Damage to this region produces a striking deficit known as \textbf{tactile agnosia}, or, when restricted to a hand, \textbf{astereognosis}. Strength is intact, primary sensation is intact, and the patient can describe the temperature, texture, and shape of an object placed in the hand, but cannot identify what the object is. The raw materials of perception are present; the act of recognition has been severed from them. As with the motor association cortex's apraxia, this is a clean dissociation between low-level capability and the higher-level integration that makes it useful.

\subsection{The Posterior Association Area}

Further back still, we reach the point at which cortex stops belonging to any one sense. Everything described so far in this section has been somatosensory: S1 registers touch, and the somatosensory association cortex interprets it, but both are working with a single modality. The \textbf{posterior association area} is different in kind. It is the large expanse of cortex where the separate sensory worlds are combined into one, and it is the reason we experience a single scene rather than three parallel reports.

Anatomically it does not respect lobar boundaries, which is precisely the point. It occupies the \textbf{parieto-temporo-occipital junction}, the territory where the three posterior lobes meet, and it takes in the \textbf{posterior parietal cortex} (the superior and inferior parietal lobules, including the \textbf{angular} and \textbf{supramarginal gyri} identified earlier) along with the adjacent posterior temporal and lateral occipital cortex. It is discussed here because its centre of mass is parietal, but a boundary drawn around it would cut across the parietal, temporal, and occipital lobes alike.

Its position explains its function. Three great streams of processed information converge on it: \textbf{somatosensory} information from the parietal cortex in front, \textbf{visual} information from the occipital cortex behind, and \textbf{auditory} information from the temporal cortex below, with vestibular input arriving alongside them. Each stream arrives already interpreted by its own association cortex. What the posterior association area adds is \textit{binding}: the felt position of the hand and the seen position of the hand are registered as one hand, and the sound of a voice and the sight of a moving mouth are registered as one person speaking.

The most conspicuous product of that binding is a representation of \textbf{space}. The posterior association area constructs where things are, where the body is among them, and how the two relate, holding this in several frames at once, some anchored to the body and some to the world. This is what allows a hand to reach accurately for a cup without being watched, a body to pass through a doorway without striking the frame, a moving object to be tracked while the head turns, and a room to be experienced as stable even though the retinal image lurches with every saccade.

Bound up with space is \textbf{attention}. The area is a major node in the brain's attention network, working with the frontal eye field to decide which part of the environment matters next. It is also the destination of the visual \textbf{dorsal stream} described in the occipital lobe section, the ``where'' and ``how'' pathway, and it is here that seeing is converted into the coordinates a movement can actually use, so that a reach or a grasp arrives where the object is rather than where it merely appeared.

Two further contributions are worth naming. The area maintains the \textbf{body schema}, the continuously updated model of the body's shape, posture, and extent that makes the body feel like a single owned object rather than a collection of limbs. And through the \textbf{angular gyrus}, which sits at the junction proper and adjoins the language cortex met in the temporal lobe, it forms the bridge from perception to \textbf{symbol}: the step at which a seen mark becomes a letter, a heard pattern becomes a word, and a quantity becomes a number.

Like the frontal lobe, the posterior association area is strongly \textbf{lateralised}. The right hemisphere carries the heavier load for spatial and attentional processing, tending to attend to both sides of space, while the left is more involved in the language-related and symbolic operations of the angular gyrus, in calculation, and in the ordered representation of learned actions. This asymmetry is functional rather than absolute, and the two sides work together through the corpus callosum, but it is a real division of labour and it explains why the region's contributions cannot be understood as a single uniform capacity.

\subsection{Summary}

The parietal lobe is where the body's sensory traffic becomes conscious sensation, where sensation becomes recognition, and where recognition is woven together with vision and movement into the felt sense of a body in space. The primary somatosensory cortex registers touch, proprioception, pain, and temperature in a contralateral, somatotopically organised map; the somatosensory association cortex turns those raw sensations into recognised objects and a coherent body schema; and the posterior association area, straddling the junction of the parietal, temporal, and occipital lobes, binds the separate sensory streams into a single scene and constructs the space, attention, and visuomotor frames within which action becomes possible. Damage at each level produces a characteristic deficit, from contralateral numbness at the first, through the tactile agnosia of the second, to disturbances of space, attention, and the body schema at the third. Taken together, the parietal lobe shows that even something as immediate as ``what the world feels like'' is the product of layered cortical processing, and that the mind's sense of being embodied in a place is not given but built.

This is the first full instance of the claim flagged at the outset, and it is worth pausing on, because the body schema is about as unpromising a candidate for a construction as could be chosen. Occupying a body does not feel like an achievement. It feels like the background condition against which everything else is experienced, the one thing that could not require assembly. Yet it can be dismantled, and the ways it comes apart are informative: a limb can be felt in the wrong place, or disowned, or duplicated, or experienced as belonging to someone else, and each of these follows damage to a particular part of the machinery described above. What cannot be done is to damage the parietal lobe and leave the sense of embodiment untouched but incorrect in some way the anatomy does not predict. The construction has parts, and the failures respect them.

\section{The Temporal Lobe}

Dropping below the lateral sulcus, we leave the cortex of touch and space and enter the cortex of \textit{sound}, \textit{smell}, and \textit{meaning}. The \textbf{temporal lobe} is where vibrations of air become recognised speech and music, where the volatile chemistry of the world becomes the pull of a familiar smell, and where, in its language regions, the heard word becomes the understood word.

If the frontal lobe turns intention into action and the parietal lobe turns sensation into space, the temporal lobe turns sensation into \textit{recognition} and \textit{meaning}. It is where the auditory and olfactory streams arrive in cortex, where they are interpreted, and where the interpretation of language as comprehensible content first becomes possible. It is also, along the way, the cortex most tightly bound to the limbic system, which is part of why sound and smell are so good at evoking memory and emotion.

\subsection{Anatomical Landmarks}

The temporal lobe sits below the lateral sulcus and forwards of the occipital lobe, occupying roughly the lower-side of each hemisphere. Its principal external landmarks are three roughly parallel gyri running along its lateral surface:

\begin{itemize}
    \item The \textbf{superior temporal gyrus}, immediately below the lateral sulcus, contains the auditory cortices and the receptive language area in the dominant hemisphere.
    \item The \textbf{middle temporal gyrus}, beneath it, contributes to higher-order auditory and visual processing, including aspects of object and motion recognition.
    \item The \textbf{inferior temporal gyrus}, lower still, contributes to high-level visual recognition, particularly of faces and objects.
\end{itemize}

Two further landmarks matter. Tucked into the floor of the lateral sulcus on the upper surface of the superior temporal gyrus lie \textbf{Heschl's gyri}, also called the \textbf{transverse temporal gyri}, which house the primary auditory cortex itself. On the medial surface, the temporal lobe curls inward to form the \textbf{uncus} and the \textbf{parahippocampal gyrus}, structures continuous with the limbic system; the primary and association olfactory cortices live here, alongside the hippocampus, the structure most closely associated with the formation of declarative memory.

As with the other lobes, the temporal lobe walks from concrete to abstract: primary sensory cortex first, association cortex behind it, and language and memory regions further back and inward.

\subsection{Primary Auditory Cortex (A1)}

The \textbf{primary auditory cortex}, traditionally labelled \textbf{A1} and corresponding to \textbf{Brodmann areas 41 and 42}, occupies Heschl's gyri on the upper surface of the superior temporal gyrus, mostly hidden inside the lateral sulcus. It is the cortical destination of the auditory pathway, where pressure waves in the ear finally become the conscious experience of sound.

The pathway that delivers sound to A1 is one of the longest and most heavily processed in the brain. Hair cells in the \textbf{cochlea} convert mechanical vibrations into neural signals carried by the \textbf{cochlear division of the vestibulocochlear nerve (CN VIII)}. From there, signals pass through the \textbf{cochlear nuclei} of the brainstem, the \textbf{superior olivary complex}, the \textbf{lateral lemniscus}, the \textbf{inferior colliculus} of the midbrain, and the \textbf{medial geniculate nucleus (MGN)} of the thalamus before finally fanning out to A1. At several points along this chain, information crosses the midline, with the result that each A1 receives input from \textit{both} ears. This bilateral organisation is why a unilateral lesion of A1 does not produce deafness in either ear; it produces only a subtle deficit in tasks like sound localisation and the perception of complex auditory scenes.

A1 is not a featureless sheet either. It is organised \textbf{tonotopically}: different patches of cortex respond to different frequencies, arranged in an orderly map that mirrors the mechanical tuning of the cochlea. Roughly, \textbf{high frequencies} are represented at one end of A1 and \textbf{low frequencies} at the other, with a smooth gradient between. This is the auditory analogue of the somatotopic organisation of S1 and the retinotopic organisation of the visual cortex: a peripheral sensory surface mapped, in orderly fashion, onto the cortex.

\subsection{Auditory Association Cortex}

Surrounding A1 on the superior temporal gyrus, and corresponding to most of \textbf{Brodmann area 22}, lies the \textbf{auditory association cortex}. Where A1 registers raw sound, the association cortex \textit{interprets} it: it recognises the sound as speech, music, a familiar voice, a doorbell, or a piece of meaningless noise. It is also the region that gives auditory experience much of its temporal structure, the perception of rhythm, melody, and the patterning of language at the level of pure sound.

Damage to the auditory association cortex produces \textbf{auditory agnosia}, a deficit analogous to the tactile agnosia of the parietal lobe. The patient hears: pure-tone audiometry is intact, A1 is doing its job, sound is consciously perceived, but recognition has been severed from perception. Specific subtypes are clinically recognised. \textbf{Pure word deafness} leaves the patient unable to understand spoken language while still able to hear, read, write, and speak normally; non-speech sounds are recognised, but speech is reduced to noise. \textbf{Amusia} leaves the patient unable to recognise melodies that were once familiar, or, in some cases, unable to perceive music as music at all. As with other association-cortex syndromes, what has been lost is not the input but the act of making the input meaningful.

\subsection{Primary Olfactory Cortex}

Smell is the odd one out among the senses, and the temporal lobe is where its peculiarity is on display. The \textbf{primary olfactory cortex} sits on the medial surface of the temporal lobe, primarily on the \textbf{uncus} and adjacent piriform cortex, and includes parts of the \textbf{piriform cortex}, periamygdaloid cortex, and amygdala. Its inputs come from the \textbf{olfactory bulb} via the \textbf{olfactory tract (CN I)}, the only cranial nerve that connects directly to the cerebrum without first passing through the brainstem.

What makes olfaction unusual is that, alone among the major senses, it does not relay through the thalamus before reaching cortex. Olfactory receptor neurons in the nasal epithelium project to the olfactory bulb, the bulb projects to the primary olfactory cortex directly, and only \textit{after} this initial cortical processing does smell information reach the thalamus and the orbitofrontal cortex for further integration. This shortcut is part of why smell has such an immediate, unmediated quality, and why it is so tightly bound to emotion and memory: the primary olfactory cortex is essentially continuous with the limbic system.

\subsection{Olfactory Association Cortex}

Adjacent to and continuous with the primary olfactory cortex, the \textbf{olfactory association cortex} occupies the \textbf{entorhinal cortex} (Brodmann area 28) and parts of the parahippocampal gyrus. Its job is to take the raw olfactory signal arriving at primary cortex and turn it into a recognisable, named, emotionally weighted experience: this smell is coffee, that one is petrol, this one is a kitchen not entered in twenty years.

The olfactory association cortex is the most direct anatomical demonstration of why smell is so tied to memory and feeling. Through the entorhinal cortex it has powerful, reciprocal connections with the \textbf{hippocampus} (the seat of declarative memory) and the \textbf{amygdala} (a key node in emotional processing). A familiar smell can summon a specific memory with a vividness that other senses rarely match, partly because the wiring that delivers it is, in a literal sense, two synapses away from the structures that store autobiographical memory and assign emotional valence.

Damage here produces the recognition counterpart to anosmia: smells are still consciously detected but cannot be identified or attached to memory, an olfactory analogue of the agnosias seen elsewhere. Entorhinal cortex is also one of the very first regions to degenerate in \textbf{Alzheimer's disease}, which is part of why a degraded sense of smell can be one of the earliest non-cognitive signs of the illness.

\subsection{Wernicke's Area}

In the \textbf{posterior part of the superior temporal gyrus} of the dominant hemisphere, almost always the left, lies \textbf{Wernicke's area}, traditionally identified with the posterior portion of \textbf{Brodmann area 22}. It is the part of the brain most closely associated with the \textit{comprehension} of language. If Broca's area, in the inferior frontal gyrus, is where the motor plans for speech are assembled, Wernicke's area is where heard and read words become understood meaning.

Anatomically, Wernicke's area sits where it does for a reason: it is right next door to the auditory association cortex, perfectly placed to interpret the patterned sounds of speech, and it is connected to Broca's area by a long-range white-matter tract called the \textbf{arcuate fasciculus}. The arcuate fasciculus is what allows the comprehended word to be turned, almost without effort, into the spoken word: it ferries information from the back of the language network to the front.

\subsection{Summary}

The temporal lobe is where the brain turns sound into meaning, smell into memory, and heard speech into understood language. The primary auditory cortex on Heschl's gyri receives the bilateral output of the auditory pathway and arranges it tonotopically; the auditory association cortex of the superior temporal gyrus turns that raw sound into recognisable speech, music, and familiar voices. The primary olfactory cortex on the uncus, alone among the senses, receives its input directly without a thalamic relay, and the olfactory association cortex of the entorhinal region binds smell to memory and emotion through its links with hippocampus and amygdala. Finally, Wernicke's area in the dominant superior temporal gyrus stands at the centre of language comprehension, paired with Broca's area through the arcuate fasciculus to form the language network. Damage at each of these levels produces a characteristic deficit, from cortical deafness and anosmia through the auditory and olfactory agnosias. Together, the temporal lobe is where the mind acquires much of its capacity for meaning.

The capacity added here is recognition, and the agnosias show it to be a separate operation rather than a strong form of perception. A patient with pure word deafness has intact hearing and intact language, and speech reaches them as noise; the sound arrives and the meaning does not, and the two can be severed because they were never the same thing. Recognition is a further step performed on a signal already fully received. This is the point at which the account acquires the difference between a system that registers what is present and one for which anything is \textit{about} something, and it is a step that primary cortex, however intact, does not take.

\section{The Occipital Lobe}

We come now to the back of the brain, and to the simplest of the cerebral lobes to describe in terms of what it does. The \textbf{occipital lobe} is, almost in its entirety, the cortex of \textit{vision}. Almost everything the brain does with the light that enters the eye, beyond the very earliest processing in the retina and thalamus, begins here.

If the temporal lobe gives sensation its meaning and the parietal lobe gives it its place, the occipital lobe gives the visual world its existence as a coherent scene at all. From the patterned firing of retinal ganglion cells, the occipital lobe reconstructs edges, surfaces, motion, depth, colour, and form, and hands the result on to the rest of the brain for use. Damage to it does not produce subtle deficits of recognition or naming; it produces, in its purest form, the inability to see.

\subsection{Anatomical Landmarks}

The occipital lobe occupies the rear pole of each hemisphere, behind the parietal and temporal lobes. Its boundaries are:

\begin{itemize}
    \item Anteriorly and superiorly, the \textbf{parieto-occipital sulcus}, visible most clearly on the medial surface, separates it from the parietal lobe.
    \item Anteriorly and inferiorly, there is no sharp external sulcus separating it from the temporal lobe; the boundary is conventional and follows imaginary lines drawn between landmarks such as the preoccipital notch.
    \item Medially, the \textbf{longitudinal fissure} divides it from its counterpart on the other side.
\end{itemize}

The single most important landmark within the occipital lobe is the \textbf{calcarine sulcus} (or calcarine fissure), a deep horizontal groove running across the medial surface of each hemisphere. The cortex lining the upper and lower banks of this sulcus is the \textbf{primary visual cortex}, and almost everything else in the occipital lobe is, in one way or another, organised around it.

Functionally, as in the other lobes, there is a gradient from concrete to abstract. The primary visual cortex around the calcarine sulcus does the earliest cortical work; the visual association cortex surrounding it on both medial and lateral surfaces takes that early work and assembles it into recognisable scenes, objects, and motion.

\subsection{Primary Visual Cortex (V1)}

The \textbf{primary visual cortex}, traditionally labelled \textbf{V1} (and also called the \textbf{striate cortex}), corresponds to \textbf{Brodmann area 17} and lies along the banks of the calcarine sulcus on the medial surface of the occipital lobe. It is the cortical destination of the visual pathway, the first place at which the patterned firing of retinal cells, after a long journey, becomes consciously perceived vision.

The pathway that reaches V1 is heavily processed before it arrives. \textbf{Retinal ganglion cells} send their axons through the \textbf{optic nerve (CN II)} to the \textbf{optic chiasm}, where fibres from the nasal halves of each retina cross to the opposite side. After this crossing, each side of the brain carries information about the \textit{contralateral half of the visual field} from \textit{both eyes}. The fibres continue as the \textbf{optic tract} to the \textbf{lateral geniculate nucleus (LGN)} of the thalamus, which relays them via the \textbf{optic radiations} to V1. Within the optic radiations, fibres carrying information from the \textit{upper} visual field loop forward through the temporal lobe (the so-called \textbf{Meyer's loop}) before arriving at the lower bank of the calcarine sulcus, while fibres from the \textit{lower} visual field travel a more direct course through the parietal white matter to the upper bank.

\subsubsection{Retinotopic Organisation}

V1 is, like S1, A1, and M1 before it, organised as a map of its peripheral surface. This is \textbf{retinotopic} organisation: each location on the retina projects to a specific location in V1, with neighbouring retinal points represented at neighbouring cortical points. The map has two clinically vital features. First, the \textbf{contralateral} representation is strict: the \textit{left} V1 represents the \textit{right} half of visual space, and vice versa. Second, the map is heavily \textbf{magnified at the centre}: the macula, the small central region of the retina that does most of our high-acuity vision, occupies a hugely disproportionate share of V1 at the occipital pole. The same principle is at work as in the motor and sensory homunculi: cortical real estate scales with the precision required of that part of the input surface, not its physical size.

A useful corollary is that the upper visual field maps to the lower bank of the calcarine sulcus and the lower visual field to the upper bank, an inversion which, when combined with the left/right contralateral flip, means that V1 represents the visual world upside-down and mirror-reversed.

\subsubsection{Microanatomy}

V1 has two anatomical features worth knowing by name. It is called the \textit{striate} cortex because of a prominent white horizontal band visible to the naked eye in cross-section, the \textbf{stria of Gennari}, which is a particularly heavy concentration of myelinated fibres carrying input from the LGN. V1 also receives that input in two parallel streams, originating in the \textbf{magnocellular} and \textbf{parvocellular} layers of the LGN: magnocellular input carries information about motion and contrast at the expense of fine detail and colour, while parvocellular input carries information about colour and fine detail at the expense of temporal resolution. These two streams are the seeds of the visual association cortex's later division of labour.

\subsection{Visual Association Cortex}

Surrounding V1, and occupying most of the rest of the occipital lobe, is the \textbf{visual association cortex}, corresponding to \textbf{Brodmann areas 18 and 19} (often referred to as \textbf{V2}, \textbf{V3}, and beyond). Where V1 registers the basic features of the visual world, edges, orientations, contrasts, motion direction, the association cortex \textit{interprets} that early activity and assembles it into the world we actually experience: recognisable objects, faces, scenes, motion, and depth.

A great deal of what we know about visual association cortex is captured by the idea that it splits, beyond the immediate surroundings of V1, into \textbf{two streams}, each running into a different lobe. The \textbf{ventral stream} (sometimes called the ``what'' pathway) runs forward from the occipital lobe into the inferior temporal cortex; it specialises in object and face recognition, colour, and form. The \textbf{dorsal stream} (the ``where'' or ``how'' pathway) runs upward into the posterior parietal cortex; it specialises in motion, depth, spatial relationships, and the visuomotor transformations that convert a seen object into a reach toward it. The two streams are not completely independent, and the simple ``what versus where'' summary is something of an oversimplification, but it captures a real anatomical and functional division.

Damage to the visual association cortex produces a family of \textbf{visual agnosias}: the patient sees, in the sense that V1 still works and the visual field is intact, but cannot make sense of what is seen. \textbf{Visual object agnosia} is the inability to recognise objects by sight, despite being able to recognise them by touch or sound. \textbf{Prosopagnosia}, classically associated with damage to the \textbf{fusiform face area} on the inferior temporal surface where the ventral stream ends, is the strikingly specific inability to recognise faces, sometimes including one's own in a mirror, while general visual recognition remains relatively preserved. \textbf{Akinetopsia} is the inability to perceive motion, with the visual world experienced as a series of still photographs; it follows damage to area V5 (also called MT) on the lateral occipital surface, a key node in the dorsal stream's motion processing. \textbf{Achromatopsia}, cerebral colour blindness, follows damage to colour-specialised regions on the ventral surface of the occipital lobe.

These syndromes, taken together, make the same point that the agnosias of the parietal and temporal lobes made: vision, like touch and hearing, is not a single function but a stack of functions, and damage at the right level produces a startlingly specific dissociation. A patient with prosopagnosia cannot recognise their spouse's face but can recognise their voice; a patient with akinetopsia can describe a stationary cup but cannot pour from it because they cannot see the rising liquid as moving. The fact that such dissociations exist at all is one of the strongest arguments for the modular, layered architecture of cortical processing.

\subsection{Summary}

The occipital lobe is where light becomes vision. The primary visual cortex around the calcarine sulcus receives the heavily processed output of the retina and lateral geniculate nucleus, lays it out as a contralateral, retinotopic, and macula-magnified map, and performs the earliest cortical analysis of edges, orientation, contrast, and motion. The visual association cortex surrounding it interprets that early work and splits into a ventral ``what'' stream that runs into the temporal lobe to recognise objects and faces, and a dorsal ``where'' stream that runs into the parietal lobe to track motion, space, and the visual control of action. Damage at each level produces a characteristic deficit: visual field defects from V1 lesions, cortical blindness from bilateral V1 damage, and the cleanly dissociable agnosias, prosopagnosia, akinetopsia, achromatopsia, object agnosia, from damage to specific patches of association cortex. The occipital lobe is the clearest demonstration we have that even something as immediate as ``seeing the world'' is not given by the eye but constructed, in stages, by cortex.

That is the same claim the parietal lobe made about the body, now made about the visual world, and the agnosias are what make it more than an assertion. If seeing were a single capacity, damage could only make it better or worse. Instead it comes apart along seams: a patient who sees faces as faces but cannot say whose, a patient for whom a poured drink is a series of stills rather than a moving stream, a patient who sees shape and depth in a world drained of colour. Each of these is a stage of the construction failing while the others continue, and their existence is difficult to account for on any view in which vision is delivered rather than assembled. The seams are invisible in a working brain, which is precisely why they have to be found in a broken one.

\section{The Insula}

Having walked the four classical lobes, frontal, parietal, temporal, and occipital, we now turn to a fifth, often forgotten one. The \textbf{insula}, or \textbf{insular cortex}, is the cortex of \textit{the body's interior}. It is where the brain registers the state of the viscera, the feel of taste on the tongue, the slow tides of hunger and thirst, the racing of the heart, and, more abstractly, the felt quality of emotion itself. Because it has no exposed surface, it is the lobe most often left off the standard surface tour, despite carrying functions that reach from taste and visceral sensation to emotion and self-awareness.

If the occipital lobe constructs the visual world and the temporal lobe gives it auditory and linguistic meaning, the insula constructs something stranger and more intimate: the felt sense of one's own body. Interoception, the perception of internal states, is the insula's signature function. It is the part of cortex where the workings of the gut, the lungs, and the heart are turned into the conscious feelings that we recognise, after the fact, as hunger, breathlessness, anxiety, or warmth.

\subsection{Anatomical Landmarks}

Unlike the other lobes, the insula has no surface visible from outside the brain. It sits hidden in the depths of the \textbf{lateral sulcus} (Sylvian fissure), buried beneath the surrounding cortex of the frontal, parietal, and temporal lobes. The flaps of overlying cortex that cover it are collectively called the \textbf{operculum} (Latin for ``little lid''), and they are conventionally divided into a \textbf{frontal operculum} above and in front, a \textbf{parietal operculum} above and behind, and a \textbf{temporal operculum} below. Pulling these apart in dissection or imaging reveals the insula like the meat of a clam.

Within the insula, a deep groove called the \textbf{central sulcus of the insula} runs obliquely across its surface and divides it into two regions:

\begin{itemize}
    \item The \textbf{anterior insula}, made up of three short gyri, sits forward of the sulcus and is more closely tied to emotion, taste, and subjective feeling.
    \item The \textbf{posterior insula}, made up of two long gyri, sits behind the sulcus and is more closely tied to interoceptive and somatosensory processing of the body's internal and bodily state.
\end{itemize}

This anterior/posterior split runs through almost everything we know about the insula's function. Information tends to enter through the posterior insula, where bodily signals arrive in something like a raw form, and to be progressively re-represented and integrated as it moves forward, finally emerging in the anterior insula as the conscious, emotionally weighted feeling of being in a body.

\subsection{Primary Gustatory Cortex}

The insula's most concretely sensory role is as the cortical home of \textbf{taste}. The \textbf{primary gustatory cortex} occupies the anterior insula along with the adjacent \textbf{frontal operculum}, corresponding to \textbf{Brodmann area 43} and parts of the surrounding insular cortex. It is the cortical destination of the gustatory pathway, the place where signals from the taste buds finally become the conscious experience of sweet, salty, sour, bitter, and umami.

The pathway that delivers taste to the insula is one of the most multimodal in the brainstem. Taste afferents from the anterior two-thirds of the tongue travel in the \textbf{facial nerve (CN VII)}, those from the posterior third in the \textbf{glossopharyngeal nerve (CN IX)}, and those from the epiglottis and pharynx in the \textbf{vagus nerve (CN X)}. All three converge in the \textbf{nucleus of the solitary tract} in the medulla, then ascend, with a relay in the \textbf{ventral posteromedial nucleus (VPM) of the thalamus}, to the gustatory cortex. As elsewhere, the projection is largely \textit{ipsilateral}: each side of the tongue projects mainly to the same side of the cortex, in contrast to the strict contralateral organisation of the somatosensory system.

Damage to the gustatory cortex produces \textbf{ageusia} (loss of taste) or, more commonly, blunted and distorted taste, often noticed first as a loss of pleasure in food. As with smell, much of what we casually call ``taste'' is in fact flavour, the integration of taste with smell, which is why insular damage and olfactory damage can both produce strikingly similar complaints from patients about food.

\subsection{Interoception and Visceral Sensation}

Beyond taste, the larger and more distinctive role of the insula is as the cortical home of \textbf{interoception}: the perception of the body's internal state. Where S1 in the parietal lobe maps the surface of the body, the posterior and mid-insula map its interior, the heart, lungs, gut, bladder, blood vessels, and the slow chemical milieu of hunger, thirst, temperature, and breathlessness.

The pathway that delivers interoceptive information is anatomically distinct from the classical somatosensory pathways. Visceral and small-fibre afferents, including the C fibres that carry slow, burning pain and warmth, ascend in the \textbf{spinothalamic} and related tracts, relay through specialised thalamic nuclei (notably parts of the \textbf{ventral medial nucleus, posterior division (VMpo)}), and project to the posterior insula. The result is a topographic map of the body's interior: a kind of visceral homunculus. Awareness of one's own heartbeat, of the urge to breathe, of fullness in the stomach, of needing to urinate, all of these depend on insular cortex.

A theme of modern cognitive neuroscience is that the posterior insula's relatively raw representation of the body is progressively \textbf{re-represented} as it is passed forward through the mid-insula to the anterior insula. The anterior insula is where the bodily state stops being a mere physiological reading and becomes a \textit{feeling}: the hot, restless, queasy, anxious quality that distinguishes one internal state from another. The anterior insula's right side in particular has been argued to host a unified ``sentient self'' representation, the felt sense of being alive in this body at this moment, although that strong claim remains debated.

\subsection{The Anterior Insula: Emotion, Disgust, and Awareness}

Whatever its precise role in selfhood, the anterior insula is firmly implicated in a cluster of emotional and social functions that go well beyond simple visceral sensation. It is one of the most consistently activated regions in neuroimaging studies of \textbf{emotion}, and especially of negative emotions such as \textbf{disgust}, \textbf{anxiety}, \textbf{empathy for pain}, and the \textbf{craving} that drives addiction. Damage to the anterior insula can blunt the experience of disgust to the point that patients no longer recognise disgusted facial expressions in others, a deficit that mirrors the way amygdala damage blunts the experience of fear.

The anterior insula is also a major node in the brain's \textbf{salience network}, working with the anterior cingulate cortex to detect and flag events that matter, whether for survival, social interaction, or goal-directed behaviour. This is part of why the insula keeps appearing in studies of conditions as varied as anxiety disorders, depression, addiction, eating disorders, and chronic pain: in each, the basic machinery for noticing and weighing internal and emotional states appears to be tuned wrongly.

\subsection{Autonomic and Vestibular Roles}

The insula does not merely \textit{perceive} the body's interior; it also helps regulate it. It has strong reciprocal connections with autonomic centres in the brainstem and hypothalamus, and stimulation of different parts of the insula can produce changes in heart rate, blood pressure, and respiration. Damage to the insula, particularly the right insula, has been linked to \textbf{autonomic instability} after stroke, including arrhythmias and sudden cardiac complications, a clinically important reason to recognise insular involvement on imaging.

A separate role worth mentioning is in \textbf{vestibular processing}: the posterior insula and the adjacent parietal operculum host parts of what is sometimes called the \textit{parieto-insular vestibular cortex}, which contributes to the conscious sense of balance and bodily orientation in space.

\subsection{Summary}

The insula is the cortex of the body's interior, the felt self, and the emotional weighting of experience. Hidden beneath the opercula of the frontal, parietal, and temporal lobes, and divided into a posterior region tied to interoception and somatosensory processing of the body's internal state, and an anterior region tied to taste, emotion, and subjective feeling, it stands at the crossroads of sensation, autonomic regulation, and affect. The primary gustatory cortex turns the chemistry of food into the experience of taste; the posterior and mid-insula construct an interoceptive map of the viscera; and the anterior insula re-represents that bodily state as feeling, contributing to disgust, empathy, craving, and the salience of events that matter. Damage produces ageusia, blunted disgust and empathy, autonomic instability, and altered awareness of one's own bodily state. The insula is, more than any other cortex, where the body becomes the self.

That last sentence is worth unpacking, because this section is where the account first arrives at self-awareness proper, and it arrives by a route that is not obvious. The preceding lobes constructed a world: a visual scene, a heard and recognised environment, a space with a body positioned in it. All of that is a model of what is outside. The insula does the same work on what is inside. Interoceptive afferents deliver the state of the heart, lungs, gut, and vasculature to the posterior insula in something close to raw form, and that reading is then re-represented forward until, in the anterior insula, it has become a felt quality rather than a measurement. The operation is the one performed by every association cortex met so far, applied to an unusual input. What comes out of it is not a percept of an object but a percept of the organism doing the perceiving.

This suggests a way of understanding what a self is that requires nothing beyond the machinery already described. Rather than a further thing that possesses the body and observes the world, the felt self can be read as another construction, built from the same materials as the visual scene and by the same sort of layered re-representation, differing only in that its subject matter is the body's own condition. That reading is not a settled result; the strong version of it, on which the right anterior insula hosts a unified representation of the sentient self, was flagged above as contested and remains so. But something like it is the reason this section sits where it does. Up to this point the notes have been assembling a model of the world. From here they are also assembling the thing that model appears to belong to, and the emotional and motivational sections that follow depend on there being an organism whose states can matter to it.


\section{The Basal Ganglia}

Every section so far has been about cortex. The basal ganglia are not cortex. They are a set of nuclei buried deep beneath it, grey matter sunk into the white matter of the hemispheres and reaching down into the midbrain, and they are wired into a loop that begins and ends in the cortex we have already described. Their business is \textbf{action selection}. At any moment the cortex has many possible movements half-prepared; the basal ganglia are the machinery that lets one of them through and holds the rest back.

The most useful thing to fix in mind before any of the anatomy is that this circuit works by \textit{subtraction}, not addition. Its output nuclei fire constantly, and what they do while firing is inhibit the thalamus. Movement is not released by switching something on; it is released by switching that inhibition off. Almost every confusion about the basal ganglia comes from forgetting that the resting state of the system is a brake held down, and that the pathways we are about to trace are competing over whether to lift it.

\subsection{Components}

The basal ganglia are made up of several nuclei, and the naming is genuinely awkward, because anatomists carved these structures up by appearance long before anyone knew what they did. The result is that a single functional unit can carry two names, and a single name can span two functional units. It is worth learning the components alongside their transmitters, because in this circuit the transmitter is what fixes the sign of every connection.

\begin{itemize}
    \item The \textbf{striatum} is the input nucleus: almost everything the cortex sends to the basal ganglia arrives here. It is split by the \textbf{internal capsule}, the great sheet of white matter carrying fibres to and from the cortex, into the \textbf{caudate nucleus} medially and the \textbf{putamen} laterally. The split is anatomical rather than functional; the two are histologically identical, and where they remain joined underneath, the tissue is called the \textbf{nucleus accumbens}. The principal cells of the striatum are the \textbf{medium spiny neurons (MSNs)}, which are GABAergic and therefore inhibitory. They come in two populations, distinguished by which dopamine receptor they carry, and that distinction turns out to be the whole story.
    \item The \textbf{globus pallidus} sits medial to the putamen and has two segments. The \textbf{external segment (GPe)} is an internal relay used by only one of the two pathways. The \textbf{internal segment (GPi)} is an output nucleus, one of the two structures through which the basal ganglia speak to the rest of the brain. Both are GABAergic.
    \item The \textbf{subthalamic nucleus (STN)} is a small lens-shaped nucleus lying beneath the thalamus. It is the exception in this circuit: alone among these nuclei it is \textit{glutamatergic}, and therefore excitatory. Whenever a chain of reasoning about the basal ganglia produces a sign that seems wrong, the STN is usually the step that was missed.
    \item The \textbf{substantia nigra} lies in the midbrain and, confusingly, does two entirely different jobs. The \textbf{pars reticulata (SNr)} is GABAergic and is the second output nucleus, the functional twin of the GPi. The \textbf{pars compacta (SNc)} is dopaminergic and is not part of the loop at all; it projects onto the striatum from the side and modulates it, which is the subject of the last two subsections of this section.
\end{itemize}

One more structure has to be named even though it is not part of the basal ganglia: the \textbf{thalamus}, and specifically its \textbf{ventral anterior (VA)} and \textbf{ventral lateral (VL)} nuclei. These are what the output nuclei act upon, and they project glutamate back up to the motor and premotor cortex. The thalamus is the thing being gated.

Putting these together, the flow of information is a loop:

\[
\text{cortex} \rightarrow \text{striatum} \rightarrow \text{GPi/SNr} \rightarrow \text{thalamus} \rightarrow \text{cortex}
\]

Two features of this loop do all the work in what follows. The first is that the output nuclei, GPi and SNr, are \textbf{tonically active}: they fire continuously, without needing to be told to, and being GABAergic they hold the thalamus under a permanent brake. The second is that when one inhibitory neuron inhibits another inhibitory neuron, the result at the far end is \textbf{disinhibition}, which is functionally indistinguishable from excitation. Two minus signs in series make a plus. The direct pathway is built out of exactly that arrangement.

\subsection{The Direct Pathway}

The direct pathway is the route by which a movement gets released. It runs from cortex to striatum to the output nuclei and out to the thalamus, and its entire trick is the double inhibition just described.

Tracing it synapse by synapse, naming the transmitter at each:

\begin{enumerate}
    \item The \textbf{cortex} fires onto the striatum, releasing \textbf{glutamate}. This is excitatory, and it activates the MSNs that carry \textbf{D1 receptors}.
    \item Those D1 MSNs project to the \textbf{GPi/SNr} and release \textbf{GABA}. This is inhibitory, and it silences output neurons that were, until this moment, firing tonically.
    \item The \textbf{GPi/SNr} normally release \textbf{GABA} onto the thalamus. Now that they have been silenced, that inhibition stops. The brake is lifted.
    \item The \textbf{thalamus}, released, fires up to the cortex on \textbf{glutamate}, exciting the motor areas that were preparing the movement.
\end{enumerate}

Written as a chain, with the sign of each synapse marked:

\[
\text{cortex} \xrightarrow{\;+\;} \text{striatum (D1)} \xrightarrow{\;-\;} \text{GPi/SNr} \xrightarrow{\;-\;} \text{thalamus} \xrightarrow{\;+\;} \text{cortex}
\]

The arithmetic is the point. Steps two and three are both inhibitory, and inhibiting an inhibitor is excitation. The striatum does not excite the thalamus; it removes something that was suppressing the thalamus, and the thalamus, which was ready to fire all along, does the rest. The cortex ends up exciting itself, by way of a detour that takes the brake off.

The net effect of activity in the direct pathway is therefore to \textbf{release the selected movement}. When a motor plan wins through this route, the thalamus is disinhibited, the loop back to cortex closes with excitation, and the action proceeds to execution.

\subsection{The Indirect Pathway}

The indirect pathway is the mirror image: it exists to hold movements back. It starts from the same cortical input and ends at the same output nuclei, but it reaches them by a longer route that inserts an extra inhibitory step and, crucially, an excitatory one. The extra step flips the sign.

Tracing it as before:

\begin{enumerate}
    \item The \textbf{cortex} again fires onto the striatum with \textbf{glutamate}, but this time it activates the MSNs carrying \textbf{D2 receptors}, a separate population from those of the direct pathway.
    \item Those D2 MSNs project to the \textbf{GPe}, not to the output nuclei, releasing \textbf{GABA} and silencing it.
    \item The \textbf{GPe} was itself holding down the subthalamic nucleus with tonic \textbf{GABA}. Silencing the GPe therefore \textit{releases} the STN.
    \item The \textbf{STN}, now disinhibited, fires onto the GPi/SNr with \textbf{glutamate}. This is the exception noted earlier, and it means the output nuclei are now being driven \textit{harder} than at rest.
    \item The \textbf{GPi/SNr}, excited, pour still more \textbf{GABA} onto the thalamus. The brake is pressed down more firmly than before.
\end{enumerate}

As a signed chain:

\[
\text{cortex} \xrightarrow{\;+\;} \text{striatum (D2)} \xrightarrow{\;-\;} \text{GPe} \xrightarrow{\;-\;} \text{STN} \xrightarrow{\;+\;} \text{GPi/SNr} \xrightarrow{\;-\;} \text{thalamus}
\]

Follow the signs and the contrast with the direct pathway becomes clear. Here too there is a double inhibition, striatum onto GPe onto STN, and here too it produces disinhibition. But the structure being disinhibited is excitatory. So instead of the released structure exciting the thalamus, the released structure excites the very nuclei that inhibit the thalamus. The disinhibition is spent on strengthening the brake rather than lifting it.

The net effect of activity in the indirect pathway is therefore to \textbf{suppress movement}. The two pathways begin at the same place and converge on the same place, and they arrive with opposite sign: the direct pathway subtracts from the thalamic brake, the indirect pathway adds to it. What the cortex actually gets back depends on the balance struck between them. That balance is not fixed, and what sets it is dopamine.

\subsection{Dopamine}

Dopamine has appeared already in these notes, listed among the neuromodulators and given the brief gloss of ``movement, reward, motivation''. It is now worth being precise about what it is and, more importantly, about how it acts, because dopamine does not work anything like the glutamate and GABA we have been tracing through the two pathways.

Dopamine is a \textbf{catecholamine}, synthesised from the amino acid \textbf{tyrosine} by way of an intermediate that will be familiar from pharmacology: tyrosine is converted to \textbf{L-DOPA}, and L-DOPA is decarboxylated to dopamine. It is loaded into synaptic vesicles, released by the calcium-triggered mechanism described earlier, and cleared from the cleft mainly by reuptake through the \textbf{dopamine transporter (DAT)}, along the lines set out in the discussion of reuptake and degradation.

The important difference lies in what happens after dopamine binds. Every synapse traced through the direct and indirect pathways used a ligand-gated ion channel, where the receptor is itself the pore and the membrane potential moves within a millisecond. Dopamine receptors are not ion channels at all. They are \textbf{G protein-coupled receptors}, working by the cascade set out when we first met graded potentials: binding throws a switch on a G protein, a GDP is swapped for a GTP, an effector enzyme is engaged, and a second messenger goes on to tune ion channels the cell already has. The consequences established there are exactly the ones that matter here. Dopamine acts slowly, in hundreds of milliseconds to seconds; its effect is amplified far beyond the amount of transmitter released; and it modulates rather than commands, changing how a striatal neuron answers its cortical input rather than firing it directly.

\subsubsection{The D1 and D2 Receptor Families}

Dopamine receptors come in two families, distinguished by which G protein they couple to, and therefore by which direction they push adenylyl cyclase.

\begin{itemize}
    \item The \textbf{D1 family} (D1 and D5) couples to \textbf{G\textsubscript{s}}, the stimulatory G protein. Its $\alpha$ subunit \textit{activates} adenylyl cyclase, cAMP rises, PKA is switched on, and the net effect on the neuron is \textbf{increased excitability}: it becomes more responsive to the glutamate arriving from cortex.
    \item The \textbf{D2 family} (D2, D3, and D4) couples to \textbf{G\textsubscript{i/o}}, the inhibitory G protein. Its $\alpha$ subunit \textit{suppresses} adenylyl cyclase, cAMP falls, PKA activity drops, and the net effect is \textbf{decreased excitability}: the neuron becomes less responsive to that same cortical input.
\end{itemize}

This is the crux, and it is worth stating plainly because it is the single fact on which the next subsection depends. Dopamine is one molecule with one shape, and it has opposite effects on two different cells. Nothing about the transmitter decides this. What decides it is which receptor the cell happens to express, and therefore which G protein sits behind that receptor. The message is identical; the interpretation is not.

\subsection{The Modulatory Pathway}

We now have everything needed to see what the substantia nigra does. Recall from the components that the pars compacta is dopaminergic and stands outside the cortex-striatum-thalamus loop. Its neurons send axons up into the striatum, spreading through it diffusely, in a projection called the \textbf{nigrostriatal pathway}.

This projection carries no motor command. It does not tell the striatum which movement to make, and it is not a third route from cortex to thalamus alongside the other two. What it does is release dopamine broadly across the striatum, where it lands on both MSN populations at once. The direct and indirect pathways select; the nigrostriatal projection sets the terms on which they compete.

The consequence follows directly from the receptor pharmacology of the previous subsection, because the two MSN populations were defined by exactly that difference:

\begin{itemize}
    \item Direct pathway MSNs carry \textbf{D1} receptors. Dopamine acting through G\textsubscript{s} raises cAMP and makes them \textit{more} excitable, so the pathway that lifts the thalamic brake is amplified.
    \item Indirect pathway MSNs carry \textbf{D2} receptors. Dopamine acting through G\textsubscript{i/o} lowers cAMP and makes them \textit{less} excitable, so the pathway that presses the brake down is damped.
\end{itemize}

These two actions are opposite in mechanism and identical in outcome. Dopamine excites the pathway that promotes movement and inhibits the pathway that prevents it, and both effects push the balance the same way. This is the elegance of the arrangement: a single diffuse chemical signal, requiring no wiring specificity whatsoever, can shift the entire system toward action, because the two competing populations have been pre-labelled with receptors of opposite sign.

The result is a \textbf{gain control} on movement. Dopamine does not choose the action; the cortex does that, by deciding which striatal neurons receive glutamate. Dopamine sets how readily whatever has been chosen gets released. When nigrostriatal dopamine is high, the threshold for releasing a movement falls and actions come easily; when it is low, the direct pathway is under-driven and the indirect pathway is under-damped, the brake stays down, and movement becomes harder to initiate.

It is worth separating the two halves of the substantia nigra once more, since the shared name invites confusion. The \textbf{pars reticulata} is a member of the loop, an output nucleus inhibiting the thalamus exactly as the GPi does. The \textbf{pars compacta} is not in the loop at all; it stands to one side and adjusts the gain of the structure that feeds it. Same anatomical name, two entirely different roles.

\subsection{Summary}

The basal ganglia are the brain's mechanism for selecting actions, and they work by releasing a brake rather than by issuing a command. Cortex feeds glutamate into the striatum; the striatum's GABAergic medium spiny neurons act on output nuclei, the GPi and SNr, which fire tonically and inhibit the thalamus without pause. The direct pathway runs striatum to GPi/SNr, and because two inhibitory synapses in series produce disinhibition, its activity lifts the thalamic brake and releases the selected movement. The indirect pathway runs striatum to GPe to subthalamic nucleus to GPi/SNr, and because the disinhibited subthalamic nucleus is glutamatergic, its activity drives the output nuclei harder and presses the brake down, suppressing movement. The two routes converge with opposite sign, and the balance between them is set by dopamine, which acts through neither of the fast ionotropic mechanisms used elsewhere in the circuit but through G protein-coupled receptors, where binding triggers a GDP-for-GTP exchange, an effector enzyme, a second messenger, and a slow, amplified, modulatory change in how the cell answers its other inputs. Because direct pathway neurons carry stimulatory D1 receptors and indirect pathway neurons carry inhibitory D2 receptors, dopamine from the substantia nigra pars compacta amplifies the pathway that permits movement while damping the one that prevents it, and a single diffuse chemical broadcast is enough to set the threshold at which the whole system consents to act.

The capacity this section adds is choosing. At any moment the cortex has many actions partly prepared, and something has to permit one and withhold the others, since a body has only one set of limbs and cannot execute competing plans at once. That is a real problem with a real solution, and the solution turns out to be a brake that is lifted rather than an instruction that is issued. It is worth noticing how little of this resembles the folk picture of deciding to move. There is no stage at which a command is composed; there is a set of candidates, a tonic inhibition holding all of them back, and a competition whose winner is disinhibited. Volition, whatever else is true of it, has at its base a mechanism of exactly this kind, and the mechanism is visible in what happens when it degrades: too little dopamine and the brake stays down, so that movements become hard to initiate despite intact strength and intact intention, which is a dissociation between wanting to move and being able to start that no account of action as simple command can easily accommodate.


\section{The Hypothalamus}

The hypothalamus weighs about four grams, roughly the size of an almond, and is by that measure one of the smallest structures given a section of its own in these notes. It is also the one whose failure the body would notice fastest. It holds temperature, fluid balance, energy intake, and the daily cycle of sleep and waking within the narrow ranges on which every other tissue depends, and it does so continuously and without conscious involvement.

What makes it worth studying as a system rather than as a list of functions is its position. The hypothalamus sits at the junction of three control systems that otherwise have no way of reaching one another: the limbic system, which determines what matters emotionally; the endocrine system, which acts on the body slowly and chemically; and the autonomic nervous system, which acts on it quickly and electrically. None of the three can act on the body in the others' currency. The hypothalamus is what converts between them, and the subsections that follow are best read as five views of that single job.

\subsection{Structure and Components}

The hypothalamus lies beneath the thalamus, from which it takes its name, forming the floor and the lower parts of the lateral walls of the \textbf{third ventricle}. Its boundaries are easier to state than most: the \textbf{optic chiasm} marks its front edge, the \textbf{mammillary bodies} its back edge, and below it narrows into the \textbf{infundibulum}, the stalk by which it hangs the pituitary gland beneath the brain.

Internally it is not a uniform mass but a collection of \textbf{nuclei}, each a cluster of cell bodies with its own job. The naming is easier to hold on to once its logic is clear. Nuclei are grouped in two directions at once: from front to back into \textbf{anterior}, \textbf{tuberal}, and \textbf{posterior} regions, and from the midline outward into \textbf{periventricular}, \textbf{medial}, and \textbf{lateral} zones. A given nucleus is located by both, so that the supraoptic nucleus, for instance, is anterior and lateral.

The nuclei that matter most for what follows are these:

\begin{itemize}
    \item The \textbf{supraoptic} and \textbf{paraventricular nuclei} contain large neurons that manufacture \textbf{vasopressin} and \textbf{oxytocin} and send their axons down the stalk into the posterior pituitary. They are the cells that make the hypothalamus an endocrine organ in its own right.
    \item The \textbf{suprachiasmatic nucleus}, sitting just above the optic chiasm as its name records, is the body's master clock. It receives a direct projection from the retina and uses it to keep an internal daily rhythm aligned with the external light cycle.
    \item The \textbf{arcuate nucleus} lies close to the median eminence and produces several of the releasing hormones that govern the anterior pituitary. It also carries receptors for circulating signals of energy state, which places it at the centre of appetite regulation.
    \item The \textbf{ventromedial nucleus} is associated with satiety, the signal to stop eating, and the \textbf{lateral hypothalamic area} with the opposite, the drive to seek food, along with a broader role in arousal and motivated behaviour.
    \item The \textbf{preoptic area}, at the front, holds the thermosensitive neurons that monitor blood temperature and the populations involved in initiating sleep.
    \item The \textbf{mammillary bodies}, two small round swellings on the underside at the back, are not endocrine or autonomic at all. They belong to a memory circuit, and they reappear in the two subsections that follow.
\end{itemize}

\subsection{Limbic Involvement}

The \textbf{limbic system} is the set of structures concerned with emotion, motivation, and memory, and the hypothalamus is best understood as its output stage. Structures such as the amygdala and hippocampus determine what a situation means; they cannot, by themselves, do anything about it. The hypothalamus is where meaning is converted into bodily change, and this is the anatomical reason emotion is felt physically rather than merely registered.

The classical description of this circuitry is the \textbf{Papez circuit}, proposed as a loop connecting emotion with memory. It runs: hippocampus, then out along the \textbf{fornix}, a prominent arching bundle of white matter, to the \textbf{mammillary bodies}; from there along the \textbf{mammillothalamic tract} to the \textbf{anterior nucleus of the thalamus}; onward to the \textbf{cingulate cortex} lying above the corpus callosum; and from the cingulate back to the hippocampus, closing the loop. The circuit as Papez conceived it is no longer taken to be the seat of emotion in the way he intended, but the connections themselves are real and the loop remains the standard way of describing how these structures relate.

The amygdala reaches the hypothalamus by two routes of its own, considered in more detail among the pathways below. Its influence is what allows a stimulus judged threatening to produce, within a second or so, the entire pattern of bodily changes associated with fear, without any deliberate decision being taken. Feeding, defensive behaviour, aggression, and sexual behaviour all depend on this limbic traffic reaching a structure that can act on the body.

\subsection{Endocrine Involvement}

The hypothalamus governs the \textbf{pituitary gland}, and it does so by two mechanisms that are entirely separate anatomically. Distinguishing them is the key to the whole subsection, because the two lobes of the pituitary are not variations on one arrangement but genuinely different kinds of structure.

The \textbf{posterior pituitary} is not really a gland. It is nervous tissue, an extension of the hypothalamus itself. The large \textbf{magnocellular} neurons of the supraoptic and paraventricular nuclei manufacture \textbf{oxytocin} and \textbf{vasopressin}, package them, and transport them down axons running through the stalk. The axon terminals sit in the posterior lobe among a bed of capillaries, and when the neuron fires, the hormone is released directly into the bloodstream. The mechanism is precisely the one described earlier for transmitter release, differing only in where the vesicle contents end up: the blood rather than a synaptic cleft. This is \textbf{neurosecretion}, and the neuron is itself the secretory cell.

The \textbf{anterior pituitary} is a true gland, made of glandular rather than nervous tissue, and no hypothalamic axon reaches it. The connection is instead vascular. Smaller \textbf{parvocellular} neurons release \textbf{releasing} and \textbf{inhibiting hormones} into a specialised set of vessels at the \textbf{median eminence}, the \textbf{hypophyseal portal system}, which carries them the short distance down to the anterior lobe. There they instruct the glandular cells to release, or to withhold, the hormones that act on targets throughout the body:

\begin{itemize}
    \item \textbf{CRH} drives release of \textbf{ACTH}, which acts on the adrenal cortex to release \textbf{cortisol}. This is the axis that matters most in these notes, and it reappears in the sections on emotion and memory.
    \item The remaining axes follow the same pattern with different targets: \textbf{TRH} drives \textbf{TSH} and the thyroid, \textbf{GnRH} drives \textbf{LH} and \textbf{FSH} and the gonads, and \textbf{GHRH} drives \textbf{growth hormone}, opposed by \textbf{somatostatin}.
    \item \textbf{Dopamine}, met earlier as a neurotransmitter, acts here as an inhibiting hormone, holding \textbf{prolactin} release in check. The same molecule doing an entirely different job in a different setting is the pattern noted when transmitters were introduced.
\end{itemize}

Two features of this arrangement are worth drawing out. The first is the economy of the portal system: by delivering hormones straight to their target through a private circulation, the hypothalamus achieves high local concentrations using quantities far too small to have any effect if released into the general bloodstream. The second is that none of these axes runs one way. The hormones produced by the peripheral glands circulate back to the hypothalamus and the pituitary and suppress the signals that called for them. Each axis is therefore a closed negative feedback loop, of the same general form as the feedback mechanisms met at the synapse, operating over hours rather than milliseconds.

\subsection{Autonomic Involvement}

The hypothalamus has been called the head ganglion of the autonomic nervous system, and the description is apt provided it is read as a statement about hierarchy. The hypothalamus does not innervate the heart, the gut, or the blood vessels itself. It issues instructions to the brainstem and spinal cord nuclei that do, chiefly through descending projections to the autonomic centres of the medulla and to the \textbf{intermediolateral cell column}, the strip of preganglionic sympathetic neurons running down the thoracic cord.

The two divisions introduced earlier are represented unevenly across the structure. Anterior and medial regions bias output toward the \textbf{parasympathetic} division, and posterior and lateral regions toward the \textbf{sympathetic}. This is a tendency rather than a map, and stimulation experiments produce mixed patterns rather than clean ones, but it is a useful first approximation.

What the hypothalamus contributes that the brainstem cannot is the \textbf{set point}. Autonomic reflexes at brainstem level are corrective but have no independent notion of what the correct value should be; the hypothalamus supplies it, and can change it. Two examples show the whole arrangement in one structure:

\begin{itemize}
    \item In \textbf{thermoregulation}, thermosensitive neurons in the preoptic area measure the temperature of the blood passing through and compare it against a set point. Deviation upward recruits sweating and cutaneous vasodilation; deviation downward recruits vasoconstriction, shivering, and, over a longer period, endocrine changes that raise metabolic rate. The sensing and the effector command sit within millimetres of one another.
    \item In \textbf{osmoregulation}, osmoreceptive neurons detect the concentration of the blood. A rise triggers vasopressin release from the posterior pituitary, which instructs the kidney to retain water, and simultaneously generates thirst, which is a behavioural rather than an autonomic response. A single measurement thus drives an endocrine, an autonomic, and a behavioural limb at once.
\end{itemize}

That last point generalises. Because the hypothalamus commands all three, its responses are typically coordinated across them rather than confined to one, which is why hunger, cold, and fear each produce a bodily state rather than an isolated adjustment.

\subsection{Pathways}

The hypothalamus is a small structure with a large number of connections, and several of its major tracts have been named individually because they are visible to the naked eye in dissection.

The \textbf{fornix} is the principal route from the \textbf{hippocampus}. It is a large arched bundle, one of the most conspicuous white matter tracts in the brain, which sweeps from the hippocampus in the temporal lobe up and forward beneath the corpus callosum before descending to terminate mainly in the mammillary bodies. It is the hippocampus's major output, and it is the reason a structure associated with memory has a direct line into one associated with drive. From the mammillary bodies the signal continues along the \textbf{mammillothalamic tract} to the anterior thalamus, as traced in the Papez circuit above.

The \textbf{amygdala} connects by two routes that differ in shape rather than in purpose. The \textbf{stria terminalis} is long and thin and takes an indirect arcing course that roughly parallels the fornix, following the curve of the lateral ventricle. The \textbf{ventral amygdalofugal pathway} takes the direct route underneath, a short and diffuse projection into the medial hypothalamus. Both carry the amygdala's assessment of emotional significance to the structure that can act on it, and their targets substantially overlap.

Connections with the \textbf{prefrontal cortex} are reciprocal, and arise mainly from its medial and orbitofrontal regions, the parts already described as most concerned with emotion, value, and social behaviour. Their significance is regulatory. The hypothalamus generates drives that are, by themselves, indifferent to circumstance, and prefrontal input is what allows context, social consequence, and long-term goals to modulate them. Hunger does not compel eating in every setting, and the anatomical basis of that restraint runs along this connection.

Two further pathways are worth naming. The \textbf{medial forebrain bundle} runs longitudinally through the lateral hypothalamic area, linking the brainstem below with the basal forebrain and cortex above, and carries much of the traffic associated with reward and motivated behaviour. The \textbf{retinohypothalamic tract} is a small, direct projection from the retina to the suprachiasmatic nucleus; it is the anatomical means by which the master clock is kept aligned to the light cycle, and it is entirely separate from the visual pathway described in the occipital lobe section, carrying information about ambient light rather than about images.

\subsection{Summary}

The hypothalamus is a structure of a few grams occupying the floor of the third ventricle, and its importance lies in its position at the junction of three systems that cannot otherwise communicate. Its nuclei, grouped from front to back and from the midline outward, each take responsibility for a different regulated variable: the supraoptic and paraventricular for fluid balance and neurosecretion, the suprachiasmatic for circadian timing, the arcuate, ventromedial, and lateral areas for energy intake, the preoptic for temperature and sleep, and the mammillary bodies for the memory circuitry they share with the hippocampus. As the output stage of the limbic system, it converts emotional significance arriving from the amygdala and hippocampus into bodily change. As an endocrine organ, it commands the pituitary by two separate mechanisms, sending axons directly into the posterior lobe to release oxytocin and vasopressin into the blood, and reaching the anterior lobe indirectly through the hypophyseal portal system with releasing and inhibiting hormones that govern the peripheral endocrine axes. As the head of the autonomic hierarchy, it supplies the set points that brainstem reflexes lack, and coordinates endocrine, autonomic, and behavioural responses to a single measurement. Its named tracts, the fornix, the stria terminalis and ventral amygdalofugal pathway, the mammillothalamic tract, the medial forebrain bundle, and the retinohypothalamic tract, are the anatomical form of those relationships. For a structure so small, the hypothalamus is where a remarkable amount of what the body does without being asked is actually decided.

The previous section supplied a mechanism for selecting among actions, and left open the question of why any action should be preferred to any other. This one answers it. A system that merely represented the world and could move would have no reason to do anything, and the hypothalamus is where reasons come from: it holds a set of variables that must stay within limits, detects departures from them, and generates the states that push behaviour toward correcting them. Hunger, thirst, cold, and the drive to sleep are not perceptions of the body but demands made upon the organism, and they are constructed here from measurements taken here. That the same structure supplies both the set point and the command is what allows a deviation to become a motive rather than merely a reading. Wanting, in its simplest and most physiological form, starts as a regulated variable drifting out of range.

\section{The Thalamus}

Having spent the previous section beneath it, we now turn to the structure the hypothalamus is named after. The \textbf{thalamus} is a pair of egg-shaped masses of grey matter, one either side of the third ventricle, and it accounts for roughly four fifths of the diencephalon by volume. It sits almost exactly at the centre of the brain, and its position is not incidental: nearly everything travelling between the cortex and the rest of the nervous system passes through it.

The standard description is that the thalamus is the brain's relay station, and as a first approximation that is right. With the single exception of smell, every sensory stream synapses in a specific thalamic nucleus before it is allowed to reach the cortex. But a relay station in the ordinary sense is a passive thing, and the thalamus is not passive. Only a minority of the synapses on a thalamic relay neuron come from the ascending pathway it is supposedly relaying; the majority come from the cortex above, the brainstem below, and the thalamus's own inhibitory shell. What arrives at the cortex is therefore not a copy of what arrived at the thalamus but a version of it that has been weighted by attention, by arousal, and by what the cortex was already expecting. The thalamus decides what gets through.

It is also more than a sensory structure. The same anatomy that carries sensation upward carries the return legs of the motor and limbic loops met in earlier sections: the basal ganglia and the cerebellum both reach the motor cortex through it, and the memory circuit running out of the hippocampus passes through it too. The three subsections that follow take those three roles in turn, and the organising claim of the section is that they are the same mechanism applied to different traffic.

\subsection{Structure and Components}

Each thalamus lies with its medial surface forming the upper wall of the \textbf{third ventricle} and its lateral surface pressed against the \textbf{posterior limb of the internal capsule}, the great sheet of white matter through which its fibres leave. In most people the two halves are joined across the ventricle by the \textbf{interthalamic adhesion}, a bridge of grey matter which, despite appearances, carries little traffic between the sides. Posteriorly the thalamus swells into the \textbf{pulvinar}, which overhangs the midbrain, and tucked beneath the pulvinar sit two small bumps that have already appeared in these notes as the last stops on the visual and auditory pathways: the \textbf{lateral} and \textbf{medial geniculate nuclei}.

Internally the thalamus is divided by a sheet of white matter, the \textbf{internal medullary lamina}, which runs front to back through the structure and splits at its anterior end into a Y. That Y is the whole of the organising scheme. The two arms enclose the \textbf{anterior group}; the space medial to the stem holds the \textbf{medial group}; and everything lateral to it forms the large \textbf{lateral group}, itself divided into a dorsal tier of association nuclei and a ventral tier of relay nuclei. Three further populations sit outside this scheme: the \textbf{intralaminar nuclei}, embedded within the lamina itself, the \textbf{midline nuclei} against the ventricle, and the \textbf{reticular nucleus}, a thin shell of cells wrapped around the lateral surface within the \textbf{external medullary lamina}.

The reticular nucleus deserves separating from the rest, because it is the exception that explains the system. Every other thalamic nucleus projects to the cortex. The reticular nucleus does not: it is \textbf{GABAergic}, and its targets are the other thalamic nuclei. Corticothalamic and thalamocortical fibres both pass through it on their way, giving off collaterals as they go, so it is informed of everything crossing in either direction and inhibits the relay nuclei accordingly. It is, in other words, the gate itself, and it is the structural reason the thalamus can be selective about what it forwards.

It is useful to sort the remaining nuclei into three functional classes rather than to memorise them by location:

\begin{itemize}
    \item \textbf{Relay} (or specific) nuclei receive a defined ascending pathway and project to a defined cortical area. The sensory relays (VPL, VPM, LGN, MGN) and the motor relays (VA, VL) are all of this kind, and their maps are preserved: neighbouring cells represent neighbouring parts of the body, the retina, or the cochlea.
    \item \textbf{Association} nuclei receive most of their input from the cortex rather than from an ascending tract, and project back to association cortex. The \textbf{pulvinar}, the \textbf{mediodorsal nucleus}, and the \textbf{lateral dorsal} and \textbf{lateral posterior} nuclei belong here. They serve the integrative regions of cortex rather than the primary ones.
    \item \textbf{Nonspecific} nuclei, the intralaminar and midline groups, project diffusely to wide areas of cortex and to the striatum. They carry arousal and attentional signals from the brainstem reticular formation rather than any particular modality, and they are the route by which the level of consciousness is set.
\end{itemize}

The fibres leaving the thalamus are gathered into four \textbf{thalamic radiations} which fan through the internal capsule to reach the frontal, parietal, occipital, and temporal cortices. Because these radiations are packed so tightly, a small lesion here can produce a combination of deficits that would require a much larger cortical lesion to reproduce.

\subsection{Limbic Involvement}

The limbic role of the thalamus is carried by the anterior and medial groups, and it is a role in memory and in the emotional colouring of experience rather than in sensation.

The \textbf{anterior nuclear group} is the thalamic member of the \textbf{Papez circuit} described in the previous section. It receives the \textbf{mammillothalamic tract} from the mammillary bodies, which are themselves fed by the fornix from the hippocampus, and it projects onward to the \textbf{cingulate cortex}, from which the loop returns to the hippocampus. The circuit is no longer read as the seat of emotion in the way Papez intended, but as a memory pathway it holds up well, and the clinical evidence is unambiguous: damage to the mammillary bodies and the anterior and mediodorsal thalamus, as occurs in the thiamine deficiency of \textbf{Korsakoff's syndrome}, produces a dense anterograde amnesia while leaving perception and intelligence largely intact. A structure usually described as a relay for sensation turns out to be a necessary component of laying down new episodic memory.

The \textbf{mediodorsal nucleus (MD)} is the largest of the medial group and is the thalamic partner of the \textbf{prefrontal cortex}. Its connections with the prefrontal cortex are reciprocal and extensive, so much so that the extent of the prefrontal cortex was once defined by where MD projects. It also receives from the \textbf{amygdala}, from the olfactory cortex, and from the hypothalamus. Two consequences follow. The first is that MD is the point at which the amygdala's assessment of emotional significance reaches the executive machinery of the frontal lobe, which is what allows an emotionally salient event to capture deliberation rather than merely producing a bodily response. The second concerns olfaction: smell is the sense that bypasses the thalamus on its way to cortex, but it does not avoid the thalamus altogether. It reaches MD after its first cortical stop, and from MD reaches the orbitofrontal cortex, which is where conscious discrimination of odours occurs. The thalamic step is postponed rather than absent.

Two smaller contributions complete the picture. The \textbf{midline nuclei} connect with the hippocampus, amygdala, and hypothalamus, and are involved in arousal and in visceral aspects of emotion. The \textbf{intralaminar nuclei} receive the spinothalamic fibres that do not terminate in VPL and project to the \textbf{anterior cingulate} and \textbf{insula}. This is the anatomical basis of a distinction met in the insula section: the sensory-discriminative component of pain, its location and intensity, travels the lateral route through VPL to S1, while the affective component, the part that makes pain unpleasant rather than merely detected, travels the medial route through the intralaminar nuclei to the cingulate and insula. The two can be separated by lesions, and patients in whom the medial route is interrupted may report that the pain is still there but no longer bothers them.

\subsection{Sensory Involvement}

The sensory role is the one the thalamus is best known for, and the arrangement is systematic: each modality has its own nucleus, each nucleus projects to its own primary cortex, and each preserves the map of the receptor surface along the way.

\begin{itemize}
    \item \textbf{Ventral posterolateral (VPL)} receives the dorsal column--medial lemniscus and spinothalamic pathways from the \textbf{body} and projects to S1, somatotopically arranged.
    \item \textbf{Ventral posteromedial (VPM)} does the same for the \textbf{face}, receiving the trigeminal pathways. Its most medial part, \textbf{VPMpc}, receives \textbf{taste} from the solitary nucleus and projects to the gustatory cortex of the insula.
    \item \textbf{Lateral geniculate nucleus (LGN)} receives the optic tract and projects through the optic radiations to V1, retinotopically arranged and layered so that input from the two eyes is kept separate until cortex.
    \item \textbf{Medial geniculate nucleus (MGN)} receives the auditory pathway from the inferior colliculus and projects to A1, tonotopically arranged.
    \item \textbf{VMpo} and neighbouring nuclei relay the interoceptive and small-fibre afferents to the posterior insula, as described in that section.
    \item \textbf{Vestibular} information relays through the ventral posterior region to the parietoinsular vestibular cortex, the one major sensory stream without a nucleus named after it.
\end{itemize}

Smell remains the sole exception, and the reason is developmental: the olfactory system is the oldest of the sensory systems and was in place before the thalamic arrangement evolved around it.

The more interesting question is what the relay actually does, since a pathway that simply copied its input to cortex would not need a synapse in the middle. Two mechanisms answer it. The first is the inhibitory gating already described: the reticular nucleus, informed by cortical feedback about what is currently relevant, suppresses relay cells carrying information that is not. Attention, at this level, is a matter of which thalamic neurons are permitted to fire. The second is that relay neurons have two modes of firing. In the \textbf{tonic mode}, which prevails during waking, they fire in a graded way that faithfully transmits the detail of their input. In the \textbf{burst mode}, which prevails during drowsiness and slow-wave sleep, they fire in stereotyped bursts that carry little information about the stimulus. The transition between them is controlled by the brainstem arousal systems acting on the membrane potential of the relay cell, and it is why the same sound will wake a sleeper who is being called by name and be entirely absent from the record of the night otherwise. The rhythmic activity of the reticular nucleus during the second mode generates the \textbf{sleep spindles} visible on the EEG.

The association nuclei of the dorsal tier work on sensory information after the primary cortices have had it. The \textbf{pulvinar}, the largest nucleus in the human thalamus, is connected chiefly with the visual, parietal, and temporal association areas and is heavily implicated in visual attention and in the selection of targets for eye movements; lesions of the pulvinar impair the ability to attend to one object among several without impairing the ability to see any of them. The \textbf{lateral posterior} and \textbf{lateral dorsal} nuclei serve the parietal and cingulate cortices in a similar integrative fashion.

The clinical consequences follow directly from the anatomy. Infarction of the VPL and VPM territory produces a contralateral loss of all sensory modalities on the body and face together, a pattern that no cortical lesion of comparable size can produce. In a proportion of these patients, the numbness is followed weeks later by \textbf{thalamic pain syndrome} (\textbf{D\'ejerine--Roussy syndrome}), in which the deafferented side becomes the source of severe spontaneous burning pain, and in which ordinary touch may be experienced as painful. Pain, on this evidence, is not merely reported by the thalamus but is shaped there.

\subsection{Motor Involvement}

The thalamus issues no motor commands. It has no descending projection to the spinal cord and no access to a motor neuron; everything it contributes to movement it contributes by talking to the motor cortex. Its importance is that it is the point at which the two great subcortical motor systems, the basal ganglia and the cerebellum, hand their output back to the cortex. Both are loops, and the thalamus is the return leg of both.

The relevant nuclei are the \textbf{ventral anterior (VA)} and \textbf{ventral lateral (VL)}, which sit in the ventral tier just in front of VPL. They receive from two sources that remain largely segregated within them:

\begin{itemize}
    \item The \textbf{basal ganglia} reach VA and the anterior part of VL by the GABAergic output of the \textbf{GPi} and \textbf{SNr}. As set out in that section, this input is tonically active and inhibitory, so the thalamic contribution to movement is made by \textit{having its brake released} rather than by being driven. When the direct pathway silences the GPi, VA/VL is disinhibited and excites the premotor and supplementary motor areas on glutamate; when the indirect pathway strengthens GPi firing, the brake is pressed harder and the movement is suppressed. The thalamus is the structure on which the whole arithmetic of that section is performed.
    \item The \textbf{cerebellum} reaches the posterior part of VL by an excitatory route. Purkinje output leaves through the \textbf{dentate} and other deep nuclei, ascends in the \textbf{superior cerebellar peduncle}, decussates in the midbrain, and terminates in VL, which projects on to the primary motor cortex. This is the anatomical reason the cerebellum's corrections, which arise from comparing intended with actual movement, are applied by adjusting the next cortical command rather than by acting on the muscle directly.
\end{itemize}

Two points are worth drawing out. The first is that the two inputs are opposite in sign and largely separate in territory, which is why disorders of the basal ganglia and of the cerebellum produce quite different movement disorders despite converging on the same nucleus. The second is that both loops are closed: the cortex projects down into the striatum and the pons, and receives back through the thalamus, so the motor cortex is in continuous conversation with both systems rather than issuing commands into them.

There is a clinical corollary that has become therapeutically important. Because VL is the funnel through which both systems act, interrupting it interrupts the abnormal signal as well as the normal one. The \textbf{ventral intermediate nucleus (VIM)}, a subdivision of VL carrying the cerebellar input, is the standard target for \textbf{deep brain stimulation} in essential tremor, and stimulation there can abolish a tremor that no cortical or cerebellar intervention would reach. Lesions of the motor thalamus more generally produce ataxia, tremor, and, when the lesion is larger, the peculiar syndrome of a limb that is neither weak nor uncoordinated but is simply not used.

\subsection{Summary}

The thalamus is a paired mass of grey matter at the centre of the brain through which almost all traffic between the cortex and the rest of the nervous system passes. Its internal medullary lamina divides it into anterior, medial, and lateral groups, with the intralaminar, midline, and reticular populations sitting outside that scheme; of these the reticular nucleus is the key structure, being the one thalamic nucleus that inhibits the others rather than projecting to cortex, and therefore the gate through which selection is exercised. Its limbic role runs through the anterior group, the thalamic member of the Papez circuit and a structure whose damage produces amnesia, and through the mediodorsal nucleus, which links the amygdala and the olfactory cortex to the prefrontal cortex; the intralaminar nuclei carry the affective component of pain to the cingulate and insula. Its sensory role is the familiar one of a modality-specific relay, VPL and VPM for body and face, LGN for vision, MGN for hearing, VPMpc for taste, with smell alone excepted, but the relay is an active one: cortical feedback and reticular inhibition determine what is forwarded, and the switch between tonic and burst firing determines whether the cortex is being informed about the world at all. Its motor role is to close the loops of the basal ganglia and the cerebellum, receiving inhibition from the GPi and SNr into VA and anterior VL and excitation from the dentate into posterior VL, and returning both to the motor cortices. The nuclei and their connections are gathered in Table~\ref{tab:thalamus}. The unifying idea is that the thalamus is not a junction box but a gate, and that consciousness of a sensation, the recollection of an event, and the release of a movement all depend on something being permitted through it.

Two of the results here bear directly on the larger argument, and both concern awareness rather than sensation. The first is that attention, described in earlier sections as a capacity of cortex, has a mechanism at this level: the reticular nucleus, informed by cortical feedback about what currently matters, inhibits the relay cells carrying what does not. Attending to something is not a matter of examining an internal image more closely. It is a matter of which thalamic neurons are permitted to fire, and the material that is not selected does not arrive in a diminished form but largely fails to arrive at all. This is why the peripheral detail of an absorbing moment is not merely forgotten afterwards; much of it was never forwarded in the first place.

The second is the distinction between tonic and burst firing, which is as close as these notes come to a physiological switch for awareness itself. The same relay cell, receiving the same input, either transmits its detail faithfully or emits stereotyped bursts carrying almost nothing about the stimulus, and which of these it does is set by the brainstem arousal systems acting on its membrane potential. The sensory pathway does not shut down during sleep; it continues to carry signals that no longer inform the cortex of anything. Whether there is experience of a stimulus therefore turns not on whether the stimulus was received, nor on whether cortex is intact, but on the mode in which an intervening cell happens to be firing.

{\footnotesize
\renewcommand{\arraystretch}{1.3}
\begin{longtable}{p{2.7cm} p{2.9cm} p{3.3cm} p{2.6cm} p{2.2cm}}
\caption{The principal thalamic nuclei, grouped as the internal medullary lamina divides them, with their major input, their cortical target, and the role each plays. The reticular nucleus is the sole entry with no cortical projection, which is what makes it the gate rather than a relay.}
\label{tab:thalamus} \\
\toprule
\textbf{Group} & \textbf{Nucleus} & \textbf{Main input} & \textbf{Main output} & \textbf{Role} \\
\midrule
\endfirsthead
\multicolumn{5}{l}{\emph{\small Table~\ref{tab:thalamus} continued from previous page}} \\
\toprule
\textbf{Group} & \textbf{Nucleus} & \textbf{Main input} & \textbf{Main output} & \textbf{Role} \\
\midrule
\endhead
\bottomrule
\endfoot
\bottomrule
\endlastfoot
Anterior & Anterior nuclei & Mammillothalamic tract & Cingulate cortex & Limbic (memory) \\
\midrule
Medial & Mediodorsal (MD) & Amygdala, olfactory cortex, hypothalamus & Prefrontal cortex & Limbic (emotion, executive) \\
\midrule
Lateral, dorsal tier & Lateral dorsal (LD) & Hippocampal, cingulate & Cingulate, parietal & Limbic association \\
 & Lateral posterior (LP) & Parietal cortex & Parietal association & Sensory integration \\
 & Pulvinar & Visual, parietal, temporal cortex & Visual and temporal association & Visual attention \\
\midrule
Lateral, ventral tier & Ventral anterior (VA) & GPi/SNr (basal ganglia) & Premotor, SMA & Motor (gating) \\
 & Ventral lateral (VL) & GPi/SNr; dentate (cerebellum) & Primary motor cortex & Motor (execution, coordination) \\
 & Ventral posterolateral (VPL) & Medial lemniscus, spinothalamic (body) & S1 & Somatic sensation \\
 & Ventral posteromedial (VPM) & Trigeminal (face); solitary nucleus (taste, via VPMpc) & S1; gustatory insula & Facial sensation, taste \\
 & Medial geniculate (MGN) & Inferior colliculus & A1 & Hearing \\
 & Lateral geniculate (LGN) & Optic tract & V1 & Vision \\
\midrule
Intralaminar & Centromedian, parafascicular & Reticular formation, spinothalamic & Striatum, diffuse cortex & Arousal, affective pain \\
\midrule
Midline & Paraventricular, others & Hypothalamus, brainstem & Hippocampus, amygdala & Visceral, emotional \\
\midrule
Reticular & Reticular nucleus & Thalamocortical and corticothalamic collaterals & Other thalamic nuclei (GABA) & Gating, sleep spindles \\
\end{longtable}
}



\section{The Limbic System}

Every section so far has taken a piece of the nervous system and asked what it does. This one works in the opposite direction. The \textbf{limbic system} is not a region that can be pointed to on a specimen, and almost nothing in it is new: the hypothalamus and the thalamus have each had a section to themselves, the amygdala and the hippocampus were met in passing in the temporal lobe, and the entorhinal cortex was introduced there as the olfactory association area. What the term names is not fresh tissue but a \textit{set of connections} among tissue already described, together with a set of functions that no one of the participating structures could carry out alone.

The name comes from Broca, who in 1878 described \textit{le grand lobe limbique}, the ``bordering lobe'', after the Latin \textit{limbus}, a hem or a margin. What he was pointing at was a ring of cortex on the medial surface of each hemisphere that curls around the corpus callosum and the upper brainstem like a collar, dividing the neocortex outside it from the diencephalon within. Broca meant the description anatomically and supposed the ring to be concerned with smell. It was \textbf{Papez} in 1937 who proposed that part of it formed a circuit for emotion, and \textbf{MacLean} in the 1950s who broadened that proposal into the ``limbic system'' as it is now generally taught.

It is worth saying plainly that the term is contested. There is no agreed membership list, no boundary that anatomy itself supplies, and no single function that all the candidate structures share; several authors have argued the phrase should be retired altogether on the grounds that it groups together structures whose connections are more varied than the label suggests. The reason it survives, and the reason it is worth a section here, is that the structures conventionally included are in fact densely interconnected with one another, and that they are jointly responsible for a set of capacities that the rest of the nervous system does not supply. The components taken up below are the ones most consistently included; the nucleus accumbens, the ventral tegmental area, and the periaqueductal grey appear in the functional subsections as close associates rather than as members.

The organising idea of the section is this. The cortex described in earlier sections tells the organism what is present: an object, a face, a sound, a smell, a limb in a particular position. That is not enough to act on. Two further questions have to be answered, and the limbic system is where each is answered. The first is \textit{what happened here before}, which is memory, and it is the business of the hippocampus and the circuit running out of it. The second is \textit{what does this mean for me}, which is emotion and value, and it is the business of the amygdala and the cortex it reports to. A third structure, the hypothalamus, converts whatever answer is arrived at into bodily change. Perception, memory, valuation, and response: the limbic system supplies the middle two and commands the last.

\subsection{Components}

\subsubsection{Limbic Lobe}

The \textbf{limbic lobe} is Broca's original ring, and it is the cortical component of the system. It is best appreciated on a midsagittal section, where it appears as a continuous band of cortex sweeping around the corpus callosum: forward and up from beneath its front end, back along its upper surface, down around its posterior end, and then forward again along the medial temporal lobe, closing the circle beneath. Three named regions make up the band.

The \textbf{cingulate gyrus} runs immediately above the corpus callosum and forms the upper half of the ring. Beneath it lies the \textbf{cingulum bundle}, the white matter tract that carries traffic along the ring and connects its ends. The gyrus is not functionally uniform, and the division between its front and back halves matters:

\begin{itemize}
    \item The \textbf{anterior cingulate cortex (ACC)} is concerned with conflict, effort, error, and the affective weighting of experience. It is the cortical partner of the anterior insula in the \textbf{salience network} described in the insula section, and it is the destination of the medial pain pathway traced in the thalamus section, which is why it carries the unpleasantness of pain rather than its location. It also has direct descending connections to autonomic centres, making it one of the few pieces of cortex that acts on the body without an intermediary.
    \item The \textbf{posterior cingulate cortex} and the \textbf{retrosplenial cortex} behind it are concerned with spatial orientation, the retrieval of autobiographical memory, and self-referential thought. The retrosplenial cortex in particular is a waypoint between the hippocampus and the parietal cortex, and it is where a memory anchored to the body's own frame of reference is translated into one anchored to the world.
\end{itemize}

The \textbf{parahippocampal gyrus} forms the lower half of the ring, on the medial surface of the temporal lobe, and is continuous with the cingulate gyrus behind by way of the narrow \textbf{isthmus}. Its anterior part is the \textbf{entorhinal cortex} (Brodmann area 28), already met as the olfactory association cortex, and it is more than that: it is the principal gateway through which the neocortex reaches the hippocampus and the route by which the hippocampus answers. Its posterior part contributes to the recognition of scenes and spatial context. At the front the gyrus hooks sharply backwards to form the \textbf{uncus}, which overlies the amygdala and carries the primary olfactory cortex.

Two smaller regions complete the lobe. The \textbf{subcallosal} or \textbf{paraterminal gyrus} lies beneath the front end of the corpus callosum and is continuous with the septal area described below. The \textbf{hippocampal formation} itself is rolled into the medial temporal lobe beneath the parahippocampal gyrus, and although it is buried it belongs to this ring rather than to the neocortex outside it.

Histologically the ring is distinct from the neocortex it borders. Neocortex has six layers; the cortex of the limbic lobe is transitional, and the hippocampal formation at its core has only three. This is one of the reasons the ring was recognised as a unit long before anyone knew what it did, and the reason it is sometimes called \textbf{allocortex} and \textbf{mesocortex} in contrast to the \textbf{neocortex} outside it.

\subsubsection{Hippocampus}

The \textbf{hippocampus} lies in the floor of the inferior horn of the lateral ventricle, curled into the medial temporal lobe, and takes its name from the seahorse it resembles in cross-section. It is the structure most closely associated with the formation of new memories, and it is the best-understood piece of the limbic system by a considerable margin.

Strictly, the term \textbf{hippocampal formation} covers three parts that differ in structure and are worth distinguishing: the \textbf{dentate gyrus}, the \textbf{hippocampus proper} (subdivided into the fields \textbf{CA1} to \textbf{CA3}, the initials standing for \textit{cornu ammonis}, Ammon's horn), and the \textbf{subiculum}, which is the transition back towards the entorhinal cortex. Unlike the six-layered neocortex, all of these are three-layered, with a single dense band of principal cells, which is what makes the internal circuitry so unusually easy to trace.

That circuitry is the \textbf{trisynaptic circuit}, and it is essentially unidirectional:

\[
\text{entorhinal cortex} \rightarrow \text{dentate gyrus} \rightarrow \text{CA3} \rightarrow \text{CA1} \rightarrow \text{subiculum}
\]

Highly processed information from association cortex throughout the brain converges on the entorhinal cortex, which sends it along the \textbf{perforant path} to the dentate gyrus. The dentate projects by \textbf{mossy fibres} to CA3, CA3 projects by the \textbf{Schaffer collaterals} to CA1, and CA1 passes to the subiculum, from which output leaves the formation. The projection is glutamatergic at every step. CA3 is additionally distinguished by a dense network of recurrent connections onto itself, an architecture well suited to completing a whole pattern from a fragment of it, which is a reasonable description of what recollection feels like from the inside.

Output leaves by two routes. The \textbf{fornix}, described in the hypothalamus section, is the great arching bundle carrying subicular and CA1 fibres forward to the mammillary bodies and the septal area. A second, quieter route runs back to the entorhinal cortex and from there to the same association cortices that supplied the input, which is the anatomical form of the hippocampus talking back to the cortex about what the cortex has just sent it.

Two properties of hippocampal neurons deserve naming here because the memory subsection depends on them. The first is \textbf{long-term potentiation (LTP)}, the lasting strengthening of a synapse that has been repeatedly and strongly driven, which was discovered in this pathway and remains the leading candidate mechanism for memory at the cellular level. The second is spatial coding: hippocampal \textbf{place cells} fire when the animal occupies a particular location, and \textbf{grid cells} in the entorhinal cortex fire in a regular hexagonal lattice across an environment, so that the structure carries a metric map of space as well as a record of events.

The hippocampus is also conspicuously fragile. CA1 is among the most hypoxia-sensitive tissue in the brain and is characteristically lost after cardiac arrest; the entorhinal cortex and hippocampus are the earliest cortical casualties of Alzheimer's disease, which is why its first symptom is a failure to lay down new memories; and the medial temporal lobe is the commonest origin of focal epilepsy in adults, where the associated \textbf{mesial temporal sclerosis} is a scarring and shrinkage of exactly these fields.

\subsubsection{Amygdala}

The \textbf{amygdala} is an almond-shaped mass of grey matter, hence the name, sitting in the medial temporal lobe immediately in front of the hippocampus and beneath the uncus. If the hippocampus asks what happened here before, the amygdala asks what it is worth, and its output is not a record but a response.

It is a collection of nuclei rather than a single structure, and three groups matter:

\begin{itemize}
    \item The \textbf{basolateral group} is the input side and is cortex-like in its cellular organisation. It receives highly processed sensory information from the association cortices and, importantly, a second and much cruder stream directly from the thalamus. It is where the association between a stimulus and its significance is actually learned.
    \item The \textbf{centromedial group}, and in particular the \textbf{central nucleus}, is the output side. It projects to the hypothalamus and brainstem and is the structure that converts a judgment of significance into a bodily response.
    \item The \textbf{cortical group} receives direct input from the olfactory bulb, one of very few places in the brain to do so, and is part of the reason smell reaches emotion so directly.
\end{itemize}

The existence of two input routes is worth dwelling on, because it explains a familiar experience. A sensory signal can reach the basolateral amygdala directly from the thalamus in a single synapse, fast but carrying little detail, or it can travel the longer way through primary and association cortex and arrive slower but fully analysed. The short route is what starts the body reacting to a shape on the path before the cortex has established that it is a stick rather than a snake; the long route is what subsequently calls the reaction off. The cost of a false alarm is trivial and the cost of a missed threat is not, and the wiring reflects that asymmetry.

Output from the central nucleus fans out to a set of targets each of which produces one component of an emotional response, and the pattern we recognise as fear is the sum of them:

\begin{itemize}
    \item To the \textbf{lateral hypothalamus}, driving the sympathetic response, the raised heart rate and blood pressure of the autonomic section.
    \item To the \textbf{paraventricular nucleus}, driving \textbf{CRH} release and with it the whole endocrine stress axis described in the hypothalamus section.
    \item To the \textbf{periaqueductal grey} of the midbrain, which organises defensive behaviour, freezing and flight.
    \item To the \textbf{locus coeruleus}, raising noradrenaline and with it vigilance, and to the \textbf{basal forebrain} cholinergic neurons, sharpening attention.
    \item To the \textbf{parabrachial nucleus}, altering the pattern of breathing.
\end{itemize}

Its two major tracts have already been named in the hypothalamus section: the long, arcing \textbf{stria terminalis} and the short, direct \textbf{ventral amygdalofugal pathway}. Connections with the \textbf{mediodorsal thalamus} and, through it, the prefrontal cortex are the route by which an emotionally significant event captures deliberation rather than merely producing a reflex, and the reciprocal projection from the \textbf{ventromedial prefrontal cortex} back onto the amygdala is the route by which deliberation restrains it.

Clinically the amygdala is best known through two syndromes. Bilateral destruction produces \textbf{Kl\"uver--Bucy syndrome}, in which animals and, rarely, patients become placid and lose the capacity for fear or appropriate aggression, along with hyperorality, indiscriminate sexual behaviour, and a failure to recognise objects visually. \textbf{Urbach--Wiethe disease}, which calcifies the amygdalae bilaterally, has yielded the striking case of a patient who could recognise every other emotion in a face but not fear, and who reported no fear herself in situations that reliably produce it in others. In the opposite direction, an overactive amygdala under weak prefrontal restraint is one of the most consistent findings in anxiety disorders and post-traumatic stress disorder.

\subsubsection{Hypothalamus}

The hypothalamus has a section of its own, and the point to carry into this one is the role it was given there: it is the \textbf{output stage} of the limbic system. The amygdala and the hippocampus determine what a situation means and what it resembles, and neither can do anything about it. The hypothalamus is the only member of the system with the connections to act, commanding the autonomic nervous system through the brainstem and spinal cord and the endocrine system through the pituitary, and it is therefore the point at which emotional significance becomes bodily change. This is the anatomical reason emotion is felt in the body rather than merely registered in the head.

Two of its parts belong to this section specifically. The \textbf{mammillary bodies} at its posterior underside are not autonomic or endocrine at all: they receive the fornix from the hippocampus and pass its signal on through the mammillothalamic tract, which makes them a memory relay sitting inside a homeostatic organ. The \textbf{lateral hypothalamic area}, through which the medial forebrain bundle runs, is the main recipient of amygdalar output and a central node in motivated behaviour.

\subsubsection{Thalamus}

The thalamus likewise has its own section, and three of its nuclear groups are limbic members rather than sensory relays.

The \textbf{anterior nuclear group} is the thalamic station of the Papez circuit, receiving the mammillothalamic tract and projecting to the cingulate cortex. It is a memory structure, and damage to it or to the mammillary bodies feeding it produces the dense anterograde amnesia of Korsakoff's syndrome.

The \textbf{mediodorsal nucleus} is the thalamic partner of the prefrontal cortex, and it receives from the amygdala, the olfactory cortex, and the hypothalamus. It is the point at which limbic information reaches the executive machinery of the frontal lobe, and it is also where olfaction takes its deferred thalamic step on the way to the orbitofrontal cortex.

The \textbf{midline} and \textbf{intralaminar nuclei} connect with the hippocampus, amygdala, and hypothalamus and carry arousal, visceral, and affective signals; the intralaminar group in particular delivers the unpleasantness of pain to the anterior cingulate and insula.

The general point is that the thalamus is the limbic system's route to the cortex in exactly the same way it is the sensory systems' route to the cortex, and it gates the traffic in the same way. What a limbic structure can make the cortex attend to is subject to the same thalamic permission as what the retina can.

\subsubsection{Septal Area}

The \textbf{septal area} is a small region of grey matter lying beneath the front end of the corpus callosum, in front of the anterior commissure and continuous above with the septum pellucidum, the thin partition between the lateral ventricles. It is easy to overlook because of its size, and it is the hinge on which two of the system's circuits turn.

Its central relationship is with the hippocampus, and it is reciprocal. The \textbf{septohippocampal pathway} runs through the fornix in both directions: the hippocampus projects to the septal nuclei, and the \textbf{medial septal nucleus} projects back, sending cholinergic and GABAergic fibres that pace hippocampal activity. What they pace is the \textbf{theta rhythm}, a four to eight hertz oscillation that dominates the hippocampus during exploration and attentive behaviour and that appears to be the timing signal under which encoding takes place. Silence the medial septum and the theta rhythm stops and hippocampal learning is impaired, which is a rare case of a memory structure being disabled by removing a clock rather than a component.

The septal area also connects forward to the prefrontal cortex, downward to the hypothalamus and midbrain through the medial forebrain bundle, and back to the \textbf{habenula} by way of the \textbf{stria medullaris thalami}. It sits immediately above the \textbf{nucleus accumbens}, with which it is often discussed, and this vicinity is where Olds and Milner in 1954 found that rats would press a lever to the exclusion of food and sleep in order to stimulate their own brains, the experiment that opened the study of reward.

Two further points are worth carrying. Septal lesions in animals produce a marked increase in defensive aggression, the phenomenon known as \textbf{septal rage}, which suggests the region normally exerts a restraining influence on the amygdala and hypothalamus. And the cholinergic cells of the medial septum belong to the same \textbf{basal forebrain} population as the nucleus basalis of Meynert, whose degeneration is a feature of Alzheimer's disease and the target of the cholinesterase inhibitors used to treat it. The role of acetylcholine in attention and memory, noted when the transmitters were first introduced, has its anatomical seat here.

\subsubsection{Habenula}

The \textbf{habenula} is a pair of small nuclei in the \textbf{epithalamus}, sitting at the posterior end of the roof of the third ventricle just in front of the pineal gland, joined across the midline by the habenular commissure. It is the smallest structure given its own heading in these notes and, until fairly recently, among the least studied. It is included because it is the system's negative pole, and because it makes direct contact with the dopamine machinery described in the basal ganglia section.

Its connections are unusually tidy. Input arrives chiefly along the \textbf{stria medullaris thalami} from the septal area and basal forebrain, and output leaves along the \textbf{fasciculus retroflexus} to the midbrain. The two habenular nuclei then diverge:

\begin{itemize}
    \item The \textbf{medial habenula} projects to the interpeduncular nucleus and is largely cholinergic. It is densely supplied with nicotinic receptors and is implicated in the aversive state of nicotine withdrawal.
    \item The \textbf{lateral habenula} is the functionally important one. It projects, mostly by way of a small GABAergic relay in the midbrain, onto the dopaminergic neurons of the ventral tegmental area and the substantia nigra pars compacta, and onto the serotonergic raphe nuclei, and its effect on all of them is inhibitory.
\end{itemize}

What makes the lateral habenula interesting is the signal it carries. Midbrain dopamine neurons fire in response to reward, and more precisely to reward that is better than expected. Lateral habenular neurons do the exact opposite: they fire when an expected reward fails to arrive, and when something aversive occurs, and by inhibiting the dopamine neurons they impose the dip in dopamine that follows disappointment. Read against the section on the basal ganglia, where dopamine was shown to set the gain on whether an action is released at all, the habenula is the structure that turns that gain down when the evidence says the action was not worth taking.

The clinical interest follows from that description. A structure whose job is to suppress dopaminergic and serotonergic activity in response to failure is an obvious candidate for involvement in depression, and hyperactivity of the lateral habenula is now among the better supported physiological findings in depression in both animal models and human imaging. It has become a subject of considerable interest as a target, both for deep brain stimulation and in accounting for the action of drugs whose antidepressant effect is faster than the autoreceptor adaptation described earlier can explain.

\subsection{Function}

\subsubsection{Olfaction}

Olfaction is where the limbic system began, both in Broca's description of it and, on the usual reading of the comparative anatomy, in evolutionary terms. The older literature calls these structures the \textbf{rhinencephalon}, the ``nose brain'', and although the name is now too narrow to be useful it records something real about how the system is put together.

The peculiarity of smell was set out in the temporal lobe section: alone among the senses it reaches cortex without relaying in the thalamus. What that section did not emphasise is that the cortex it reaches is \textit{limbic} cortex. The olfactory tract projects to the piriform and periamygdaloid cortex on the uncus, to the entorhinal cortex, and directly to the cortical nucleus of the amygdala. Every one of those targets is a structure in the present section. A smell is therefore two synapses from the hippocampus and one from the amygdala, while a sight or a sound must pass through a thalamic relay, a primary cortex, and an association cortex before it arrives anywhere comparable.

This is the anatomy behind an experience everyone recognises: that a smell can produce a specific autobiographical memory, complete with its emotional colouring, more abruptly and more completely than any other sensory cue. It is also why olfactory hallucinations are a characteristic aura of medial temporal lobe seizures, the so-called uncinate fits, and why the loss of smell is an early sign in Alzheimer's disease and in Parkinson's disease, both of which affect these structures early.

The evolutionary reading, which should be marked as an interpretation rather than a fact, is that the machinery came first and the functions were layered onto it. A chemical arriving at the nose of a simple vertebrate had to be classified as food, threat, or mate, remembered, and acted on, and that is a requirement for valuation, memory, and autonomic output built around an olfactory input. On this account emotion and memory as we experience them are elaborations of a system that originally existed to decide what a smell meant, which is why the structures that serve them are still wrapped around the olfactory pathway.

\subsubsection{Memory}

Memory is not one faculty, and the first step is to divide it, because the divisions map onto different structures. \textbf{Declarative} (or explicit) memory is memory for facts and events that can be brought to mind and stated, and it divides into \textbf{episodic} memory for experienced events and \textbf{semantic} memory for knowledge detached from the occasion of learning it. \textbf{Non-declarative} (or implicit) memory covers everything that is expressed in performance rather than in recollection: motor skills and habits, which depend on the basal ganglia and cerebellum; conditioned emotional responses, which depend on the amygdala; and priming, which depends on the cortex directly. Only the declarative branch is the hippocampus's business, and that is why hippocampal damage produces such a strange clinical picture.

The circuit that carries it is the \textbf{Papez circuit}, met twice already and worth setting out once as a whole:

\begin{align*}
\text{hippocampus} &\xrightarrow{\;\text{fornix}\;} \text{mammillary bodies} \xrightarrow{\;\text{mammillothalamic tract}\;} \text{anterior thalamus} \\
&\rightarrow \text{cingulate} \xrightarrow{\;\text{cingulum}\;} \text{parahippocampal cortex} \rightarrow \text{hippocampus}
\end{align*}

Papez proposed it as a circuit for emotion and it is no longer read that way, but as a memory circuit it has held up remarkably well, and the clinical evidence is that damage at any point along it produces amnesia. Lesions of the hippocampus itself do so, and so do lesions of the fornix, the mammillary bodies, and the anterior thalamus, which is what makes \textbf{Korsakoff's syndrome} such a clean demonstration: thiamine deficiency damages the mammillary bodies and medial thalamus specifically, sparing the hippocampus, and produces a dense amnesia nonetheless.

The classical case at the hippocampal end is the patient known during his lifetime as \textbf{H.M.}, whose medial temporal lobes were removed bilaterally in 1953 to control intractable epilepsy. The result was a profound \textbf{anterograde amnesia}: he could not form new declarative memories at all, and forgot a conversation within minutes of having it. What was spared is as informative as what was lost. His intelligence and language were intact. His working memory was intact, so he could hold a number in mind while attending to it. His memories from well before the surgery were largely intact. And he could learn new motor skills perfectly well while denying each time that he had ever practised them. Three conclusions were drawn, and they have survived: that declarative memory is separable from other kinds, that the hippocampus is required for laying down new memories but is not where old ones are stored, and that immediate holding of information is a different operation altogether.

That second conclusion is the process called \textbf{systems consolidation}. On the standard account the hippocampus binds together the scattered cortical fragments of an experience, the sight and sound and place and feeling of it, and by repeatedly reactivating them, including during sleep, gradually strengthens the direct cortical connections among them until the binding structure is no longer required. A memory therefore migrates, in effect, from being hippocampus-dependent to being cortically self-sufficient. The prediction this makes is that retrograde amnesia after hippocampal damage should be \textbf{temporally graded}, with recent memories lost and remote ones spared, and that is broadly what is observed.

At the cellular level the mechanism is \textbf{long-term potentiation}, and it draws directly on the receptor material of the firing section. The key element is the \textbf{NMDA receptor}, which is a coincidence detector: it is a glutamate-gated channel, but its pore is physically plugged by a magnesium ion at resting potential, and the block is only expelled when the membrane is already depolarised. It therefore opens only when two conditions are met at once, that the presynaptic cell has released glutamate and that the postsynaptic cell is already active. When both hold, calcium enters through the NMDA receptor and triggers a cascade whose short-term effect is to insert more AMPA receptors into the postsynaptic membrane, and whose long-term effect, by way of the activity-dependent gene expression noted when the neuron's nucleus was described, is structural change in the synapse itself. A synapse that has been part of a successful coincidence is thereafter stronger. Hebb's much-quoted formulation, that cells which fire together wire together, turns out to have a molecule behind it.

Finally, memory and emotion are not independent, and the anatomy says why. The amygdala projects to the hippocampus, and emotional arousal releases adrenaline peripherally and cortisol through the axis described in the hypothalamus section, both of which act back on these structures. Emotionally significant events are consequently encoded more strongly and retained better than neutral ones, which is adaptive in most circumstances and pathological in a few: the intrusive, vivid, involuntary memories of post-traumatic stress disorder are, on one widely held view, the same consolidation mechanism running at excessive gain.

\subsubsection{Emotional Responses}

An emotional response can be decomposed into four stages, and each has an anatomical home in the structures above. There is \textbf{appraisal}, the judgment that something matters and in what way. There is \textbf{output}, the bodily and behavioural change that follows. There is \textbf{feeling}, the conscious experience that accompanies it. And there is \textbf{regulation}, the modification of all three in the light of context.

\textbf{Appraisal} is the amygdala's, working with the orbitofrontal and ventromedial prefrontal cortex. The basolateral amygdala learns which stimuli predict which outcomes, and the best-worked-out case is \textbf{fear conditioning}: a neutral tone paired with an aversive event acquires, within a few trials, the capacity to produce the full defensive response by itself, and the synaptic change underlying it occurs in the basolateral amygdala by the same NMDA-dependent potentiation described for memory. The amygdala is not exclusively a fear structure, and it responds to appetitive and socially significant stimuli too, but fear is where the circuit has been traced most completely.

\textbf{Output} is the central nucleus's, by way of the projections listed earlier, and it has three limbs which are worth keeping distinct because they operate on different timescales. The autonomic limb runs through the hypothalamus and brainstem and acts within a second. The endocrine limb runs through CRH, ACTH, and cortisol and acts over minutes to hours. The behavioural limb runs through the periaqueductal grey, which organises the coordinated patterns of freezing, flight, and defensive aggression, and which is also the origin of the descending analgesia that makes injuries sustained in an emergency go unnoticed.

\textbf{Feeling} is the least settled of the four and is where this section makes contact with the insula. The classical dispute is between the \textbf{James--Lange} position, that the bodily change comes first and the feeling is the perception of it, and the \textbf{Cannon--Bard} position, that the two arise in parallel from a central structure. Neither is now held in its original form, but the insula section supplied something that reads as a partial vindication of the first: the posterior insula constructs an interoceptive map of the body's state and the anterior insula re-represents it as a felt quality. On that account the racing heart is not merely a symptom of fear but part of what is read to produce the experience of it, and the anterior insula and anterior cingulate are where the reading happens.

\textbf{Regulation} runs from the prefrontal cortex back onto the amygdala, chiefly from the ventromedial and orbitofrontal regions described in the frontal lobe section. Its clearest experimental expression is \textbf{extinction}: a conditioned fear response that is no longer reinforced does not fade because the original association has been erased, but because a second, inhibitory association is learned, and it is the ventromedial prefrontal cortex that holds it and imposes it on the amygdala. This has two consequences worth knowing. It explains why extinguished fears return under stress or in a new context, since the inhibitory memory is more fragile and more context-bound than the original. And it explains why the combination of an overactive amygdala with an underactive ventromedial prefrontal cortex recurs so persistently in the imaging literature on anxiety disorders and post-traumatic stress disorder: the appraisal is intact and the brake is not.

\subsubsection{Motivation and Behaviours}

The behaviours the limbic system organises are traditionally summarised as feeding, fighting, fleeing, and reproduction, and although the mnemonic is crude it identifies the right set. What unifies them is that each is a case of the organism acting to change its own state rather than responding to the world, and each requires the same three things: a signal of what is needed, a valuation of what is available, and a mechanism for selecting among the options.

The valuation machinery is the \textbf{reward system}, and its core is dopaminergic. The \textbf{ventral tegmental area (VTA)} of the midbrain projects along the medial forebrain bundle to the \textbf{nucleus accumbens} in the ventral striatum, the \textbf{mesolimbic pathway}, and to the prefrontal cortex, the \textbf{mesocortical pathway}. These are the same cells, and the same transmitter, whose motor counterpart was traced from the substantia nigra pars compacta in the basal ganglia section, and the parallel is not a coincidence. There is a \textbf{limbic loop} through the basal ganglia with exactly the architecture of the motor loop set out there: limbic and prefrontal cortex projects to the ventral striatum, the ventral striatum to the ventral pallidum, the ventral pallidum to the mediodorsal thalamus, and the thalamus back to the anterior cingulate and orbitofrontal cortex. The same disinhibitory arithmetic, the same dopaminergic gain control, applied not to the selection of a movement but to the selection of a goal.

Two refinements of the popular account are worth having. The first is that dopamine does not signal pleasure. What midbrain dopamine neurons encode is closer to a \textbf{reward prediction error}, the difference between the reward obtained and the reward expected: they fire to an unexpected reward, fall silent when an expected one fails to arrive, and come to fire to the cue that predicts a reward rather than to the reward itself. It is a teaching signal, not a pleasure signal, and the lateral habenula described above supplies its negative half. The second, related, refinement is the distinction between \textbf{wanting} and \textbf{liking}. Dopamine supports the incentive salience of a goal, the pull towards it and the effort spent pursuing it, while the hedonic impact of actually obtaining it depends on opioid and endocannabinoid signalling in small regions of the accumbens and pallidum. The two can be dissociated experimentally, and the dissociation is the most useful single fact about addiction: a drug can come to be wanted enormously while being liked very little.

Addiction follows from that architecture rather than from any special property of the drugs. The agents described in the mimicry subsection reach the same endpoint by different doors: nicotine as an agonist at nicotinic receptors sitting on VTA dopamine neurons, opioids as agonists at $\mu$ receptors that disinhibit them, and cannabinoids at CB1. Each raises accumbens dopamine more reliably and more steeply than any natural reward, and the teaching signal, which evolved on the assumption that such a discrepancy indicates something of real value, duly does its job.

The remaining behaviours are best traced as circuits already assembled. \textbf{Feeding and drinking} run from the arcuate, ventromedial, and lateral hypothalamic nuclei described earlier, and the lateral area's role in seeking food is inseparable from its position on the medial forebrain bundle: it is a hunger centre and a reward pathway at the same time. \textbf{Defence and aggression} run from the amygdala to the medial hypothalamus and thence to the periaqueductal grey, with the septal area exerting the restraining influence whose removal produces septal rage. \textbf{Sexual behaviour} depends on the medial preoptic area and the amygdala together, under endocrine control through the axes of the hypothalamus section.

The last point to make is the one the frontal lobe section anticipated. The limbic system generates drives that are, in themselves, indifferent to circumstance: hunger does not know whether it is an appropriate moment to eat, and fear does not know whether the room contains people whose opinion matters. What supplies that knowledge is the prefrontal cortex, through its reciprocal connections with the amygdala, the hypothalamus, and the ventral striatum. Self-control, in the anatomical sense, is not the absence of limbic drive but the presence of a projection capable of modulating it, and most of the conditions in which behaviour goes conspicuously wrong, addiction, impulsivity after orbitofrontal damage, the disinhibition of frontotemporal dementia, are failures of that projection rather than excesses of the drive itself.

\subsection{Summary}

The limbic system is a set of structures defined by their connections rather than their location, forming a ring of transitional cortex around the corpus callosum together with a group of subcortical nuclei that ring encloses. The \textbf{limbic lobe} supplies the cortical component: the cingulate gyrus, whose anterior part weights experience for salience and affective pain and whose posterior part serves spatial and autobiographical memory, and the parahippocampal gyrus, whose entorhinal cortex is the gateway to the hippocampus. The \textbf{hippocampus}, three-layered and wired as an almost unidirectional trisynaptic circuit, is required for laying down new declarative memories and carries a map of space. The \textbf{amygdala} appraises significance, learning associations in its basolateral group and issuing them from its central nucleus to the hypothalamus, brainstem, and endocrine axis. The \textbf{hypothalamus} is the output stage that turns those judgments into bodily change; the \textbf{thalamus}, through its anterior, mediodorsal, midline, and intralaminar nuclei, is the route by which limbic signals reach the cortex and are gated on the way; the \textbf{septal area} paces the hippocampal theta rhythm and restrains the amygdala; and the \textbf{habenula} supplies the system's negative pole, inhibiting midbrain dopamine when an expected reward fails to arrive.

Functionally, the system is built around olfaction in the sense that smell reaches it directly and without a thalamic relay, which is why odour recruits memory and emotion more immediately than any other sense. Its memory role runs through the Papez circuit, damage anywhere along which produces amnesia, and rests cellularly on NMDA-dependent long-term potentiation, a coincidence detector requiring transmitter and depolarisation together. Its emotional role decomposes into appraisal by the amygdala, output through autonomic, endocrine, and behavioural limbs of quite different timescales, feeling constructed interoceptively in the insula, and regulation imposed from the ventromedial prefrontal cortex. Its motivational role runs through a limbic loop of the basal ganglia with precisely the architecture of the motor loop, in which dopamine from the ventral tegmental area encodes not pleasure but the discrepancy between reward expected and reward obtained, and sets the gain on whether a goal is pursued. The components and their principal connections are gathered in Table~\ref{tab:limbic}. Taken together, the limbic system is where the nervous system stops asking what is present and starts asking what it is worth and what happened last time, and it is the point at which those two questions are converted into a body doing something about it.

With this the ascent set out at the beginning is as complete as these notes can currently make it. The substrate was established in the first four sections, the construction of a world in the next five, and the machinery of selection, drive, gating, and valuation in the last four. The worked example that follows runs a single event through the whole of it, from photons and pressure waves to a body that has already changed and a feeling that arrives afterwards, and the closing pages take stock of what such an account does and does not deliver.

{\footnotesize
\renewcommand{\arraystretch}{1.3}
\begin{longtable}{p{2.6cm} p{3.2cm} p{3.2cm} p{4.0cm}}
\caption{The principal components of the limbic system, with their major connections and the role each plays. The named tracts are those visible in dissection and used throughout this section; several structures appear here and in earlier sections under other headings, and the cross-references in the prose identify where.}
\label{tab:limbic} \\
\toprule
\textbf{Component} & \textbf{Main input} & \textbf{Main output} & \textbf{Role} \\
\midrule
\endfirsthead
\multicolumn{4}{l}{\emph{\small Table~\ref{tab:limbic} continued from previous page}} \\
\toprule
\textbf{Component} & \textbf{Main input} & \textbf{Main output} & \textbf{Role} \\
\midrule
\endhead
\bottomrule
\endfoot
\bottomrule
\endlastfoot
Cingulate gyrus & Anterior thalamus; intralaminar nuclei & Cingulum to parahippocampal gyrus; autonomic centres & Salience, conflict, affective pain (anterior); spatial and autobiographical memory (posterior) \\
\midrule
Parahippocampal gyrus (entorhinal) & Association cortex; olfactory tract & Perforant path to hippocampus & Gateway to and from the hippocampus; olfactory association; grid cells \\
\midrule
Hippocampus & Entorhinal cortex (perforant path) & Fornix to mammillary bodies and septal area & Consolidation of declarative memory; spatial map \\
\midrule
Amygdala & Thalamus (direct); association cortex; olfactory bulb & Stria terminalis and ventral amygdalofugal pathway to hypothalamus and brainstem & Appraisal of significance; fear conditioning; emotional output \\
\midrule
Hypothalamus & Fornix; amygdalar pathways; prefrontal cortex & Brainstem and spinal autonomic centres; pituitary & Output stage: autonomic, endocrine, and behavioural response \\
\midrule
Thalamus (anterior, MD, midline, intralaminar) & Mammillothalamic tract; amygdala; hypothalamus & Cingulate; prefrontal cortex; insula & Route to cortex for memory, emotion, and affective pain \\
\midrule
Septal area & Fornix from hippocampus; hypothalamus & Septohippocampal pathway; stria medullaris to habenula & Hippocampal theta rhythm; restraint of aggression; reward \\
\midrule
Habenula & Stria medullaris from septal area and basal forebrain & Fasciculus retroflexus to VTA, SNc, raphe & Inhibition of dopamine on reward omission; aversive signalling \\
\end{longtable}
}


\subsection{A Worked Example: The Bear in the Woods}

Everything in this section has been presented structure by structure, which is how it has to be taught and is not how it ever happens. It is worth ending by running a single event through the whole system in order, because the components were described separately but they operate together and on quite different timescales, and the interleaving is most of what makes the response feel the way it does.

The situation is this. You are walking a wooded path. Twenty or thirty metres ahead and slightly to one side, a branch cracks, and a large dark shape moves in the undergrowth. It is a black bear. What follows traces that event from photons and pressure waves through to a body that has already changed and a mind that is only beginning to catch up.

\subsubsection{Arrival: The Sensory Pathways}

Nothing about the bear enters the nervous system as a bear. What enters is a pattern of light on two retinas and a pattern of pressure on two basilar membranes, and both must travel the pathways set out in the occipital and temporal lobe sections before they mean anything at all.

The visual signal leaves the retinal ganglion cells along the optic nerve, crosses partially at the optic chiasm so that each hemisphere carries the contralateral half of the field, and arrives at the \textbf{lateral geniculate nucleus}. The auditory signal leaves the cochlea along the cochlear division of CN VIII and climbs through the cochlear nuclei, superior olivary complex, lateral lemniscus, and inferior colliculus to the \textbf{medial geniculate nucleus}. Both are now in the thalamus, roughly twenty to forty milliseconds after the event, and neither has yet been identified as anything. The LGN holds a retinotopic pattern of contrast and motion; the MGN holds a tonotopic pattern with a broadband transient in it.

One response has already happened without troubling any of this. The crack of the branch drives a short brainstem reflex arc from the cochlear root neurons to the caudal pontine reticular formation and thence to motor neurons, producing the \textbf{acoustic startle} within about ten to twenty milliseconds: the whole-body flinch that precedes any knowledge of what caused it. It is worth naming because the amygdala does not produce it but does modulate it, and the magnitude of a startle is a standard laboratory measure of how frightened an animal already was. Yours will be larger than it would have been an hour ago on an open road.

\subsubsection{Two Roads to the Amygdala}

From the thalamus the signal takes two routes at once, and the difference between them is the heart of the example.

The \textbf{low road} is a direct projection from the thalamus to the basolateral amygdala, bypassing cortex entirely:

\[
\text{retina, cochlea} \rightarrow \text{LGN, MGN} \rightarrow \text{basolateral amygdala}
\]

It is fast and it is crude. What it can deliver is something on the order of \textit{large, dark, low, close, moving} — enough to specify that a big object is nearby and approaching, and nothing whatever about what the object is. The amygdala is therefore beginning to respond while the visual cortex is still assembling edges.

The \textbf{high road} runs the full course described in the earlier sections:

\[
\text{LGN} \rightarrow \text{V1} \rightarrow \text{ventral stream} \rightarrow \text{inferior temporal cortex} \rightarrow \text{basolateral amygdala}
\]

V1 responds at perhaps forty to sixty milliseconds and performs its analysis of edges, orientation, and motion; the ventral ``what'' stream carries the result forward into inferior temporal cortex, where identification occurs somewhere around a hundred to two hundred milliseconds; the dorsal ``where and how'' stream simultaneously establishes distance, heading, and whether the shape is closing. In parallel the auditory signal reaches A1 on Heschl's gyri and passes to the auditory association cortex, which classifies the sound as a large animal in undergrowth rather than as wind. The \textbf{posterior association area} then does the job it was given in the parietal lobe section and binds these into one object: the seen shape and the heard sound are registered as a single thing at a single place, rather than as two unrelated reports.

The asymmetry between the roads is the point, and the justification is the asymmetry of the costs. Reacting to a stick as though it were a bear costs a few seconds and some adrenaline; failing to react to a bear costs everything. A system facing that payoff structure should commit early on evidence that would not survive scrutiny, and then withdraw the commitment when scrutiny arrives, and that is precisely what the two roads implement.

It should be said that the human low road is less firmly established than the textbook diagrams imply. The subcortical route was worked out in rodents, where it is not in doubt; in humans its existence and its importance are both actively disputed, with some evidence that rapid amygdala responses could be carried by fast cortical routes instead. The functional principle, that a coarse fast appraisal precedes a fine slow one, is not in question. The anatomy by which it is achieved in our own brains is less settled than the story usually told.

\subsubsection{Appraisal: What It Means}

The basolateral amygdala now holds both streams, and appraisal is the convergence of three things upon it.

The first is the sensory evidence just described. The second is \textbf{context}, supplied by the hippocampus and entorhinal cortex: this is a wood, in bear country, and there was a sign at the trailhead. Context is what makes the same dark shape terrifying here and unremarkable behind glass at a zoo, and the division of labour is clean, since a cue can be conditioned by the amygdala alone but a context cannot be represented without the hippocampus. The third is \textbf{prior learning}, which in this case is worth a remark. You have almost certainly never been attacked by a bear, so the association is not one you acquired by conditioning in the way the laboratory paradigm describes. It arrives instead through semantic memory, from cortex, and possibly through the readiness with which certain classes of stimulus are learned as threatening at all. The amygdala's business is valuation, and it does not much care where the valuation was learned.

The output of that convergence goes to the \textbf{central nucleus}, and from this point the example stops being about perception and starts being about response.

\subsubsection{The Autonomic Limb: Under a Second}

The central nucleus projects to the lateral hypothalamus, which commands the descending autonomic pathways of the hypothalamus section, reaching the autonomic centres of the medulla and the \textbf{intermediolateral cell column} of the thoracic cord.

The first change is not an addition but a subtraction: \textbf{vagal withdrawal}. Resting parasympathetic tone on the heart is removed, and because that only requires an ongoing inhibition to stop rather than a new signal to arrive, the heart rate begins to climb within a beat or two. Sympathetic drive follows over the next several seconds and adds to it. The pattern that results is the one set out when the autonomic nervous system was introduced, now with its causes attached:

\begin{itemize}
    \item Heart rate and contractility rise; blood pressure rises; vasoconstriction in the skin and viscera diverts blood towards skeletal muscle.
    \item The pupils dilate, the sympathetic supply to dilator pupillae overriding the parasympathetic constrictor fibres carried by CN III, which widens the aperture at the cost of depth of focus.
    \item The airways dilate and respiration deepens and quickens.
    \item Gut motility and salivation are suppressed, which is where the dry mouth comes from, and sweating begins, by way of the sympathetic fibres that are anomalously cholinergic.
    \item Glycogenolysis and lipolysis are driven, and within seconds the \textbf{adrenal medulla} releases adrenaline and noradrenaline into the circulation, sustaining all of the above hormonally after the neural signal has done its part.
\end{itemize}

None of this has yet been decided on, and none of it has yet been noticed.

\subsubsection{The Behavioural Limb: The Same Second}

In parallel the central nucleus drives the \textbf{periaqueductal grey}, and the PAG determines which defensive behaviour is produced. Its columns are functionally distinct, and the selection between them tracks how imminent the threat is:

\begin{itemize}
    \item The \textbf{ventrolateral column} produces freezing and quiescence, and dominates when the threat is detected but distant and has not yet oriented on you.
    \item The \textbf{dorsolateral and lateral columns} produce active flight or defensive aggression, and dominate when the threat is close or closing.
\end{itemize}

At thirty metres, with the bear apparently unaware of you, the ventrolateral column has it, and you stop dead. Freezing is not a failure to act; it is an action, selected because movement is what most reliably attracts a predator's attention. Had the bear been at five metres and coming, the same nucleus would have produced the opposite behaviour from a neighbouring column.

The PAG also drives \textbf{descending analgesia}, which is why an injury sustained in the next thirty seconds may not be felt until afterwards. And the movement, when it does come, is released by the machinery of the basal ganglia section, with the thalamic brake lifted through VA and VL and dopamine setting the threshold at which the whole system consents to act.

\subsubsection{The Arousal Limb: Hundreds of Milliseconds}

The central nucleus simultaneously drives the ascending modulatory systems, and their effect is to change how the rest of the brain operates rather than to produce any behaviour of their own. It reaches the \textbf{locus coeruleus}, which broadcasts noradrenaline across the cortex and raises vigilance while narrowing the focus of attention; the \textbf{basal forebrain} cholinergic neurons, which sharpen cortical processing; and the \textbf{raphe} and \textbf{ventral tegmental area}.

Two consequences follow in structures described earlier. The thalamic relay neurons are pushed firmly into \textbf{tonic mode}, the high-fidelity firing pattern of full waking, so that sensory detail is transmitted as faithfully as the system can manage. And the \textbf{reticular nucleus} adjusts its inhibitory gate so that what is relayed is what concerns the bear, which is what the narrowing of attention consists of at the level of cells. The reason peripheral detail is missing from the account you give later is not that you failed to remember it but that it was never forwarded.

\subsubsection{The Endocrine Limb: Minutes to Hours}

The slowest limb starts at the same moment as the others and takes far longer to arrive. The central nucleus drives the \textbf{paraventricular nucleus}, which releases \textbf{CRH} into the hypophyseal portal system; the anterior pituitary releases \textbf{ACTH}; the adrenal cortex releases \textbf{cortisol}, which peaks in the blood something like fifteen to twenty minutes after the branch cracked, long after the bear has gone.

Cortisol mobilises glucose, modulates immune function, and acts back on the hippocampus, hypothalamus, and prefrontal cortex, both to shut the axis down by the negative feedback described in the hypothalamus section and, as the memory subsection noted, to strengthen the consolidation of what is being encoded. This limb is the reason you are still shaking twenty minutes later on a path where nothing is happening: the neural signal ended long ago and the chemical one has not.

\subsubsection{Feeling: Where the Fear Is Actually Experienced}

So far nothing in this account has required you to feel anything, and that is deliberate, because the feeling arrives late and is assembled from the response rather than preceding it.

The bodily changes just produced are themselves sensed. Interoceptive and small-fibre afferents from the heart, lungs, gut, and vasculature ascend to \textbf{VMpo} and neighbouring thalamic nuclei and reach the \textbf{posterior insula}, which holds them as something close to a raw physiological reading. That reading is re-represented forward through the mid-insula to the \textbf{anterior insula}, where it becomes a felt quality: the pounding, the tightness, the dry mouth, the hollow sensation in the abdomen. The \textbf{anterior cingulate}, the anterior insula's partner in the salience network, flags the whole episode as the thing that matters and suppresses everything else competing for attention.

At the same time the amygdala's signal reaches the \textbf{mediodorsal nucleus} of the thalamus and through it the prefrontal cortex, which is the step at which the event becomes available to deliberation rather than merely producing a response. This is roughly where you would say the fear began, and it is late, several hundred milliseconds to a second or two after your heart rate had already changed.

The insula section's material is what makes this more than a curiosity. On the interoceptive account, the racing heart is not simply a symptom of the fear but part of what is read in order to produce the experience of it, which is the James--Lange position given an anatomy. The Cannon--Bard objection retains its force, since the amygdala clearly initiates the response rather than inferring it, and the honest summary is that appraisal, response, and feeling are three processes running on overlapping timescales rather than a chain with a single direction.

\subsubsection{Encoding: Why You Will Remember This}

While all of the above is happening, the episode is being written down. The bound percept assembled by the posterior association area reaches the \textbf{entorhinal cortex} and enters the trisynaptic circuit, through the perforant path to the dentate gyrus, the mossy fibres to CA3, and the Schaffer collaterals to CA1 and the subiculum, with the \textbf{theta rhythm} paced by the medial septum running underneath it and NMDA-dependent potentiation strengthening the synapses that carried it.

Two things then make this memory unusually durable. The amygdala projects directly to the hippocampus and enhances encoding there. And the circulating adrenaline, which cannot itself cross the blood--brain barrier, acts on vagal afferents that reach the nucleus of the solitary tract and from there the locus coeruleus, so that the peripheral hormonal response feeds back into the central one; cortisol does the same more slowly. The result is the phenomenon everyone recognises, that you will be able to describe this thirty seconds of your life in detail in twenty years' time while remembering nothing whatever of the hour that preceded it.

Two kinds of conditioning are also being laid down at once, and they can be separated by lesion. The \textbf{cue}, the particular sound and shape, is conditioned by the amygdala. The \textbf{context}, this stretch of path, is conditioned by the hippocampus. Walking here next year will produce an attenuated version of the entire response before you have consciously recalled why, which is the same mechanism working correctly, and is also, when the original event was severe enough and the extinction never occurs, the mechanism of post-traumatic stress disorder.

\subsubsection{Regulation: Calling It Off}

Finally the slow route completes what the fast route started. The cortex has by now established that this is a black bear rather than a grizzly, that it is thirty metres away, that it has not oriented on you, and that black bears are very rarely predatory towards humans. That knowledge is held in the \textbf{ventromedial prefrontal cortex}, which projects back onto the amygdala and inhibits the central nucleus. Sympathetic drive falls, vagal tone returns, the freeze releases into a slow deliberate retreat, and the cortisol continues circulating regardless because nothing can call that back.

Two observations close the example. The first is that this regulation requires knowledge as well as anatomy: the vmPFC can only impose an inhibitory association it actually holds, so a person who knows nothing about bears has the same projection and nothing to send along it, and the same is true of the extinction learning that reduces a fear over repeated safe encounters. The second is that the entire episode is a negotiation between speed and accuracy. The low road commits the body to a response on evidence that would not survive a moment's scrutiny; the high road, the hippocampus, and the prefrontal cortex supply the scrutiny and withdraw the commitment. Fear is what it feels like to be inside that negotiation while it is being conducted.

One structure has been silent throughout and is worth a closing sentence, since it is the component this example would otherwise leave out. The \textbf{lateral habenula} responds to aversive events by inhibiting the midbrain dopamine neurons, which is the negative teaching signal described earlier: alongside everything else, the system is recording that this path, at this hour, was a worse outcome than expected, and adjusting downward the value it will assign to walking it again.

The sequence, with the rough timings that make its structure legible, is gathered in Table~\ref{tab:bear}.

{\footnotesize
\renewcommand{\arraystretch}{1.3}
\begin{longtable}{p{2.4cm} p{4.6cm} p{7.2cm}}
\caption{The bear in the woods, traced in order of arrival. The timings are approximate and are drawn largely from animal work; what matters is not the numbers but the ordering, and in particular that the autonomic response is complete before the feeling of fear has been assembled and long before the endocrine response has arrived.}
\label{tab:bear} \\
\toprule
\textbf{Time} & \textbf{Structure} & \textbf{Event} \\
\midrule
\endfirsthead
\multicolumn{3}{l}{\emph{\small Table~\ref{tab:bear} continued from previous page}} \\
\toprule
\textbf{Time} & \textbf{Structure} & \textbf{Event} \\
\midrule
\endhead
\bottomrule
\endfoot
\bottomrule
\endlastfoot
0 & Retina, cochlea & Transduction of light and pressure \\
10--20 ms & Cochlear root, pontine reticular formation & Acoustic startle, amygdala-potentiated \\
20--40 ms & LGN, MGN & Thalamic relay; nothing yet identified \\
40--80 ms & Basolateral amygdala (low road) & Crude appraisal: large, dark, close \\
40--60 ms & V1, A1 & Primary cortical analysis begins \\
100--200 ms & Ventral and dorsal streams, posterior association area & Identification as a bear; binding of sight and sound \\
150--300 ms & Basolateral to central amygdala; hippocampus & Appraisal completed with context \\
$<$1 s & Vagal withdrawal; PAG (ventrolateral) & Heart rate climbs; freezing \\
1--5 s & IML, sympathetic chain, adrenal medulla; locus coeruleus & Full sympathetic pattern; adrenaline; vigilance and narrowed attention \\
1--10 s & Posterior to anterior insula; anterior cingulate & The bodily state becomes a felt quality \\
5--30 s & Mediodorsal thalamus to PFC; vmPFC to amygdala & Deliberation, then regulation and release of the freeze \\
15--20 min & PVN, CRH, ACTH, adrenal cortex & Cortisol peaks, long after the event \\
Hours to days & Hippocampus, amygdala & Consolidation, enhanced by adrenaline and cortisol \\
Weeks onward & Neocortex & Systems consolidation; cue and context conditioned \\
\end{longtable}
}

\phantomsection
\section*{Taking Stock}
\addcontentsline{toc}{section}{Taking Stock}

The introduction proposed to start from ions crossing a membrane and arrive at a mind, and it is worth ending by asking honestly how much of that distance has been covered.

The ledger is substantial. A system built from the parts described in these notes can discriminate one stimulus from another, and we know how: transduction at a receptor surface, a labelled line to a thalamic nucleus, and a topographic map in primary cortex. It can bind separate sensory streams into a single scene, and we know roughly where, at the parieto-temporo-occipital junction, though rather less about how. It can recognise, which the agnosias show to be a further operation on an already complete percept rather than a better version of perceiving. It can attend, and here the mechanism is unusually concrete: the reticular nucleus decides which relay cells are permitted to fire. It can be more or less awake, and the switch is the mode in which those same cells fire. It can store an event and retrieve it, through a circuit whose anatomy is known in detail and whose synaptic mechanism, in the coincidence detection performed by the NMDA receptor, is known in chemistry. It can assign value, through a structure that learns what predicts what and a signal that reports the difference between what was expected and what arrived. It can want things, because a regulated variable drifting out of range is converted into a demand. It can choose among competing actions by disinhibiting one and holding back the rest. It can model its own body, both from the outside as a shape in space and from the inside as a state that feels like something. And it can regulate all of the above from the prefrontal cortex, which is what makes any of it look like judgment rather than reflex.

That list is worth taking seriously, because a century ago almost none of it was available and much of it would have seemed unavailable in principle. It is also worth being exact about what kind of item each entry is, because they are not all the same. Some are genuine mechanisms, specified well enough that the phenomenon follows from the parts: the resting potential falls out of the Nernst equation and a dominant potassium permeability, and nothing further is required. The disinhibition arithmetic of the basal ganglia is of this kind, as is the opposite action of dopamine on D1 and D2 populations, and as is long-term potentiation. Others are localisations wearing the clothes of mechanisms. We can say that binding happens at the posterior association area without being able to say how separate streams become one object, and saying that the anterior insula re-represents bodily state is a description of a transformation whose actual operation is not specified. The honest position is that the account is dense in some places and thin in others, and that the thin places are mostly the ones nearest to what we most want explained.

Several things are missing entirely and are still to come as the lecture series continues. The cerebellum has appeared only as a comparator and deserves its own treatment. The language network has been described from its two ends without the middle. Sleep has been mentioned as a state of thalamic firing and not as a phenomenon. There is nothing on development, which is where most of the wiring described here actually comes from, and little on the cortical microcircuit, the repeated six-layer motif whose uniformity across regions is one of the more remarkable facts about the brain. Psychiatric illness has appeared only as evidence for how intact machinery works.

There remains the question set aside at the beginning, and it should be left where it was put. Everything above is an account of what the machinery does. None of it addresses why the machinery should be accompanied by any experience at all, and it is worth being clear that no amount of further detail of the kind these notes contain would obviously address it either. Knowing precisely which thalamic cells were firing tonically, and precisely how the anterior insula transformed its input, would tell us with great accuracy when there is an experience and what it is an experience of, and would not tell us why there is something it is like to be the system in question. Some think that gap will close as the mechanisms fill in, and others think it is a different sort of gap. These notes take no position, and readers should notice that the question has been left standing rather than answered quietly along the way.

What does survive, and survives more strongly than it began, is the claim named at the outset: that nothing about the mind is given. Thirteen sections have now made it about different subject matter, and the pattern held every time. The visual world is assembled in stages and comes apart along the seams of its assembly. The sense of occupying a body is built and can be dismantled. Recognition is separable from perception, meaning from sound, the unpleasantness of pain from its location, wanting something from liking it, and the feeling of an emotion from the response that produced it. In each case what feels like a single indivisible act of experience turns out to be a compound with parts that can be lost independently, and the parts correspond to tissue.

The immediacy is not an illusion. It is a real feature of what it is like to have a mind, and it is exactly what one would expect of a system whose constructive work is not available to itself. A brain does not perceive its own processing; it perceives the products, and the products arrive finished. That is why the whole subject has to be approached the way these notes have approached it, from the ions upward and from the lesions inward, because the one vantage point from which the machinery is systematically invisible is the inside.

\end{document}
